\documentclass[../main.tex]{subfiles}
\begin{document}
%==========================================TIÊU ĐỀ TẬP========================================
\begin{table}[h]
  \begin{adjustwidth}{-1cm}{}
    \begin{tabular}{>{\centering\arraybackslash}p{6cm}|>{\raggedright\arraybackslash}p{12.5cm}}
      \multirow{5}{6cm}
      % Cột bên trái: ÉPISODE và số tập
      {\\[-40pt]
        \flaregothic\fontsize{20pt}{20pt}\selectfont \centering
        {\contour{errorfive}{\textcolor{black}{ÉPISODE}}}\color{black}\\
        \burbank\fontsize{72pt}{72pt}\selectfont
        \includegraphics[height=3.12359cm]{../../../../Image/Book?/number/num37.png}
        \\[18pt]
        % \flaregothic\fontsize{14pt}{14pt}\selectfont
        % \color{epattr}BIENVENUE, \uline{\color{Banpire} Mel} !
        % \flaregothic\fontsize{14pt}{14pt}\selectfont{\contour{errorfive}{\textcolor{white}{BIENVENUE, }}
        %     {\color{black} Naomi}
        %     {\contour{errorfive}{\textcolor{white} {!}}}}\\
      }
       &
      % Dòng thứ nhất: tên tập tiếng Nhật
      {\fotskip\fontsize{18pt}{18pt}\selectfont \ruby{祭}{まつ}りの\ruby{熱}{ねつ}\ruby{狂}{きょう}}                   \\[6pt]
      % Dòng thứ hai: tên tập tiếng Anh
       & {\cafeteria\fontsize{18pt}{18pt}\selectfont Festival Fever}                                       \\[4pt]
      % Dòng thứ ba: tên tập tiếng Việt
       & {\vnmsans\fontsize{18pt}{18pt}\selectfont Ngày lễ hội cuồng nhiệt}                                \\[8pt]
      % Dòng thứ tư: nhân vật xuất hiện
       & {\caviar{\fontsize{14pt}{14pt}\selectfont Nhân vật: Tất cả các nhân vật xuất hiện từ đầu series}} \\[8pt]
      % \\[6pt]
      % Dòng cuối cùng: Thời gian diễn ra của tập (thường sẽ là ngày phát hành)
      % & {\continuum\fontsize{16pt}{16pt}\selectfont Ngày phát hành: 11/06/2022}\\[8pt]
    \end{tabular}
  \end{adjustwidth}
\end{table}
%===========================================NỘI DUNG==========================================
\color{black}
\pov{Hikawa Kuon}

Tiếng ve sầu bắt đầu râm ran trên những rặng cây quanh trường trung học Aogami, báo hiệu rằng mùa hè đã thực sự gõ cửa thị trấn.

Từ cửa sổ lớp học, tôi có thể nhìn thấy cả một góc thị trấn đang khoác lên mình tấm áo của mùa mới. Những hàng cây trắc bá và thông Nhật rải rác dọc con đường chính vẫn xanh thẫm như thường lệ, nhưng ánh nắng chiếu qua vòm lá giờ đây mang thêm một sắc vàng chói chang và mồ hôi ướt áo mà vài tháng trước còn chưa có. Dọc theo con kênh nhỏ chảy ngang qua trung tâm Aogami, những bụi lau và cỏ dại đã vươn cao tới mức đủ để che khuất mặt nước lấp lánh bên dưới. Một đàn chim sẻ đậu trên dây điện trước ngôi đền gần cổng trường, líu lo rồi vù bay khi một cậu học sinh đi xe đạp lướt qua bên dưới.

Ở góc sân trường, khóm hoa cẩm tú cầu màu xanh tím nhạt đang nở rộ bên cạnh bờ tường phủ rêu, một màu sắc đặc trưng mà người Nhật hay gọi là màu của mùa mưa chuyển sang hè. Mấy cô gái từ lớp bên đang tụ tập chụp ảnh với khóm hoa đó, tiếng cười nói rộn ràng vọng vào tận lớp. Xa hơn một chút, phía bên kia hàng rào sắt của trường, những mái ngói đen của khu phố cũ nằm im lặng dưới bầu trời xanh cao vời vợi. Một bầu trời kiểu Nhật điển hình vào đầu tháng Bảy, trong vắt đến nỗi những đám mây trắng bồng bềnh trôi trên đó trông như được ai đó vẽ lên bằng bút lông.

Thị trấn Aogami vẫn yên bình như thế. Vẫn nhỏ nhắn, chậm rãi và dịu dàng như thế này, như thể chưa có một thế lực siêu nhiên nào tác động tới nơi này vậy.

Sau những ngày căng thẳng của kỳ thi học kỳ, bầu không khí trong trường như được trút bỏ một gánh nặng khổng lồ. Ánh nắng rực rỡ chiếu xuống sân trường, làm nổi bật lên những gương mặt rạng rỡ của đám học sinh đang háo hức bàn tán về kỳ nghỉ hè sắp tới.

Tuy nhiên, đối lập với sự náo nhiệt ấy là dãy hành lang của khối năm ba. Ở đó, không khí vẫn mang chút gì đó trầm mặc và căng thẳng. Những anh chị năm cuối đang phải gồng mình cho những buổi ôn thi đại học khắc nghiệt, và đối với họ, lễ hội trường học sắp tới không chỉ là một sự kiện giải trí, mà còn là dấu ấn cuối cùng của thời học sinh trước khi bước vào cuộc đời trưởng thành.

Nhưng ở lớp 1-C chúng tôi, tinh thần lễ hội đã lan tỏa mạnh mẽ hơn bất cứ đâu.

``\textbf{Cả lớp trật tự nào! Mọi người chú ý lên bảng nhé!}''

Ayame đứng trên bục giảng, tay gõ thước xuống bàn đầy dứt khoát. Cạnh cô, Wata đang mỉm cười hiền hậu, tay cầm một xấp tài liệu dày cộp. Ngồi ở bàn giáo viên, Yuko-sensei chống cằm quan sát chúng tôi với ánh mắt trìu mến, thỉnh thoảng lại gật đầu khích lệ.

Tôi khẽ liếc nhìn Suzuki đang ngồi cạnh Naomi. Cả hai đều đang nhìn lên bảng với vẻ mặt đầy tò mò.

``Onii-sama này,'' Suzuki ghé sát tai tôi, thì thầm. ``\hl{Em có chút thắc mắc. Chẳng phải lễ hội trường học trong các truyện tranh em đọc thường tổ chức vào mùa thu, khoảng tháng 10 hay tháng 11 sao? Tại sao ở Aogami lại tổ chức ngay giữa mùa hè thế này, lại còn ngay trước Lễ hội mùa hè của thị trấn nữa?}''

Tôi cũng đang tự hỏi điều tương tự. Ở thế giới cũ của chúng tôi, quy luật này gần như là bất di bất dịch. ``\hl{Anh cũng đang thắc mắc y chang em đây. Nhưng mà nghĩ kỹ lại,}'' tôi đáp lại, ``\hl{nếu ở thế giới này, ngày Lễ tình nhân rơi vào tháng Năm thì có thể Lễ hội trường học rơi vào tháng Sáu, tháng Bảy cũng không có gì lạ, phải không?}''

``\hl{Thì đúng thế, nhưng mà\dots}''

Nghe thấy tiếng thì thầm của hai chúng tôi, Hội trưởng cười nhếch mép, vẻ đầy đắc thắng: ``Câu hỏi hay đấy, Hikawa-san! Đây chính là điểm đặc biệt chỉ có ở Aogami. Lễ hội trường học của chúng ta được đẩy lên sớm để tạo thành một chuỗi sự kiện `Đại Lễ Hội Aogami' kéo dài hai ngày liên tiếp!''

``Đại Lễ hội Aogami?'' chúng tôi đồng thanh hỏi lại.

Touka-san, lúc này đang lật giở một cuốn sách lịch sử của trường, tiếp lời với giọng điệu điềm tĩnh thường ngày: ``Đúng vậy. Theo truyền thống của thị trấn, việc gộp Lễ hội của trường và Lễ hội mùa hè của đền Aogami thành một sự kiện lớn sẽ tạo ra sự cộng hưởng linh khí mạnh mẽ nhất. Hơn nữa, vì nơi này từng là một thị trấn nhỏ, việc tập trung mọi nguồn lực vào một thời điểm sẽ khiến lễ hội trở nên hoành tráng và thu hút nhiều du khách hơn, và truyền thống đó kéo dài tới tận bây giờ. Đó là sự kết hợp giữa hiệu quả kinh tế và tín ngưỡng tâm linh lâu đời.''

Tôi khẽ gật đầu, thầm cảm thán sự độc đáo của thế giới này.

``\vie{Ra là vậy. Một lễ hội kép sao? Nghe có vẻ sẽ tốn rất nhiều sức lực đây,}'' tôi lẩm bẩm, bàn tay vô thức chạm vào mặt đồng hồ.

Suzuki mỉm cười, đôi mắt lấp lánh: ``\vie{Nhưng mà em thích ý tưởng này. Nó giống như một buổi đại tiệc kéo dài không nghỉ vậy. Nếu đây là mùa hè cuối cùng\dots\ thì chúng ta nên làm cho nó thật rực rỡ chứ nhỉ?}''

Nhìn nụ cười của em và sự hào hứng của cả lớp, tôi cảm nhận được một luồng nhiệt huyết đang lan tỏa. Đúng vậy, bất kể thời gian có lệch lạc so với những gì chúng tôi từng biết, thì sự chân thành và niềm vui của những người xung quanh là hoàn toàn thực tế.

``Vậy thì, hỡi những con dân của lớp 1-C!'' Ayame vung tay, giọng đầy kịch tính. ``Hãy cho ta thấy sự sáng tạo của các bạn đi! Chúng ta nên làm gì để biến lễ hội văn hóa năm nay thành một huyền thoại đây?''

Cả lớp ngay lập tức bùng nổ. Những cánh tay giơ lên liên hồi kèm theo những tiếng hò hét phấn khích.

``\textbf{Thưa Hội trưởng!}'' Shion đứng bật dậy, chiếc khăn quàng cổ (dù trời đang nóng) tung bay theo một vũ điệu kỳ quặc. ``Ta đề xuất biến lớp học thành `Đền thờ Hắc ám của Vương quốc Bóng tối'! Chúng ta sẽ triệu hồi những linh hồn vất vưởng và ban phát lời nguyền cho khách tham quan!''

``Bác bỏ!'' Hội trưởng phũ phàng dập tắt ý tưởng ngay lập tức. ``Lễ hội là để vui chơi, không phải để trù ẻo khách, Shion-san.''

Shion nghe vậy liền xụ mặt xuống, khoanh tay trước ngực với vẻ dỗi hờn: ``Hừ, các ngươi thật chẳng có chút nhãn quan nghệ thuật nào cả. Bóng tối là khởi nguồn của mọi sức mạnh, vậy mà lại bị coi nhẹ thế này sao?''

Kaigou đang ngồi cạnh tôi khẽ thở dài, buông một lời bình luận đầy thực tế: ``Nếu cậu biến lớp học thành đền thờ hắc ám và bắt đầu ban phát lời nguyền thật, tôi e là cảnh sát thị trấn sẽ ghé thăm chúng ta trước khi lễ hội kết thúc đấy, Shion-san.''

Suzune giơ tay, đôi mắt lấp lánh như thể vừa tìm thấy kho báu: ``Mình muốn mở một quầy bánh bao hấp! Bánh bao nhân thịt, nhân đậu đỏ, nhân cà ri\dots\ mình có thể giúp nếm thử mọi loại nhân để đảm bảo chất lượng!''

Touka nghe vậy liền thở dài, buông một câu xanh rờn: ``Nếu cậu làm chủ quầy, mình e là chúng ta sẽ phá sản trước khi bán được cái bánh đầu tiên vì `chi phí nếm thử' của cậu đấy, Suzune-san.''

Suzune phồng má, hai tay chống hông tỏ vẻ hậm hực như một đứa trẻ bị từ chối kẹo. Touka thấy vậy liền khẽ xoa đầu cô bạn, khiến Suzune nhanh chóng quên luôn nỗi buồn mà cười hì hì trở lại.

Yuko-sensei khẽ cười hiền hậu: ``Tôi thấy ý tưởng của Suzune-chan rất đáng yêu, nhưng có lẽ chúng ta cần một cái gì đó có quy mô hơn cho một lớp học đầy năng lượng như 1-C chúng ta.''

Bầu không khí trong lớp 1-C lúc này náo nhiệt hơn bao giờ hết, những tiếng bàn tán xôn xao vang lên không ngớt. Có người đề xuất mở trung tâm trò chơi điện tử, người lại muốn làm một buổi trình diễn thời trang cosplay, nhưng dường như chưa có ý tưởng nào đủ sức thuyết phục được số đông như mong đợi.

Sau đó, các nhóm câu lạc bộ bắt đầu đưa ra kế hoạch riêng của mình để phối hợp với lớp.

``Nhóm Âm nhạc bọn tớ (Azuki, Honoka và Yuki) đã chuẩn bị xong danh sách các bài hát rồi,'' Honoka hớn hở nói. ``Nhóm sẽ biểu diễn ở sân khấu chính, nhưng nếu lớp có gian hàng, bọn tớ có thể biểu diễn live để thu hút khách!''

``Mỹ thuật thì có bọn này lo phần trang trí nội thất và vẽ biển hiệu nhé!'' Hotaru và Mari đồng thanh, tay đã bắt đầu múa may như thể đang vẽ trên không trung.

Hanabi, Miku và Mitsuki từ câu lạc bộ Thời trang thì tự tin khẳng định: ``Dù lớp chọn làm gì, trang phục cứ để chúng ta lo. Chúng ta sẽ thiết kế những bộ đồ độc nhất vô nhị!''

Không khí thảo luận càng lúc càng nóng hơn cả cái nắng ngoài kia. Cuối cùng, Nanase - người nãy giờ vẫn im lặng quan sát - khẽ hắng giọng.

``Tôi có một đề xuất thế này,'' cô nói, giọng điềm tĩnh nhưng đầy uy lực khiến cả lớp im bặt. ``Chúng ta có thể kết hợp nhiều yếu tố lại với nhau. Phòng học chính của chúng ta sẽ biến thành một `Café Hầu gái và Quản gia' (Maid \& Butler Café). Với sự khéo tay của nhóm Touka-san và Saki-san, menu sẽ rất tuyệt vời. Còn phòng chuẩn bị kế bên, chúng ta sẽ biến nó thành một `Nhà ma mini' do nhóm của Shion-san và Sakura-san phụ trách. Vừa có chỗ nghỉ chân ăn uống, vừa có trò cảm giác mạnh, chắc chắn sẽ rất hút khách.''

Shion nghe xong liền chống cằm suy nghĩ, rồi đôi mắt cô nàng bỗng rực sáng lên đầy phấn khích. ``Ồ, một mê cung sợ hãi nằm ngay cạnh một quán trà thư giãn sao? Sự tương phản giữa nỗi khiếp nhược và sự yên bình\dots\ ý tưởng này thực sự mang đậm phong vị của một phù thủy bậc thầy đấy! Ta chấp nhận!''

Wata gật đầu tán thành: ``Ý tưởng của Nanase-chan rất toàn diện. Nó tận dụng được thế mạnh của hầu hết các nhóm trong lớp chúng ta.''

``Đúng vậy!'' Ayame đập bàn. ``Sự kết hợp giữa ánh sáng và bóng tối, giữa sự phục vụ tận tình và nỗi sợ hãi tột cùng! Tôi thích nó! Cả lớp thấy sao?''

Một tràng pháo tay giòn giã vang lên thay cho câu trả lời. Yuko-sensei mỉm cười gật đầu: ``Vậy là đã chốt nhé. Lớp 1-C sẽ tổ chức `Café \& Nhà ma'. Bây giờ, chúng ta sẽ đi vào phần quan trọng nhất: phân công công việc cụ thể cho từng người.'' Cô chủ nhiệm gõ nhẹ lên bảng để thu hút sự chú ý của mọi người một lần nữa.

Ánh nắng chiều rực rỡ chiếu xuyên qua những ô cửa kính, tạo nên những vệt sáng dài lấp lánh trên sàn gỗ lớp học. Bầu không khí im lặng bao trùm trong giây lát, chỉ còn nghe tiếng ve sầu kêu râm ran từ phía rặng thông ngoài cửa sổ.

``Được rồi, bây giờ chúng ta sẽ bắt đầu chia nhóm công việc cụ thể cho dự án `Café \& Nhà ma' của lớp 1-C chúng ta nhé.''

Ayame cầm viên phấn mới, bắt đầu ghi tên các hạng mục chính lên bảng với một nụ cười đầy bí hiểm. Cô nàng dường như đã âm thầm chuẩn bị sẵn một danh sách các ``nạn nhân'' trong đầu từ trước. ``Đầu tiên là nhóm Quản gia cho quán cà phê. Ai có đủ tự tin để đảm nhận trọng trách này nào?''

Cả lớp bỗng chốc im bặt, những ánh mắt tinh quái bắt đầu đổ dồn về phía cậu nam sinh hiếm hoi trong lớp. Đúng rồi đấy, là tôi.

Tôi thở dài một tiếng thật khẽ, thầm tự nhủ rằng dù mình có ngồi im thin thít thì kiểu gì cũng bị lôi vào cuộc thôi.

Thế là, trước khi bất kỳ ai kịp lên tiếng chỉ định, tôi đã giơ tay lên một cách vô cùng dứt khoát: ``Để tôi làm quản gia cho. Đằng nào thì mọi người cũng định gọi tên tôi thôi đúng không?''

Cả lớp bỗng ồ lên một tràng đầy ngạc nhiên, ngay cả Hội trưởng cũng phải nhướng mày vì sự chủ động bất ngờ này của tôi. ``Ồ, Kuon-kun hôm nay hăng hái quá nhỉ? Cứ ngỡ cậu sẽ tìm đủ mọi cách để trốn tránh như mọi khi chứ.''

``Thay vì bị các bạn ép vào những vai diễn kỳ quặc hơn, tôi thà tự chọn vai mình cảm thấy ít phiền phức nhất,'' tôi đáp lại với vẻ mặt thản nhiên nhất có thể, dù trong lòng thầm biết rõ mình chẳng thể nào thoát được số phận.

Ayame cười khoái chí, nhanh tay ghi tên tôi vào mục Quản gia với những nét chữ to và vô cùng rõ ràng. ``Được thôi, tinh thần tự giác rất đáng khen! Vậy Quản gia sẽ gồm có Kuon-kun và\dots\ Kaigou-chan nhé!''

Tôi khựng lại một nhịp khi nghe tên cô nàng mắt xanh đó bị xếp vào nhóm Quản gia. Chưa kịp để tôi hay bất cứ ai lên tiếng, Kaigou đã khẽ nhíu mày, lên tiếng đầy dứt khoát: ``Khoan đã, Hội trưởng. Nếu phải chọn giữa Quản gia và Hầu gái, tôi sẽ chọn Hầu gái.''

Cả lớp bỗng chốc lặng đi. Ngay cả Ayame cũng phải tròn mắt ngạc nhiên: ``Ồ? Kaigou-chan mà lại chủ động muốn làm hầu gái sao? Tôi cứ ngỡ cậu sẽ thích phong cách quản gia mạnh mẽ hơn chứ?''

Cô nàng tóc đuôi ngựa thản nhiên khoanh tay, đưa mắt nhìn ra ngoài cửa sổ như thể điều cô vừa nói là hiển nhiên nhất trên đời. ``Chỉ là tôi thấy mặc váy thoải mái hơn là mặc quần dài thôi. Hơn nữa, mặc đồ hầu gái cũng không có gì tệ, miễn là nó không quá rườm rà.''

Tôi hắng giọng, nhìn cô nàng với vẻ hơi lo ngại. ``Cô chắc chứ? Cô nên nhớ là khách hàng của chúng ta không chỉ có các học sinh trong trường đâu. Lễ hội này mở cửa cho cả người dân thị trấn, và tệp khách hàng ghé quán café hầu gái thường là nam giới đấy.''

Tôi khẽ nghiêng người về phía cô, nói nhỏ: ``Tôi biết cô vốn chẳng ưa gì đám con trai, nhưng làm ơn, nếu không thể nhỏ nhẹ thì ít nhất đừng quá thô lỗ với khách, kẻo lại mang tiếng đấy.''

Ayame nghe thấy vậy liền cười phá lên, điệu cười đầy vẻ ranh mãnh như thể vừa bắt được thóp của ai đó. ``Đúng đấy Kaigou-san! Tôi vẫn chưa quên ở hội thao lần trước, cậu đã nhìn đội nam bên kéo co đáng thương thế nào đâu. Ánh mắt của cậu lúc đó cứ như muốn đóng băng luôn cả sợi dây thừng lẫn đối thủ của mình vậy!''

Kaigou khoanh tay trước ngực, quay mặt đi chỗ khác với vẻ mặt không hề có chút thay đổi hay lay chuyển nào: ``Tôi chỉ làm những gì cần thiết để mang lại chiến thắng và vinh quang cho đội chúng ta mà thôi. Việc phải đi phục vụ và tỏ ra ngọt ngào với đám đàn ông cặn bã đó chưa bao giờ nằm trong từ điển của tôi cả.''

\dots Cô nàng ghét đàn ông tới mức đó sao. Nhất là sau những gì mà đám côn đồ kia đã làm với em gái của cô ấy. Tôi hiểu được cảm xúc của Kaigou lúc đó, nhưng thực sự nếu cứ giữ khư khư suy nghĩ như vậy thì sẽ chẳng tốt lành gì. Dù thế nào thì đây cũng là lần đầu tiên mà lớp chúng ta có một sự kiện lớn như vậy. Không thể để nó đổ bể chỉ vì những lý do cá nhân được.

Tình hình có vẻ như đang rơi vào bế tắc cho đến khi Touka bỗng nhiên lên tiếng với vẻ mặt cực kỳ điềm tĩnh. Vị thần đồng của lớp ta dường như vừa nảy ra một ý tưởng mang tính đột phá từ những quan sát nhạy bén của mình.

``Thực ra, nếu Kaigou-san cảm thấy khó khăn trong việc thay đổi thái độ, chúng ta có thể biến nó thành một phong cách.''

Cả lớp bỗng đổ dồn mọi ánh mắt về phía Touka, tò mò không biết cô nàng định nói điều gì mang tính vĩ mô tiếp theo.

Cô nàng tóc xanh tuyết chậm rãi đứng dậy, ánh mắt ánh lên vẻ tinh ranh của một người nắm bắt được những xu hướng tâm lý kỳ lạ: ``Ở các thành phố lớn hiện nay, có một kiểu quán café hầu gái nơi nhân viên sẽ vào vai những cô nàng kiêu kỳ. Họ đối xử lạnh nhạt, thậm chí là có phần `độc hại' và khinh thường đối với chính khách hàng của mình.''

Cả lớp bỗng chìm vào im lặng mất vài giây để cố gắng tiêu hóa lượng thông tin có phần dị hợm này. Cô tiếp tục, giọng có chút nghiêm túc hơn: ``Và điều bất ngờ là, có một lượng lớn khách hàng nam giới lại cực kỳ yêu thích cảm giác được phục vụ như thế đấy.''

Tôi và Suzuki đứng chết lặng, nhìn Touka như thể cô vừa nói ra một thứ ngôn ngữ đến từ một hành tinh xa xôi nào đó. Bọn tôi vừa nghe cái gì thế?

``\vie{Cái gì cơ? Trên đời này lại thực sự tồn tại những người có sở thích kỳ quặc đến mức đó sao?}'' tôi lẩm bẩm trong miệng, cảm thấy kiến thức về xã hội của mình dường như vẫn còn quá nhỏ bé và hạn hẹp.

Suzuki thì ôm lấy đầu mình, giọng rên rỉ vì sự sụp đổ hoàn toàn của hình tượng hầu gái trong sáng trong tâm trí: ``Touka-san, mình không nghĩ là ở một thị trấn yên bình như Aogami lại có người thích bị mắng mỏ đâu\dots''

Nhưng Hội trưởng thì ngược lại hoàn toàn, cô nàng vỗ tay bộp một cái, đôi mắt sáng rực lên như vừa đào được vàng. ``\textbf{Thiên tài! Touka-san đúng là một thiên tài bẩm sinh!}'' cô hét lên đầy vẻ phấn khích và hào hứng. ``Vậy là chốt phương án nhé! Kaigou-san sẽ là `Hầu gái Nữ vương' của quán café chúng ta! Cậu cứ việc giữ nguyên cái nhìn băng giá đó, thậm chí là hãy thể hiện sự khinh miệt ra mặt với khách nam đi!''

Kaigou khẽ nhướng mày, có vẻ như vẫn chưa hiểu hết sự tình nhưng thấy mình không cần phải ép buộc bản thân nên cũng ổn.
``Nếu chỉ cần đứng đó và tỏ ra khó chịu mà vẫn mang lại hiệu quả công việc, thì tôi sẽ tham gia.'' Tôi chỉ biết thở dài ngao ngán, thầm cầu nguyện cho những vị khách nam xấu số nào đó sẽ vô tình ghé thăm gian hàng lớp mình.

Tiếp theo, Ayame bắt đầu phân vai cho các hầu gái còn lại để tạo nên một sự đa dạng đặc sắc cho quán.

``Suzuki-chan, với gương mặt ngây thơ và giọng nói đó, cậu chắc chắn phải vào vai `em gái nhỏ' rồi! Cậu sẽ gọi khách là `Onii-chan' với một giọng điệu nũng nịu và dễ thương nhất có thể cho tôi nhé!''

Suzuki nghe đến đó liền đỏ bừng mặt từ tận chân tóc đến tận mang tai, tay chân quơ quạn rối rít như người mất trí: ``Không đời nào! Chuyện đó thật sự quá sức xấu hổ! Làm sao mình có thể gọi một người lạ như thế được chứ!''

Dĩ nhiên, cô chị Kaigou là người phản ứng dữ dội nhất. Cô đứng bật dậy, đập bàn một cái rầm khiến cả lớp giật mình, rồi nói như khạc ra lửa: ``Tuyệt đối không! Tôi không đời nào cho phép Suzuki phải đóng cái vai nũng nịu, ưỡn ẹo đó trước mặt bất cứ gã đàn ông nào cả!''

Tôi nhìn vào đôi mắt đang bốc lửa của cô ấy, rồi lại nhìn sang Suzuki đang bối rối cúi đầu. Tôi hiểu cảm giác đó. Sau những gì đã xảy ra với em ấy ở những vòng lặp trước, hay thậm chí là vụ rắc rối với đám côn đồ vừa rồi, việc Kaigou muốn bao bọc em gái mình trong một lớp vỏ thép là điều hoàn toàn dễ hiểu. Nhưng nếu cứ như vậy, Suzuki sẽ mãi chẳng thể thực sự hòa nhập được với mọi người.

Tôi hắng giọng, nhìn thẳng vào Kaigou: ``Này, cô bình tĩnh lại chút đi.''

Cô nàng quay sang lườm tôi: ``Cậu bảo tôi bình tĩnh sao, Kuon-kun? Cậu cũng hiểu rõ tại sao tôi lại phản đối mà, đúng không?''

``Vì đây là lớp 1-C chứ không đơn thuần chỉ trong gia đình nữa,'' tôi đáp lại bằng giọng điềm tĩnh nhất có thể. ``Ở đây, Suzuki không chỉ là em gái của cô, mà còn là một thành viên của tập thể này. Cô nhìn xem, chẳng phải em ấy cũng muốn đóng góp cho lớp sao? Nếu cô cứ cấm cản vì nỗi sợ của riêng mình, chẳng phải cô đang vô tình tách em ấy ra khỏi bạn bè à?''

Tôi tiến lại gần thêm một bước, hạ thấp giọng: ``Hơn nữa, tôi sẽ làm quản gia ngay tại đó. Tôi cam đoan với cô, nếu có bất cứ tên nào dám có hành động hay lời nói khiếm nhã với Suzuki, tôi sẽ là người tống khứ hắn ra khỏi cửa đầu tiên. Cô không tin tôi sao?''

Kaigou khựng lại, ánh mắt cô dao động giữa sự bảo bọc cực đoan và sự lưỡng lự. Những lời của tôi dường như đã chạm đúng vào điểm mấu chốt. Các bạn trong lớp cũng bắt đầu lên tiếng ủng hộ.

``Đúng đấy Kaigou-san! Bọn mình sẽ luôn ở cạnh Suzuki-chan mà!'' Honoka hăng hái tiếp lời. Hotaru cũng gật đầu chắc nịch: ``Tôi sẽ lo phần trang phục sao cho vừa dễ thương lại vừa kín đáo, cậu yên tâm nhé.''

Trước áp lực từ cả lớp và cái nhìn kiên định của tôi, Kaigou thở dài một tiếng thườn thượt, rồi ngồi xuống với vẻ mặt vẫn còn chút hậm hực. ``Được rồi\dots\ tôi sẽ quan sát thật kỹ. Chỉ cần một tên nào đó đi quá giới hạn, đừng trách tôi tại sao quán café lại biến thành võ đài.''

Suzuki khẽ ngước lên, đôi mắt lấp lánh nhìn tôi đầy biết ơn. Tôi chỉ khẽ gật đầu, thầm nghĩ rằng vai trò quản gia của mình có lẽ sẽ bận rộn hơn tôi tưởng nhiều đây.

``Honoka-chan sẽ đảm nhận vai hầu gái năng động, thỉnh thoảng hơi vụng về để tạo tiếng cười nhé!'' Hội trưởng tiếp tục phân vai. Cô nàng bờm Cáo sa mạc xem chừng cũng tâm đắc với vai này. Ayame ghi chép liên hồi lên bảng, danh sách các `hình mẫu' hầu gái bắt đầu lấp đầy những khoảng trống còn lại. ``Suzune-chan sẽ là hầu gái ngây thơ ham ăn, vừa bưng bê phục vụ vừa lén lút ăn vụng bánh ngọt của khách. Hanabi-chan sẽ vào vai hầu gái lôi cuốn, và Miku-san sẽ là một hầu gái trưởng thành, thanh lịch và quý phái.''

Mỗi người mang một vẻ đẹp và phong cách riêng, tạo nên một đội hình hầu gái đầy màu sắc mà chắc chắn sẽ không ai bỏ qua. Lớp học bỗng chốc trở thành một xưởng sản xuất những hình mẫu nhân vật kinh điển thường thấy trong các bộ anime.

Trong khi nhóm café đang vô cùng sôi nổi, nhóm Nhà ma cũng bắt đầu lên kế hoạch hành động dưới sự chỉ đạo của Shion. Cô nàng phù thủy tự xưng đứng hẳn lên ghế, vung tay đầy kịch tính như đang đọc một bài diễn văn triệu hồi hắc ám. ``Hỡi những đồng đội trung thành của bóng tối! Chúng ta sẽ tạo ra một `Mê cung Đêm trường Vĩnh cửu' tại đây! Nơi mà mỗi bước đi sẽ là một sự tra tấn tinh thần khủng khiếp đối với những kẻ phàm trần gan lì nhất thị trấn!''

Yuka, người thường ngày luôn trông như sắp chìm vào giấc ngủ ngàn thu, bỗng mở mắt ra với ánh nhìn đầy sắc lẹm: ``Tôi có thể giúp chỉ ra những góc chết và những nơi mà cảm giác rùng rợn luôn hiện hữu rõ rệt nhất trong căn phòng đó\~{} Tôi sẽ sử dụng năng lực đặc biệt để đặt những `cái bẫy' tâm lý đúng vào lúc khách hàng đang sơ hở nhất\~{}''

Sakura cũng vô cùng hào hứng tham gia, trên tay cô là một cuốn sổ ghi chép đầy những truyền thuyết kinh dị địa phương. Cô nói với một tông giọng của một kẻ thích kể chuyện ma: ``Tôi sẽ đảm nhận phần không khí! Chúng ta sẽ lồng ghép những câu chuyện rùng rợn nhất về cánh rừng Aogami. Tiếng gió hú lạnh lẽo, tiếng than khóc oán hận của những linh hồn bị bỏ lại\dots\ đảm bảo họ sẽ chạy tới mức khóc thét luôn!''

Aku, người nãy giờ vẫn cố gắng thu mình lại một góc nhỏ bé, bỗng bật dậy và hét lên một tiếng thất thanh đầy hoảng loạn: ``Không! Tuyệt đối không bao giờ! Mình không muốn bước chân vào cái căn phòng đáng sợ đó đâu!'' Cô nàng hậu đậu tội nghiệp bắt đầu chạy quanh lớp học để tìm đường thoát thân nhưng mọi lối ra đều đã bị chặn lại.

Shion khoác vai cô bạn thân với một nụ cười gian xảo mang đậm phong cách ác ma khiến Aku càng thêm run rẩy dữ dội: ``Yên tâm đi mà Aku-chan, ngươi sẽ có một vai diễn vô cùng quan trọng và không thể thay thế trong mê cung này. Ngươi sẽ là `người dẫn đường tội nghiệp' luôn bị những con ma đáng sợ đuổi sát sau lưng xuyên suốt cả mê cung! Tiếng hét thất thanh chân thực nhất của ngươi chính là nhạc nền sống động nhất cho nhà ma của chúng ta đấy!''

Aku nghe xong liền suy sụp hoàn toàn, cô ngồi bệt xuống sàn nhà với đôi mắt vô hồn nhìn vào khoảng không vô định. Nhìn cô bạn Hầu gái đang run lên như cầy sấy, hai hàng lệ rơi xuống, tôi chỉ có thể cười chát với tình cảnh éo le của cô, đành lên tiếng động viên: ``Thôi nào Shion-san, đừng dọa cô ấy nữa. Aku-san, chúng ta hãy cùng cố gắng nhé.''

Đáp lại lời của tôi là cái gật đầu khiên cưỡng của Aku. Sau khi đã chuẩn bị xong, phân vai đầy đủ, nhóm nhà ma cũng bắt đầu kế hoạch của mình. Họ di chuyển đến một căn phòng trống, nơi được chọn để trở thành ``Mê cung Đêm trường Vĩnh cửu'' của họ.

Cũng hơi tiếc khi Aku lại không ở trong nhóm café hầu gái. Cô ấy thực sự là phù hợp nhất với những chuyện kiểu như thế này - chỉ trừ đúng tính hậu đậu và vụng về giống như phiên bản gốc mà thôi.

Trong lúc đó, Hotaru từ nhóm Mỹ thuật cũng đã hoàn thiện những bản phác thảo cuối cùng cho chiếc áo đồng phục lớp. Cô nàng giơ cao tờ giấy vẽ lên để tất cả mọi người trong lớp đều có thể nhìn thấy rõ ràng từng chi tiết nhỏ nhất. ``Đây là mẫu thiết kế mà tớ đã dành rất nhiều tâm huyết. Tớ muốn nó vừa mang nét hiện đại vừa đậm tính biểu tượng.''

Chiếc áo đồng phục lớp sẽ có màu xanh hải quân sang trọng, với điểm nhấn là những đường viền trắng tinh tế ở cổ và tay áo. Ở phía ngực trái là logo hình hoa cẩm tú cầu được cách điệu một cách tài tình, lồng ghép bên trong một chiếc đồng hồ cát. Biểu tượng đó gợi nhắc cho chúng ta về sự chuyển giao của thời gian và vẻ đẹp dịu dàng nhưng đầy kiên cường của lớp 1-C.

Mọi người trong lớp đều trầm trồ khen ngợi không ngớt, ai nấy đều cảm thấy tự hào khi sắp được khoác lên mình chiếc áo đó.

``Đẹp quá Hotaru-san! Màu xanh này mặc lên chắc chắn sẽ tôn dáng và trông vô cùng chuyên nghiệp đấy!'' Hanabi thốt lên.

Saki cũng gật đầu tán thưởng sự tinh tế trong việc lựa chọn màu sắc chủ đạo và các biểu tượng ý nghĩa của cô bạn mình.

Từng hạng mục công việc đã dần được lấp đầy trên bảng đen, không khí bận rộn bắt đầu bao trùm lấy tất cả mọi người.

Cuối cùng, Naomi nãy giờ vẫn im lặng theo dõi mọi diễn biến từ đầu, bỗng khẽ nghiêng đầu hỏi một câu đầy ngây ngô: ``Vậy còn em? Em sẽ đóng vai trò gì trong những kế hoạch lớn lao này của mọi người ạ?''

Ánh mắt của toàn bộ lớp học bỗng chốc đổ dồn về phía cô bé mắt xanh với một sự quan tâm vô cùng đặc biệt. Cô bé luôn mang một vẻ đẹp rất khác biệt, một sự thản nhiên và pha chút bí ẩn luôn thu hút mọi ánh nhìn xung quanh.

Ayame bước lại gần, đặt tay lên vai cô bé nhỏ nhắn với một nụ cười rạng rỡ và ấm áp nhất từ đầu buổi họp đến giờ: ``Naomi-chan sẽ là `Kanban-musume' - cô bé biểu tượng đại diện cho toàn thể lớp 1-C chúng ta tại lễ hội lần này! Em sẽ mặc bộ trang phục lộng lẫy nhất, làm nhiệm vụ chào đón tất cả các khách tham quan. Sự hiện diện của em chính là lời mời gọi tuyệt vời nhất khiến không ai có thể nỡ lòng đi ngang qua gian hàng lớp mình.''

Tôi khẽ nhíu mày, không khỏi lo lắng lên tiếng: ``Khoan đã, Ayame-san. Để một đứa trẻ như Naomi đứng ở những chỗ đông người liệu có an toàn không? Lễ hội sẽ cực kỳ đông đúc, lỡ em ấy bị lạc hay quá mệt thì sao?''

Suzuki cũng gật đầu liên tục, vòng tay ôm lấy em như thể sợ ai bắt mất: ``Đúng đó. Em ấy vẫn còn nhỏ. Chưa kể sẽ có rất nhiều người lạ xung quanh nữa.''

Nhưng trái với sự lo lắng của chúng tôi, Naomi chỉ khẽ nghiêng đầu. Đôi mắt lục bảo của cô bé ánh lên những tia sáng kiên định hiếm thấy: ``Chào đón khách hàng và thu hút mọi ánh nhìn sao? Em muốn làm công việc đó.''

``Em chắc chứ, Naomi?'' tôi hỏi lại.

Cô bé gật đầu nhẹ nhàng, vẻ mặt thản nhiên nhưng vô cùng chân thành: ``Vâng. Mọi người đã giúp đỡ em rất nhiều để em có thể ở lại đây. Em cũng muốn được đóng góp cho lớp 1-C. Hơn nữa, em không còn yếu ớt như trước nữa đâu.''

Nghe những lời từ cô bé nhỏ nhắn, Hotaru mỉm cười tiến lại gần, ánh mắt ánh lên sự yêu chiều với những thứ dễ thương: ``Nếu Kuon lo lắng về sự an toàn, chúng ta có thể thiết kế một góc đón khách riêng, trang trí thật đẹp mắt ngay cạnh gian hàng để Naomi-chan ngồi thay vì phải đứng ngoài cổng hay đi lại nhiều.''

Mitsuki cũng gật đầu đồng tình, giọng nói dịu dàng nhưng đầy chắc chắn: ``Đúng thế, na. Thêm nữa, kể cả khi em ấy nhận việc, thì cũng không thể để Naomi-chan ở một mình được. Cần phải có người thay phiên nhau trông chừng và chăm sóc em ấy suốt thời gian diễn ra lễ hội.''

Saya khẽ chắp tay lại, lên tiếng với nhịp điệu chậm rãi và êm ái đặc trưng của mình: ``Vậy thì để bọn mình lo việc đó cho\~{} Mình, Hotaru-san và Mitsuki-san\dots\ sẽ luân phiên nhau ở cạnh Naomi-chan\~{} Bọn mình sẽ đảm bảo em ấy luôn được an toàn\~{}''

Nghe những lời đề nghị chu đáo từ các bạn học, nỗi lo trong lòng tôi vơi đi quá nửa. Suzuki cũng thở phào nhẹ nhõm, mỉm cười vuốt tóc mềm mại của Naomi: ``Nếu có ba người các cậu giúp đỡ thì mình và Onii-sama yên tâm rồi.''

``Vậy là quyết định xong!'' Ayame vỗ tay cái bốp, chốt lại vấn đề. ``Naomi-chan sẽ là linh vật tối cao của `Café \& Nhà ma' lớp 1-C nha!''

Tiếng phấn viết bảng lách cách vang lên dồn dập khi Ayame nhanh tay lấp đầy những khoảng trống còn lại trong danh sách phân công. Touka và Kanade được giao trọng trách quán xuyến khu vực bếp, nơi Touka lập tức bắt đầu liệt kê danh mục nguyên liệu với một vẻ mặt vô cùng tập trung.

Phần kỹ thuật cho nhà ma được giao cho Koishin và Azuki, hai người họ bắt đầu thảo luận về các hiệu ứng ánh sáng và thời điểm hù dọa sao cho hiệu quả nhất. Kaoru và Kana đảm nhận vai trò phát thanh, cả hai đã sớm lên kế hoạch cho những thông báo sôi động để giữ cho không khí lễ hội luôn ở mức cao trào.

Với những công việc nặng nhọc và hậu cần, Yae và Shiina là sự lựa chọn hiển nhiên với Kaigou có thể hỗ trợ một phần nào đó; cả hai cùng giơ ngón tay cái ra hiệu sẵn sàng cho những thử thách thể lực phía trước. Akane và Uzuki kiểm tra lại máy ảnh, sẵn sàng ghi lại mọi khoảnh khắc đáng nhớ với tư cách là những nhiếp ảnh gia chính thức (nàng bờm Chó vẫn cần phải huấn luyện thêm một chút về kỹ năng chụp ảnh, nhưng có vẻ như cô nàng cũng sẵn sàng cho một mùa lễ hội đầy bùng nổ sắp tới).

Nhìn vào bảng phân công chi tiết với đầy đủ tên các bạn, tôi biết đây sẽ không phải là một lễ hội mùa hè đơn thuần. Chúng tôi đang cùng nhau viết nên những dòng đầu tiên của một bản trường ca mang tên ``Tuổi trẻ'', nơi mỗi cái tên, mỗi vai trò, dù lớn hay nhỏ, đều góp phần tạo nên một bản hòa ca trọn vẹn và đáng nhớ.

\ct

Buổi họp lớp kết thúc khi ráng chiều đỏ rực nhuộm kín bầu trời phía sau dãy núi Aogami. Sau khi dọn dẹp xong xuôi, bốn chúng tôi rời khỏi trường, cất bước trên con đường quen thuộc trở về nhà.

Naomi lon ton chạy lên phía trước vài bước, thỉnh thoảng lại dừng lại ngắm nghía những bông hoa dại ven đường. Bóng lưng nhỏ bé của em in trên mặt đường nhựa trải dài dưới ánh hoàng hôn. Ba chúng tôi đi chậm lại phía sau một chút để giữ khoảng cách, đủ để cô bé không nghe thấy những lời nói tiếp theo.

``\vie{Hôm nay mới chỉ xong phần của Lễ hội trường học thôi nhỉ, Kuon-kun,}'' Kaigou lên tiếng trước, khoanh tay trước ngực. ``\vie{Nghe Hội trưởng nói, ngày mai cậu ấy cùng mấy người trong Hội học sinh sẽ sang tận nhà chúng ta để bàn tiếp về Lễ hội thị trấn cơ đấy.}''

Suzuki khẽ gật đầu, nở một nụ cười nhẹ: ``\vie{Vâng. Gia đình Nanase-san là người đứng ra tổ chức Lễ hội thị trấn mà, chắc chắn sẽ còn nhiều việc phải làm. Cơ mà\dots\ Ayame-san sang nhà mình họp, hy vọng cậu ấy không bày trò quậy phá gì.}''

Tôi bật cười trước sự lo lắng của em gái. Nhìn về phía hình bóng nhỏ nhắn của Naomi đang vui vẻ bắt một chú bướm đêm vừa thức giấc, tôi khẽ nói, giọng chùng xuống: ``\vie{Nhưng mà, hai người thấy đó\dots\ Đây là lễ hội đầu tiên mà chúng ta tham gia kể từ khi đến thế giới này.}''

Kaigou hiểu ý tôi, ánh mắt cô cũng trở nên trầm ngâm, dõi theo Naomi: ``\vie{Và cũng là lễ hội cuối cùng của chúng ta ở đây. Thời gian không còn nhiều nữa.}''

``\vie{Lễ hội đầu tiên và cuối cùng\dots}'' Suzuki lặp lại, bàn tay khẽ siết chặt lấy quai cặp. ``\vie{Chính vì vậy, chúng ta phải làm cho nó thật rực rỡ, không để lại bất kỳ sự nuối tiếc nào. Không chỉ vì lớp 1-C, mà còn vì Naomi nữa.}''

Tôi gật đầu đồng tình. Quyết định rời đi không bao giờ là dễ dàng, nhất là khi chúng tôi đã có quá nhiều kỷ niệm và sự gắn kết với nơi này, mặc dù sau những biến cố xảy ra trước khi chúng tôi có được sự yên bình bây giờ. Nhưng dù kết cục có ra sao, thì hiện tại mới là điều mà tôi muốn trân trọng.

``\vie{Đúng vậy,}'' tôi quả quyết. ``\vie{Hãy dốc hết sức cho những ngày sắp tới. Lễ hội trường học, rồi Lễ hội thị trấn, và cả Nghi thức Đêm trăng rằm nữa. Chúng ta sẽ tạo ra một kỷ niệm đẹp nhất cho tất cả mọi người.}''

Naomi dường như nghe thấy tiếng thì thầm, cô bé quay lại, nghiêng đầu nhìn chúng tôi với đôi mắt lục bảo lấp lánh: ``\vie{Mọi người đang nói chuyện gì bí mật thế ạ?}''

Kaigou vội vàng xua tay, nở một nụ cười hiếm hoi: ``\vie{Không có gì đâu. Chị chỉ đang bàn xem tối nay chúng ta ăn gì thôi. Em thích ăn gì nào, Nao?}''

Cô bé mắt xanh mỉm cười rạng rỡ, chạy ngược lại nắm lấy tay cô chị lớn: ``\vie{Em muốn ăn trứng cuộn ạ!}''

Nhìn bóng dáng ba chị em họ vui vẻ bước đi dưới ráng chiều, tôi khẽ mỉm cười.

``\vie{Cố gắng lên nhé, mọi người,}'' tôi lẩm bẩm, như một lời hứa với chính bản thân mình và cả tập thể lớp 1-C. Một huyền thoại mới sắp được bắt đầu, và chúng tôi đã hoàn toàn sẵn sàng.

\begin{center}
  \uline{Ngày hôm sau.}
  \pov{Hikawa Kaigou}
\end{center}

Ánh nắng buổi sáng len lỏi qua những tán lá trắc bá, chiếu vào phòng khách nhà Hikawa chúng tôi. Tôi thở dài, khẽ chỉnh lại mấy tách trà trên bàn. Cũng giống như bao ngày cuối tuần khác, hôm nay, căn nhà của chúng tôi lại trở thành một điểm tụ tập của các bạn học lớp 1-C.

Ayame, Wata và Yuko-sensei đã có mặt từ sớm. Theo sau họ là một vài thành viên nòng cốt khác của lớp 1-C. Có vẻ như dư âm của buổi họp lớp hôm qua về Lễ hội trường học vẫn chưa hề hạ nhiệt, và giờ đây, chúng tôi lại tiếp tục bị cuốn vào vòng xoáy của một sự kiện còn lớn hơn: Lễ hội thị trấn Aogami.

``Chào mọi người! Rất xin lỗi vì đã làm phiền gia đình Hikawa vào ngày nghỉ thế này nhé!'' Hội trưởng cười rạng rỡ, nhưng đôi mắt thì vẫn ánh lên sự tinh quái thường ngày.

Trong khi đó, Hội phó khẽ cúi đầu xin lỗi: ``Bọn mình thực sự cần một không gian yên tĩnh để bàn bạc kỹ hơn về sự phối hợp giữa lớp và thị trấn. Mong mọi người thông cảm nhé.''

Yuko-sensei ngồi xuống ghế, mỉm cười nhìn chúng tôi: ``Tôi thấy rất mừng vì các em lại nhiệt tình như thế. Đây là cơ hội tốt để thắt chặt tình cảm của lớp chúng ta đấy.''

Tôi khẽ liếc nhìn Kuon, người đang ngồi ở góc phòng, tay cầm cuốn sổ ghi chép, vẻ mặt vẫn điềm tĩnh như mọi khi. Nao thì ngồi bệt dưới sàn, đôi mắt lục bảo tò mò nhìn xấp tài liệu mà Nanase đang trải ra trên bàn. Suzu thì loay hoay trong bếp để chuẩn bị thêm một ít bánh ngọt cho khách. Cảnh tượng này\dots\ có chút gì đó giống như một gia đình thực thụ đang chuẩn bị cho một sự kiện trọng đại.

Nanase hắng giọng, mở đầu cuộc họp với phong thái của một nữ doanh nhân thực thụ: ``Như mọi người đã biết, Lễ hội thị trấn mọi năm đều do công ty xây dựng của gia đình tôi thầu toàn bộ phần mặt bằng và hạ tầng, và năm nay cũng không ngoại lệ. Điều đó có nghĩa là lớp 1-C sẽ có những đặc quyền nhất định về vị trí đặt các gian hàng. Chúng ta cần những ý tưởng có thể thu hút khách du lịch từ các thị trấn lân cận, thay vì chỉ phục vụ học sinh như ở Lễ hội văn hóa ở trường.'' % yatai: gian hàng

Ngay lập tức, Honoka giơ tay cao vút, khuôn mặt tràn đầy sự phấn khích không hề giấu diếm: ``Tớ có đề xuất! Chúng ta nên tổ chức một cuộc 'Đua xe bò' quy mô toàn thị trấn! Người thắng cuộc sẽ được nhận một phiếu ăn miễn phí trọn đời tại gian hàng của lớp mình! Hãy tưởng tượng cảnh hàng chục con bò lao vun vút qua những con phố cổ của Aogami, thật là hào hùng biết bao!''

Cả căn phòng im bặt trong vài giây. Ayame ôm đầu, vẻ mặt như thể vừa nghe thấy một kế hoạch hủy diệt thế giới: ``Honoka-san, cậu định lấy đâu ra bò cho cả thị trấn đua? Và quan trọng nhất là ai sẽ đi dọn\dots\ hậu quả của chúng sau cuộc đua? Chưa kể đến việc phá hoại cảnh quan cổ kính của thị trấn nữa.''

Wata cười khổ, cố gắng đưa mọi thứ trở lại quỹ đạo: ``Honoka-san, ý tưởng này có vẻ hơi\dots\ khó thực thi về mặt vệ sinh, an toàn và cả tính nhân văn đối với động vật nữa.''

Chưa dừng lại ở đó, cô nàng Cáo sa mạc tiếp tục với một đề xuất còn táo bạo và điên rồ hơn: ``Vậy thì\dots\ 'Đại tiệc ném takoyaki nóng hổi'! Chúng ta sẽ thi xem ai ném takoyaki trúng miệng đối phương nhiều nhất! Vừa vui nhộn, vừa có thể tiêu thụ được một lượng lớn hàng tồn kho!''

Yuko-sensei vội vàng can thiệp trước khi Ayame tiếp tục ý kiến: ``Lãng phí thức ăn là không tốt đâu, Omori-san. Hơn nữa, bánh takoyaki nóng rất dễ gây bỏng cho khách tham quan. Chúng ta cần những ý tưởng an toàn hơn.''

Shion, lúc này đang quấn chiếc khăn quàng cổ dù trời khá nóng, đứng bật dậy với một điệu bộ kịch tính: ``Hỡi những kẻ phàm trần! Tại sao chúng ta không tổ chức một 'Lễ hội pháo hoa hắc ám'? Chúng ta sẽ nhuộm đen bầu trời bằng những làn khói tăm tối nhất, rồi từ đó, những tia sáng tím lạnh lẽo sẽ nổ tung, gieo rắc sự khiếp nhược rực rỡ vào lòng người dân thị trấn!''

``Cậu định biến lễ hội thành hiện trường vụ hỏa hoạn đấy à, Shion-san?'' Touka buông một lời bình luận sắc lẹm như thường lệ. ``Pháo hoa màu đen về cơ bản là không thể nhìn thấy trong đêm tối. Đó là một ý tưởng hoàn toàn phản khoa học.''

Suzune thì hào hứng đề xuất mở một sạp bánh bao hấp hình các linh vật thị trấn, nhưng ngay lập tức bị cô nàng Thần đồng dập tắt bằng một câu: ``Nếu cậu làm chủ sạp, mình e là chúng ta sẽ chỉ còn lại những cái lồng hấp không trước khi lễ hội bắt đầu vì cậu đã `kiểm tra chất lượng' hết rồi.''

Trong khi những ý tưởng kỳ quặc vẫn liên tiếp được đưa ra, một cuộc tranh luận khác gay gắt và mang tính định hướng hơn đã nổ ra giữa Sakura và Nanase. Đây là hai luồng tư tưởng hoàn toàn đối lập trong lớp chúng tôi. Cứ nhìn cách hai cô nàng dạy học cả lớp hồi ôn thi là biết.

``Chúng ta không thể biến lễ hội thành một khu thương mại thuần túy được, Nanase-san!'' Sakura đập tay xuống bàn, đôi mắt ánh lên vẻ kiên định của một người bảo vệ truyền thống. ``Đây là Lễ hội mùa hè của đền Aogami. Trọng tâm phải là các nghi thức tâm linh, diễu hành Mikoshi và những trò chơi dân gian mang tính truyền thống. Tôi đề xuất lớp chúng ta nên phụ trách khu vực viết lời ước và tổ chức một buổi triển lãm nhỏ về các truyền thuyết thần thoại của thị trấn. Đó mới là giá trị cốt lõi mà chúng ta cần gìn giữ.''

Nanase khoanh tay, đưa ánh mắt sắc sảo nhìn đối phương, không hề nhượng bộ: ``Sakura-san, cậu phải thực tế một chút đi. Lễ hội muốn hoành tráng, muốn có nhiều người tham gia thì phải có sự đổi mới và kinh phí. Việc lớp chúng ta đảm nhận các sạp trò chơi hiện đại với thiết kế bắt mắt và quà tặng hấp dẫn sẽ thu hút được nhiều khách du lịch hơn. Từ đó, chúng ta mới có thêm tiền để gây quỹ lớp và có một khoản đóng góp đáng kể cho việc tu sửa đền. Cậu có biết chi phí bảo trì cái cổng torii cũ kỹ kia là bao nhiêu không?''

``Tiền, tiền, cậu lúc nào cũng chỉ biết có hiệu quả kinh tế thôi sao?'' Sakura-san phản pháo, giọng cao lên một tông. ``Lễ hội là để gắn kết con người với thần linh, là dịp để bày tỏ lòng biết ơn, không phải là nơi để các cậu thực hành kỹ năng kinh doanh!''

``Cậu nói hay lắm, nhưng nếu không có sự đầu tư và tài trợ từ phía gia đình tôi thì cái sân đền của cậu giờ chắc vẫn còn đang rêu phủ đầy và tối tăm như một cái hang động đấy!'' Nanase-san cũng không vừa, đáp lại một cách đanh thép.

Không khí trong phòng khách bỗng chốc trở nên đặc quánh sự căng thẳng. Những thành viên khác của lớp bắt đầu xì xào bàn tán, chia thành hai phe rõ rệt. Inari thì loay hoay giữa hai người bạn học, cố gắng nói vài câu làm dịu tình hình nhưng giọng của cậu ấy hoàn toàn bị lấn át bởi sự cương quyết của hai cô nàng kia.

Tôi khẽ huých tay Suzu, ra hiệu cho em ấy cùng tôi bước ra phía hành lang gỗ sau nhà, nơi Kuon cũng đã lặng lẽ bước ra từ lúc nào để tránh khỏi mớ hỗn độn đó. Nao thì vẫn đang mải mê nhìn hai chị lớn cãi nhau với vẻ mặt thích thú, có lẽ cô bé coi đây là một phần của 'lễ hội chuẩn bị' chăng?

Ba chúng tôi đứng ở hành lang, nhìn ra khu vườn nhỏ xanh mướt. Tiếng ve sầu vẫn râm ran không dứt, hòa lẫn với tiếng tranh luận gay gắt bên trong phòng khách vọng ra.

``\vie{Suzu, Nao, hai đứa thấy sao về buổi họp này?}'' tôi khẽ hỏi, giọng đủ để ba người nghe thấy.

Suzu khẽ thở dài, dựa lưng vào cột nhà, đôi mắt thoáng chút trầm ngâm: ``\vie{Em không ngờ là chuẩn bị cho lễ hội thị trấn lại phức tạp và gây tranh cãi đến thế, Onee-sama. Mọi người đều có những cái lý riêng của mình, nhưng em thấy ai cũng đang thực sự dốc lòng vì tập thể cả. Chỉ là cách tiếp cận khác nhau quá thôi.}''

Kuon khẽ nhếch mép, đôi mắt đỏ nhìn về phía những rặng núi mờ xa phía chân trời: ``\vie{Sự nhiệt huyết đó đúng là rất đáng quý. Nhưng có một điều mà có lẽ chỉ ba chúng ta cảm nhận rõ nhất\dots\ đây là lần đầu tiên, và cũng chính là lần cuối cùng chúng ta được thực sự hòa mình vào một không khí chuẩn bị lễ hội náo nhiệt như thế này tại thị trấn Aogami.}''

Nghe Kuon nói, lòng tôi bỗng chùng xuống một nhịp đau nhói. Khoảng thời gian êm đềm này, hóa ra lại ngắn ngủi đến vậy. ``\vie{Đúng vậy. Lễ hội trường học rồi sẽ qua, Lễ hội thị trấn cũng sẽ kết thúc\dots\ mọi thứ diễn ra dồn dập như một bản nhạc cao trào trước khi hạ màn. Kuon-kun, cậu có lo lắng về Nghi thức vào đêm cuối cùng không?}''

Kuon im lặng một lúc lâu, rồi khẽ chạm vào mặt đồng hồ đang đeo trên tay, vẻ mặt trở nên nghiêm nghị: ``\vie{Game đã phát ra những cảnh báo về sự biến động năng lượng không gian. Nghi thức Đêm trăng rằm không chỉ đơn thuần là để ổn định linh hồn cho Nao, mà nó còn là chìa khóa then chốt để chúng ta có thể rời khỏi thế giới này. Tôi sẽ không cho phép bất cứ một sai sót nhỏ nào xảy ra.}''

Suzu nhìn vào phòng khách qua lớp cửa kính, nơi Ayame-san đang bắt đầu dẹp loạn bằng uy quyền của Hội trưởng: ``\vie{Em chỉ mong\dots\ những ngày sắp tới sẽ trôi qua thật chậm thôi. Em muốn được cùng các bạn ấy dựng nên những gian hàng tốt nhất, được cùng Nao đi dạo dưới ánh pháo hoa rực rỡ, được mặc yukata và cùng nhảy điệu Bon Odori\dots\ Em muốn lưu giữ tất cả những khoảnh khắc này vào sâu trong tim.}''

``\vie{Chúng ta sẽ làm được thôi, Suzu,}'' tôi khẽ xoa đầu em gái, cảm nhận sự mềm mại của tóc em. ``\vie{Hãy coi như tất cả những nỗ lực này là món quà cuối cùng chúng ta dành tặng cho Aogami, và cũng là món quà vô giá mà Aogami dành cho chúng ta trước khi chia xa.}''

Bên trong phòng, tiếng tranh luận đã bắt đầu dịu đi. Có vẻ như Yuko-sensei và Wata-san đã tìm ra một phương án thỏa hiệp trung lập, kết hợp cả yếu tố truyền thống lẫn hiện đại. Tôi nhìn Nao, cô bé bỗng quay đầu lại, đôi mắt lục bảo lấp lánh nụ cười rạng rỡ với ba chúng tôi.

Tôi hít một hơi thật sâu, cảm nhận mùi hương nồng nàn của cỏ cây mùa hè và sự náo nhiệt, đầy sức sống của tuổi trẻ đang bao trùm lấy ngôi nhà nhỏ vốn dĩ yên tĩnh này.

``\vie{Đi thôi. Chúng ta vào trong để xem Hội trưởng định 'ban phát' những trọng trách nào cho nhà Hikawa chúng ta nào,}'' Kuon lên tiếng, đôi mắt cậu đã lấy lại vẻ kiên định và sắc sảo thường lệ.

Khi chúng tôi bước trở lại phòng khách, bầu không khí đặc quánh lúc trước đã biến mất, thay vào đó là một sự đồng thuận đầy tích cực. Yuko-sensei đang mỉm cười gật đầu với bản phác thảo mới trên bàn.

``Mọi người đã thỏa hiệp xong rồi sao?'' cậu lên tiếng, đặt cuốn sổ ghi chép xuống.

Ayame cười rạng rỡ, chỉ vào bản đồ thị trấn: ``Đúng vậy! Nanase-san sẽ phụ trách khu vực sạp hàng hiện đại dọc theo con đường chính để đảm bảo kinh tế. Còn Sakura-san sẽ quán xuyến khu vực đường dẫn lên đền với các hoạt động truyền thống, Ema và nghi thức cầu nguyện. Và điểm kết nối giữa hai thế giới đó\dots\ chính là 'Con đường Linh hồn' - phiên bản nâng cấp của Nhà ma lớp 1-C do Shion-san đạo diễn!''

Nanase khoanh tay, vẻ mặt đầy tự tin: ``Gia đình tôi sẽ hỗ trợ toàn bộ phần dựng rạp và điện nước. Các bạn lớp 1-C sẽ chia ra hỗ trợ cả hai khu vực.''

Tôi nhìn vào danh sách công việc đang dày đặc dần lên, khẽ nhíu mày: ``Nanase-san, cậu vừa phải quán xuyến việc kinh doanh của gia đình, vừa điều phối các bạn trong lớp cho cả hai lễ hội\dots\ Cậu không thấy mình đang ôm quá nhiều việc sao?''

Nanase khựng lại một chút, rồi nở một nụ cười đầy kiêu hãnh: ``Đừng có coi thường sức mạnh và sự bền bỉ của người thừa kế tập đoàn Furukawa nhé, Kaigou-san! Với lại, tôi không làm một mình. Ayame-san, Wata-san và mọi người đều sẽ hỗ trợ mà.''

Hội trưởng Hội học sinh vỗ vai cô: ``Đúng thế! Bọn mình sẽ là hậu phương vững chắc cho cậu!''

Việc phân vai nhanh chóng được chốt lại. Tiểu thư nhà Furukawa điều phối chung các gian hàng, trong khi Sakura-san lo phần lễ nghi. Ayame và Wata làm cầu nối với Hội đồng thị trấn, còn Touka và Suzune quản lý chất lượng thực phẩm (dù cô nàng Thần đồng phải canh chừng cô bạn thân rất gắt gao). Honoka, Azuki và Yuki phụ trách văn nghệ trên sân khấu chính, còn Inari thì chạy đôn chạy đáo hỗ trợ hậu cần cho tất cả mọi người.

Về phía nhà Hikawa, Hội trưởng cũng không quên giao trọng trách: ``Kuon-kun sẽ là Trưởng ban Hậu cần và An ninh hỗ trợ cho cả lớp, vì tính cậu vốn cẩn thận và đáng tin. Kaigou-san và Suzuki-chan sẽ giúp Sakura-san chuẩn bị khu vực đền và hỗ trợ sạp hàng ăn uống khi cần thiết. Còn Naomi-chan\dots\ em chính là 'Nàng thơ' của lễ hội, nhiệm vụ của em là mặc yukata thật đẹp và mỉm cười với tất cả mọi người nhé!''

Nao nghe thấy tên mình liền nhảy cẫng lên vui sướng. Nhưng sự chú ý của mọi người bỗng chốc đổ dồn về phía Shion khi cô nàng tiến lại gần Kuon-kun, đôi mắt trốn dưới mái tóc dày tỏa ra một luồng khí tức kỳ quái.

``Kuon à\dots\ để 'Con đường Linh hồn' có thể thực sự chạm đến ranh giới giữa sự sống và cái chết, ta cần một giọng nói từ vực thẳm\dots\ một thanh âm có thể khiến linh hồn của kẻ gan dạ nhất cũng phải run rẩy,'' Shion chỉ tay vào chiếc đồng hồ đeo tay của Kuon-kun.

Bỗng nhiên, một giọng nói trầm khàn, mang theo chút âm hưởng kim loại và đầy vẻ mỉa mai phát ra từ chiếc đồng hồ: ``Ồ? Cái cô bé mặc đồ đen này muốn mượn giọng của ta sao? Chậc, lũ phàm trần các ngươi đúng là hay bày vẽ. Nhưng mà\dots\ nghe cũng có vẻ thú vị đấy. Được nhìn thấy chúng khóc thét vì sợ hãi thì cũng không tệ chút nào.'' Bình thường ông ta vốn rất 'nhây' và hay đưa ra những lời bình phẩm cay độc, nhưng có vẻ ý tưởng dọa người của Shion đã khơi dậy sự hứng thú của ông ta.

Kuon-kun nhìn tôi, rồi nhìn Shion, khẽ thở dài: ``Nếu Game đã đồng ý và việc này giúp ích cho quy mô của lễ hội, thì tôi không phản đối. Cậu có thể mượn chiếc đồng hồ trong thời gian chuẩn bị.''

Shion đón lấy chiếc đồng hồ với vẻ mặt sùng bái như đang nhận một thánh vật cổ xưa. Cô nàng vung tay áo, cười đầy vẻ ma mị: ``Fufufu\dots\ Tuyệt vời! Kuon đã ban cho ta 'Tử khí Vĩnh hằng'! Với thánh vật được gia trì sức mạnh từ không gian thứ tư này, lực sát thương tâm linh của nhà ma chúng ta sẽ đạt tới mức tuyệt diệt tối thượng! Kẻ nào bước vào đó, hãy chuẩn bị sẵn tinh thần để linh hồn bị nghiền nát đi!''

Tôi lắc đầu cười khổ trước sự chuunibyou của Shion, nhưng cũng thầm mừng vì sự phối hợp này sẽ khiến lễ hội trở nên hoành tráng hơn. Lúc này, Nao kéo kéo áo Suzu, khẽ hỏi: ``Suzuki-oneechan, Lễ hội thị trấn thực sự sẽ có những gì ạ? Nó có giống như trong những câu chuyện chị kể không?''

Suzu mỉm cười dịu dàng, quỳ xuống ngang tầm mắt của Nao: ``Sẽ có nhiều thứ lắm, Naomi-chan ạ. Có những gian hàng trò chơi này, có những món ăn ngon bốc khói nghi ngút này, và quan trọng nhất là vào buổi tối, mọi người sẽ cùng mặc yukata thật đẹp, cùng nhảy điệu Bon Odori dưới ánh đèn lồng rực rỡ\dots''

``Và sẽ có pháo hoa nữa chứ ạ?'' Nao hỏi, đôi mắt lấp lánh sự mong chờ.

Wata chen vào, giọng đầy ấm áp: ``Đúng thế, Nao-chan. Pháo hoa ở Aogami là đẹp nhất vùng này đấy. Chúng sẽ nổ tung trên bầu trời đêm, tạo thành những bông hoa rực rỡ đủ màu sắc.''

``Và chắc chắn là không có pháo hoa nào màu đen như Shion-san đề xuất ban nãy đâu,'' Kuon nhìn Shion, mỉm cười.

``Hứ, chỉ là mấy trò con nít thôi, chẳng có gì đáng để ồn ào,'' cô nàng bị nhắc tới khoanh tay, hứ một tiếng nhẹ, rõ ràng là vẫn đang dỗi vì bị bác bỏ ý kiến.

Nhìn vẻ mặt hào hứng của Nao và những lời mô tả đầy màu sắc của mọi người, tôi bỗng thấy Kuon khựng lại. Ánh mắt chúng tôi chạm nhau, và tôi biết cậu cũng đang nhớ về cùng một điều.

Đó là lời nói của Hanabi ở vòng lặp đầu tiên, khi mà cô nàng đã từng rủ chúng tôi tới Lễ hội của thị trấn, và đó cũng là điểm khởi đầu cho mọi tai ương sau đó chỉ vì cú đập mạnh từ Suzune vào lưng của Kuon.

Lúc đó, chúng tôi vẫn còn là những người xa lạ, loay hoay trong vòng lặp vô tận. Còn giờ đây, chúng tôi đang đứng ở đây, cùng với những người bạn thực sự, chuẩn bị cho một kỷ niệm mà có lẽ sẽ là đẹp nhất đời mình.

Tôi hít một hơi thật sâu, cảm nhận sự náo nhiệt đang dần lan tỏa khắp căn phòng.

Một chương mới, đầy ắp tiếng cười, những giọt mồ hôi chuẩn bị và cả những nỗi niềm chia ly ẩn giấu, đã thực sự bắt đầu.

\pov{Hikawa Kuon}

Tiếng bước chân của các bạn học xa dần ngoài cổng, trả lại cho ngôi nhà Hikawa sự tĩnh lặng vốn có của nó. Nhưng cái tĩnh lặng này không còn giống như trước đây. Nó mang theo hơi thở của mùa hè, của những dư vị từ cuộc thảo luận sôi nổi vừa rồi, và cả một sự bồn chồn thầm kín đang lớn dần trong lồng ngực tôi.

Tôi đứng tựa cửa, nhìn theo bóng dáng Shion đang vừa đi vừa lẩm bẩm với chiếc đồng hồ của tôi trên tay. Cô nàng trông có vẻ cực kỳ đắc ý, như thể vừa nhận được chìa khóa mở ra kho tàng hắc ám của cả vũ trụ.

Bỗng nhiên, chiếc đồng hồ trên cổ tay cô nàng Phù thủy tự xưng - lúc này đang ở khá xa tôi - bỗng phát ra một luồng sáng nhẹ. Giọng nói của Game vang lên, không quá lớn nhưng đủ để lọt vào tai tôi qua làn gió chiều, trầm khàn và đầy vẻ cảnh báo:

``Này, Kuon-user, Kaigou-user và Suzuki-user. Đừng có mải mê với mấy cái gian hàng mà quên mất 'Kịch bản gốc'. Trăng sẽ đạt đỉnh vào đêm cuối cùng của chuỗi sự kiện này. Nếu 'Nhân vật chính' (Naomi) không được hiệu chuẩn hoàn toàn về mặt linh hồn trước khi màn đêm buông xuống, thì 'Sân khấu' (thế giới này) sẽ đóng lại mãi mãi mà không có lối thoát cho bất cứ ai đâu. Nhớ lấy, sự cân bằng đang mỏng manh lắm rồi.''

Shion giật mình, rồi cô nàng cười khoái chí: ``Fufufu\dots\ Giọng nói của vực thẳm đang thúc giục ta chuẩn bị cho đại lễ! Ta hiểu rồi, linh hồn của ta đã sẵn sàng cho sự cộng hưởng này!''

Tôi khẽ thở dài. Game vẫn là Game, lão ta không bao giờ nói thẳng, nhưng lời nhắc nhở đó thì rõ ràng như ban ngày. Nghi thức Đêm trăng rằm không chỉ là để cứu Naomi, mà còn là thời khắc quyết định sự tồn vong của chúng tôi tại nơi này. Lão ta không nhắc đến ``cánh cổng về thế giới thực'', nhưng cụm từ ``đóng lại mãi mãi'' đủ khiến tôi phải rùng mình.

``Kuon-san?''

Tiếng gọi nhẹ nhàng cắt đứt dòng suy nghĩ của tôi. Yuko-sensei đang đứng trong phòng khách, tay cầm mấy chiếc tách trà trống không. Cô là người duy nhất còn nán lại, viện cớ là giúp chúng tôi thu dọn một chút trước khi ra về.

Tôi bước vào trong, đỡ lấy khay trà từ tay cô. ``Để em làm cho, Sensei. Cô cứ nghỉ ngơi đi ạ.''

Yuko-sensei nhìn tôi, ánh mắt cô sâu thẳm và dịu dàng, như thể cô có thể nhìn thấu qua lớp mặt nạ điềm tĩnh mà tôi luôn cố gắng duy trì. Cô không ngồi xuống, mà bước lại gần cửa sổ, nhìn ra rặng núi Aogami đang dần chìm trong bóng tối.

``Tôi có cảm giác\dots\ thời gian ở Aogami đang bắt đầu trôi nhanh hơn bình thường,'' cô nói, giọng trầm xuống. ``Nhìn các em hào hứng chuẩn bị cho lễ hội, tôi thấy rất mừng. Nhưng trong ánh mắt của em, của Kaigou-san và Suzuki-san\dots\ có một điều gì đó rất khác. Nó giống như sự trân trọng cuối cùng của một người khách sắp phải đi xa.''

Tôi khựng lại. Sự nhạy cảm của Yuko-sensei luôn khiến tôi phải kinh ngạc. Tôi biết mình không thể giấu cô mãi được, nhất là khi lễ hội này có thể là lần cuối cùng chúng tôi gặp nhau.

``Sensei\dots\ đúng là như vậy ạ,'' tôi nói, đặt khay trà xuống bàn, giọng nói trở nên nghiêm túc hơn bao giờ hết. ``Sau khi lễ hội kết thúc, và sau khi Nghi thức Đêm trăng rằm hoàn tất\dots\ chúng em sẽ phải rời khỏi nơi này.''

Căn phòng bỗng chốc trở nên im ắng lạ thường. Chỉ còn tiếng kim đồng hồ treo tường nhích từng nhịp đều đặn. Yuko-sensei không quay lại, nhưng tôi thấy vai cô hơi rung nhẹ.

``Rời đi sao\dots'' cô lặp lại, như thể đang cố gắng tiêu hóa sự thật đó. ``Tôi đã luôn cảm nhận được các em không thuộc về nhịp sống bình thường của thị trấn này. Các em mang theo một sứ mệnh, một sức mạnh mà Aogami chưa từng thấy. Nhưng tôi cứ hy vọng\dots\ hy vọng rằng sự gắn kết của lớp 1-C có thể giữ chân các em lâu hơn một chút.''

``Chúng em cũng rất muốn ở lại, Sensei. Thực sự đấy,'' tôi nói, cảm nhận một sự nghẹn lại nơi cổ họng. ``Nhưng nếu chúng em không đi, Naomi sẽ không bao giờ có được một cuộc sống bình thường, và sự tồn tại của chúng em sẽ trở thành gánh nặng cho không gian này.''

Đúng lúc đó, một bóng người âm thầm xuất hiện nơi cửa ra vào hành lang gỗ. Là Shino. Cô nàng đứng đó, bóng tối của mái hiên che khuất nửa khuôn mặt, nhưng đôi mắt xanh đậm của cô thì sáng rực lên một cách kỳ lạ.

``Cậu ấy nói đúng đấy, Yuko-sensei,'' Shino bước vào phòng, bước chân không hề phát ra tiếng động.

Yuko-sensei giật mình quay lại: ``Misora-san? Em chưa về sao?''

Shino khẽ lắc đầu, cô nhìn tôi với ánh mắt đầy sự thấu hiểu và cả một chút kính trọng. Cô quay sang nhìn Yuko-sensei, giọng nói mang theo sự trang nghiêm của một người phục vụ thần linh: ``Họ là những người lữ hành từ nơi rất xa, được định mệnh đưa tới đây để hóa giải những oán hận và sai lệch của Aogami. Sự hiện diện của họ là một phép màu, nhưng phép màu nào cũng có lúc phải kết thúc. Sau khi họ đã hoàn thành Nghi thức, họ cũng sẽ không còn lý do gì để ở lại đây được nữa.''

Yuko-sensei nhắm mắt lại, một giọt nước mắt khẽ lăn dài trên má. Cô hít một hơi thật sâu, rồi mỉm cười - một nụ cười bao dung và đầy nghị lực. ``Tôi hiểu rồi. Vậy ra đây là lý do các em dốc hết sức cho lễ hội này\dots\ Đó là lời chào tạm biệt của lớp 1-C này cho các em.''

Tôi cúi đầu: ``Vâng. Chúng em muốn để lại những ký ức đẹp nhất. Để dù chúng em không còn ở đây, thì lớp 1-C và Sensei vẫn có thể nhớ về chúng em như một phần của mùa hè Aogami.''

Shino bước lại gần, đặt tay lên vai tôi. Một luồng linh lực ấm áp truyền qua, khiến tâm trí tôi bình ổn lại đôi chút. ``Mình và Sakura-san sẽ lo phần Nghi thức ở đền. Bọn mình sẽ đảm bảo linh khí của thị trấn ổn định nhất để hỗ trợ các cậu. Kuon-san, hãy cứ tập trung cho những ngày hội sắp tới. Đừng để nỗi buồn làm lu mờ đi ánh sáng của pháo hoa.''

Tôi nhìn hai người phụ nữ - một là người cô giáo luôn bảo bọc chúng tôi bằng tình thương, một là người bạn đồng hành đã cùng chúng tôi vượt qua những gian khó. Sự ủng hộ của họ tiếp thêm cho tôi một sức mạnh to lớn.

``Cảm ơn hai người. Thực sự\dots\ cảm ơn rất nhiều,'' tôi nói, lòng tràn đầy sự cảm kích.

Khi Shino và Yuko-sensei cuối cùng cũng rời đi, tôi đứng một mình giữa phòng khách vắng lặng. Ánh trăng non bắt đầu hiện ra phía chân trời, nhợt nhạt và xa xăm.

Tôi khẽ đưa tay lên nhìn cổ tay mình, nơi chiếc đồng hồ vốn hiện hữu. Dù nó đang ở trong tay Shion, nhưng sợi dây liên kết giữa tôi và Game vẫn rung động không ngừng.

Thời gian không còn nhiều. Mỗi giây phút trôi qua giờ đây đều quý giá như vàng ngọc. Tôi quay vào trong, nhìn thấy Kaigou đang dọn nốt mấy chiếc gối ôm trên sàn, Suzuki đang kiểm tra lại bếp, và Naomi thì đã ngủ thiếp đi trên ghế sofa với một nụ cười mãn nguyện.

Đây chính là những gì tôi muốn bảo vệ. Dù cái giá phải trả là sự chia ly, dù tương lai phía trước vẫn còn là một ẩn số đại diện cho thế giới thực của chúng tôi, tôi cũng sẽ không bao giờ hối hận.

``Chuẩn bị thôi,'' tôi tự nhủ, ánh mắt rực cháy một sự quyết tâm mãnh liệt. ``Lễ hội cuối cùng\dots\ sẽ là một lễ hội rực rỡ nhất mà Aogami từng chứng kiến.''

\ct

Những ngày sau đó, thị trấn Aogami như biến thành một công trường khổng lồ đầy náo nhiệt. Tiếng cưa máy, tiếng búa gõ và tiếng cười nói vang vọng khắp mọi ngóc ngách, từ cổng trường trung học cho đến tận những bậc thang dẫn lên đền Aogami.

Lớp 1-C chúng tôi thực sự đã tạo nên một cơn sốt. Bất cứ nơi nào chúng tôi đi qua, người dân thị trấn cũng đều nhìn với ánh mắt tò mò và đầy thán phục trước sự nhiệt huyết của đám trẻ.

Tôi, với vai trò Trưởng ban Hậu cần, dường như không có một phút nghỉ ngơi. Sáng sớm, tôi phải có mặt ở trường để kiểm tra tiến độ trang trí cho quán café và nhà ma. Hotaru và Mari đã biến những bức tường hành lang cũ kỹ thành những tác phẩm nghệ thuật mang hơi hướng cổ điển châu Âu, với những họa tiết ren và hoa văn tinh xảo.

``Kuon-kun! Cậu xem tấm biển hiệu này đặt ở đây đã ổn chưa?'' Hotaru gọi vẫy, tay cầm một tấm gỗ lớn vẽ hình một cô hầu gái đang mỉm cười chào đón.

Tôi bước lại gần, điều chỉnh vị trí tấm biển thêm một chút cho cân đối. ``Ổn rồi đấy. Nó sẽ rất nổi bật khi khách bước từ cầu thang lên.''

Đến chiều, tôi lại phải chạy sang khu vực đền để hỗ trợ nhóm của Sakura. Công việc ở đây đòi hỏi sự tỉ mỉ và tôn trọng truyền thống hơn nhiều. Chúng tôi phải giúp các tu sĩ dọn dẹp sân đền, treo hàng trăm chiếc đèn lồng giấy dọc theo lối đi dẫn lên chính điện.

Suzuki và Kaigou cũng bận rộn không kém. Tôi thường thấy hai chị em họ đang loay hoay với những chiếc ema (thẻ điều ước), giúp cô nàng Vu nữ sắp xếp chúng sao cho thật đẹp mắt. Naomi thì như một chú chim nhỏ, bay nhảy khắp nơi, thỉnh thoảng lại giúp mọi người bê những món đồ nhẹ hoặc đơn giản là cổ vũ bằng những nụ cười rạng rỡ.

Nhưng điều khiến tôi chú ý nhất chính là Shion. Cô nàng phù thủy tự xưng dường như đã hoàn toàn chìm đắm vào việc xây dựng ``Con đường Linh hồn''. Nhờ có sự giúp đỡ của Game - thông qua chiếc đồng hồ - cô ấy đã tạo ra những hiệu ứng âm thanh và không gian mà tôi cũng phải thấy rùng mình khi đi ngang qua phòng chuẩn bị.

``Fufufu\dots\ Kuon! Ngươi có nghe thấy tiếng than khóc của những linh hồn bị giam cầm không?'' Shion cười đầy vẻ ma mị khi tôi mang cho cô ấy một ít nước uống. ``Nữ hoàng Bóng tối ta đây sẽ khiến những kẻ phàm trần dám bước vào đây phải khiếp sợ và run rẩy quỳ gối trước sức mạnh vĩ đại của ta!''

Game, từ chiếc đồng hồ, buông một lời mỉa mai: ``Tôi đã phải hạ thấp tần số sóng âm xuống mức tối thiểu rồi đấy. Nếu không, mọi người ở đây sẽ ngất xỉu ngay từ bước chân đầu tiên mất. Con bé này đúng là một kẻ nghiện công việc đầy hắc ám.''


Azuki và Koishin thì đang vây quanh một đống dây cáp và loa đài ở góc phòng chuẩn bị. Cả hai trông vô cùng tập trung, thỉnh thoảng lại trao đổi với nhau bằng những thuật ngữ kỹ thuật mà tôi khó lòng hiểu hết.

``Koishin-san, kiểm tra lại bộ cảm biến ở góc cua thứ hai đi,'' Azuki nói, tay đang thoăn thoắt điều chỉnh bảng điều khiển âm thanh. ``Tiếng hét phải vang lên đúng lúc khách dẫm chân vào tấm ván lỏng, nếu không sẽ mất đi hiệu ứng giật mình đấy.''

Koishin gật đầu, tay cầm chiếc đèn pin soi vào các bảng mạch: ``Đã rõ. Tôi cũng đã cài đặt lại hệ thống đèn strobe để nó chỉ nháy với cường độ 10\% thôi. Bóng tối càng dài thì nỗi sợ càng lớn, đúng không?''

Tôi bước lại gần, không khỏi cảm thấy thán phục trước sự chuyên nghiệp của hai người họ: ``Hai cậu thực sự định biến căn phòng này thành một cơn ác mộng công nghệ cao đấy à?''

Cô nàng Ca sĩ nở một nụ cười đầy tự tin: ``Tất cả nằm ở sự chính xác, Kuon-san. Một giây lệch nhịp cũng có thể phá hỏng cả một bầu không khí sợ hãi mà Shion-san đã dày công xây dựng đấy.''

Tôi chỉ biết cười trừ. Sự nhiệt huyết của mọi người dường như đã lan tỏa và xua tan đi phần nào nỗi buồn chia ly đang âm ỉ trong lòng tôi.

Ở một góc khác của sân trường, Akane đang loay hoay với chiếc máy ảnh chuyên nghiệp của Uzuki. Cô nàng trông có vẻ khá lúng túng khi phải đối mặt với vô số nút bấm và vòng xoay trên ống kính.

``Uzuki-san! Sao cái hình này cứ mờ tịt đi thế? Mình đã cố nhắm vào bông hoa kia rồi mà!'' nàng bờm Khuyển kêu lên, đôi tai trên đầu thỉnh thoảng lại vểnh lên đầy vẻ sốt ruột.

Cô nàng Nhiếp ảnh khẽ chỉnh lại gọng kính, kiên nhẫn hướng dẫn: ``Cậu phải điều chỉnh lại tốc độ màn trập và giữ tay thật chắc, Akane-san. Đừng có hưng phấn quá mà làm rung máy. Đây không giống cậu chụp trên điện thoại đâu, nó đòi hỏi sự tỉ mỉ hơn nhiều.''

Tôi bước lại gần, nhìn vào màn hình xem lại của chiếc máy ảnh: ``Ghi lại những khoảnh khắc của lễ hội là một trọng trách lớn đấy. Cố lên nhé, Akane-san.''

``Cứ tin ở mình!'' Akane nhe răng cười, vẻ đầy tự tin dù mồ hôi đang lấm tấm trên trán. ``Mình sẽ trở thành nhiếp ảnh gia xuất sắc nhất mà Aogami từng có! Cậu cứ chờ mà xem, Kuon-san!''

Trong khi đó, tại khu vực nhà bếp tạm thời của quán café, Saya đang phải đối mặt với một thử thách hoàn toàn khác: kiềm chế cơn thèm ăn vô tận của Suzune.

``Suzune-san\dots\ cậu đã nếm thử đến cái bánh ngọt thứ tư rồi đấy\dots\ Nếu cứ thế này thì chúng ta sẽ không còn đủ mẫu cho buổi thử thực đơn đâu\~{}'' cô nàng Thiên văn lên tiếng với nhịp điệu chậm rãi và êm ái đặc trưng, cố gắng nhẹ nhàng lấy đĩa bánh ra khỏi tay cô bạn.

Cô bạn ``Linh vật'' phồng má, ánh mắt vẫn không rời khỏi những chiếc bánh kem đang tỏa hương thơm phức: ``Nhưng tớ cần phải chắc chắn chứ! Kiểm soát chất lượng là khâu quan trọng nhất mà! Nếu tớ không nếm, làm sao biết khách hàng có hài lòng hay không?''

Saya khẽ thở dài, mỉm cười bất lực: ``Vậy thì để tớ giúp cậu bày đĩa nhé\dots\ như vậy sẽ an toàn hơn cho số bánh còn lại của lớp mình đấy\~{}''

Nhìn cảnh tượng dở khóc dở cười đó, tôi thầm nghĩ rằng dù bận rộn và có chút lộn xộn, nhưng chính sự hồn nhiên và nhiệt huyết của mỗi người mới là điều làm nên linh hồn của lễ hội này.

Và dĩ nhiên, không thể thiếu buổi thử trang phục đầy ``sóng gió'' tại phòng chuẩn bị của câu lạc bộ Thời trang.

Tôi đứng trước gương, loay hoay với chiếc cà vạt nơ của bộ đồ quản gia. Bộ vest đen tuyền được cắt may vừa vặn đến từng milimet, đi kèm với chiếc áo sơ mi trắng muốt và đôi găng tay trắng tinh khôi. Tôi cảm thấy mình cứng nhắc như một pho tượng đá trong bảo tàng. Chưa bao giờ tôi lại phải mặc bộ đồ bảnh bao như thế này; đến cả hồi nhận việc cho chức Tổng quản ở \hll\ tôi cũng chỉ mặc một bộ đồ lịch sự thôi.

``Ồ! Kuon-kun nhìn bảnh quá nhỉ!'' Ayame vỗ tay tán thưởng, đôi mắt ranh mãnh quét qua tôi từ đầu đến chân. ``Đúng là 'người đẹp vì lụa', nhìn cậu chẳng khác gì một quản gia trưởng từ một dinh thự cổ ở Anh quốc cả.''

``Tôi cảm thấy như mình sắp ngạt thở vì cái cổ áo này rồi đây,'' tôi than vãn, cố gắng nới lỏng nó ra một chút nhưng vô ích. Hotaru đứng bên cạnh, tay cầm kim chỉ, gõ nhẹ vào vai tôi: ``Đừng có làm hỏng form áo của tôi chứ! Cậu phải đứng thẳng lưng lên, đó mới là phong thái quản gia!''

Nhưng sự chú ý của mọi người ngay lập tức bị hút về phía góc phòng, nơi Kaigou vừa bước ra từ sau tấm rèm thay đồ.

Cô nàng đứng đó trong bộ trang phục hầu gái với tông màu đen trắng kinh điển. Những đường ren tinh xảo dọc theo tạp dề và chiếc băng đô nhỏ trên mái tóc đuôi ngựa làm tôn lên vẻ đẹp sắc sảo nhưng cũng đầy lạnh lùng của cô. Tuy nhiên, trái ngược với vẻ ngoài xinh đẹp đó là một gương mặt đang cau có đến cực độ.

Kaigou gầm gừ, đôi tay không ngừng kéo gấu váy xuống. ``Tại sao tôi lại phải mặc thứ này để đi phục vụ đám đàn ông cặn bã ngoài kia chứ?''

Hanabi, người nãy giờ vẫn đang ghi chép trong sổ tay, bỗng reo lên đầy phấn khích: ``Tuyệt vời! Kaigou-san, chính là nó! Cái ánh nhìn khinh bỉ đó, cái thái độ không chút nhân nhượng đó! Cậu đang vào vai 'Hầu gái Nữ vương' cực kỳ chuẩn bài luôn!''

Kaigou quay sang lườm cô bạn một cái cháy mặt: ``Tôi không có `vào vai' gì hết! Tôi đang thực sự cảm thấy ghê tởm đấy!''

Cô nàng hậm hực bước lại gần tôi, hạ thấp giọng xuống mức chỉ đủ cho hai chúng tôi nghe thấy, nhưng ánh mắt thì vẫn giữ nguyên vẻ băng giá: ``\vie{\hl{Nói thật nhé Kuon-kun, nếu bắt tôi phải tỏ ra ngọt ngào với đám rác rưởi ngoài kia, tôi thà chỉ phục vụ mình cậu còn thấy dễ chịu hơn vạn lần.}}''

Tôi thở dài, cảm thấy hơi nóng bắt đầu bốc lên mặt. Câu nói thẳng thừng đó của cô ấy khiến tôi chẳng biết nên phản ứng thế nào cho phải, chỉ biết hắng giọng rồi quay mặt đi chỗ khác để tránh cái nhìn thiêu đốt của cô nàng. Dường như với Kaigou, ranh giới giữa 'thật' và 'diễn' đã hoàn toàn biến mất trong bộ đồ hầu gái này rồi.

Trong khi đó, Hanabi xoay một vòng, chiếc váy của cô tung bay nhẹ nhàng. Khác với bộ của Kaigou, bộ trang phục của Hanabi có những đường cắt xẻ tinh tế ở phần vai và eo, làm nổi bật vóc dáng năng động và đầy sức sống của cô. ``Cậu thấy sao, Kuon-san? Mình đã tự tay chỉnh sửa lại bộ này để nó trông 'lôi cuốn' đúng như vai diễn của mình đấy! Những đường xếp ly này sẽ tạo hiệu ứng rất tốt khi mình di chuyển giữa các bàn.''

Miku thì lại chọn cho mình một bộ đồ mang phong cách cổ điển và kín đáo hơn, với phần váy dài quá gối và chiếc tạp dề được thêu ren thủ công tỉ mỉ. Cô vuốt lại nếp nhăn trên áo, toát lên một vẻ thanh lịch và quý phái khó cưỡng, đúng chất của một 'hầu gái trưởng'. ``Mình thì thấy phong cách này thoải mái và phù hợp với mình hơn. Nó mang lại cảm giác của một người quản lý thực thụ trong quán café. Cậu thấy nó có hợp với mình không, Kuon-san?''

Tôi chỉ biết gật đầu lia lịa, thầm thán phục sự đầu tư và tâm huyết mà các cô gái dành cho lễ hội lần này. ``Mọi người đều trông rất tuyệt. Chắc chắn gian hàng của lớp chúng ta sẽ là nơi nổi bật nhất trường cho xem.''

Trong khi không khí ở phòng chuẩn bị đang vô cùng căng thẳng (và có phần ngại ngùng), thì từ phía phòng nhạc, những giai điệu sôi động của mùa hè đã bắt đầu vang vọng khắp các dãy hành lang.

Nhóm nhạc của Azuki, Honoka và Yuki đang dốc hết sức cho buổi tập luyện cuối cùng. Tiếng ghi-ta điện của Honoka gầm rú đầy máu lửa, hòa cùng nhịp trống dồn dập của Azuki và những âm thanh điện tử biến ảo từ chiếc keyboard của Yuki.

``\textbf{Một, hai, ba! Lại đoạn điệp khúc nào!}'' tiếng hét đầy năng lượng của Honoka át cả tiếng nhạc.

Tôi đứng từ xa nhìn vào qua khe cửa lớp nhạc. Cả ba cô gái đều đang đẫm mồ hôi, nhưng ánh mắt họ thì rực cháy một niềm đam mê mãnh liệt. Họ đang chơi một bản rock có tiết tấu cực nhanh, một bản nhạc mà họ gọi là 'Festival Fever' - linh hồn của sân khấu chính năm nay.

Sự cuồng nhiệt đó, sự nỗ lực đó\dots\ tất cả đang hòa quyện vào nhau, tạo nên một bầu không khí tiền lễ hội đầy phấn khích. Mọi thứ đang dần vào guồng, mỗi mảnh ghép đang tìm thấy vị trí của mình để tạo nên một bức tranh hoàn hảo nhất.

Trong khi công việc ở trường đã dần đi vào ổn định, thì tại khu vực quảng trường và sân đền, không khí chuẩn bị cho Lễ hội thị trấn cũng đang ở giai đoạn nước rút. Dù ban đầu lớp chúng tôi chỉ định tập trung cho lễ hội văn hóa, nhưng cuối cùng một số người cũng tự nguyện tham gia hỗ trợ cho sự kiện lớn nhất của Aogami.

Shiina và Yae - những người được giao trọng trách hậu cần nặng nhọc nhất - đang cùng các công nhân nhà Furukawa khuân vác những tấm gỗ lớn để dựng khung cho các sạp hàng.

``Phù! Mấy tấm ván này nặng hơn vẻ ngoài của chúng đấy,'' Shiina thở hắt ra, quẹt mồ hôi trên trán trong khi vẫn đang giữ chặt một đầu thanh xà gỗ. ``Yae, cậu giữ chắc đầu kia nhé, tôi chuẩn bị hạ xuống đây.''

Yae mỉm cười, đôi tay săn chắc giữ chặt lấy thanh gỗ mà không hề lộ vẻ mệt mỏi: ``Yên tâm đi! Bọn mình còn mười sáu cái khung sạp nữa phải dựng trước khi mặt trời lặn đấy. Cố lên nào!''

Tôi bước tới định giúp một tay nhưng Shiina đã nhanh chóng xua tay: ``Thôi nào Kuon, cậu còn phải đi kiểm tra danh mục hàng hóa cho Nanase nữa mà. Chút việc chân tay này cứ để bọn này lo. Dù sao thì vận động một chút cũng tốt hơn là ngồi lì trong lớp.''

Yae gật đầu đồng tình, giọng đầy hào hứng: ``Đúng thế! Coi như đây là buổi tập thể lực đặc biệt đi. Đến lúc lễ hội bắt đầu, bọn tớ sẽ khỏe đến mức có thể khiêng cả cái kiệu Mikoshi chạy quanh thị trấn cho xem!''

Những giọt mồ hôi rơi xuống trên sàn gỗ trường học, trên nền đá đền cổ kính, tất cả đều là minh chứng cho một tình bạn, một sự gắn kết mà tôi biết mình sẽ mang theo mãi mãi.

Chúng tôi đang sống trong những ngày rực rỡ nhất của mùa hè, nơi mà ranh giới giữa thực và ảo, giữa quá khứ và tương lai dường như mờ nhạt đi trước sức nóng của lòng nhiệt huyết.

Aogami đang thay da đổi thịt từng ngày, sẵn sàng cho thời khắc cho một lễ hội cuồng nhiệt bùng nổ nhất. Và chúng tôi, những người lữ hành, cũng đang chuẩn bị cho màn trình diễn cuối cùng và hoành tráng nhất của mình.

\begin{center}
  \uline{Nhiều ngày sau.}
\end{center}

Bình minh của ngày hội văn hóa rốt cuộc cũng đã ló rạng.

Khi chúng tôi bước qua cánh cổng trường Aogami, một bầu không khí rực rỡ và náo nhiệt đến nghẹt thở lập tức ùa tới, xua tan hoàn toàn cái se lạnh còn sót lại của sương sớm. Ngôi trường vốn dĩ yên tĩnh nay đã lột xác hoàn toàn; những dải băng rôn đa sắc treo lủng lẳng từ các cửa sổ lớp học, những tấm bảng hiệu vẽ tay sặc sỡ dựng dọc khắp các lối đi, và mùi thơm của bánh crepe, yakisoba bắt đầu lan tỏa trong không gian.

Các lớp học khác cũng không hề kém cạnh trong việc phô diễn sự sáng tạo. Lớp 2-A đã dựng hẳn một gian hàng ``VR Dungeon'' đầy mời gọi với những mô hình quái vật khổng lồ; trong khi lớp 3-B lại chọn phong cách trầm mặc hơn với một ``Trà quán Truyền thống'' ngay giữa sân trường để thu hút những vị khách lớn tuổi.

Chúng tôi đi giữa dòng người đang tấp nập qua lại. Suzuki nắm chặt lấy tay Naomi, đôi mắt em ấy lấp lánh sự hào hứng không giấu giếm: ``\vie{Naomi-chan, em thấy không? Mọi thứ còn đẹp hơn cả những gì chị đã tưởng tượng nữa!}''

Naomi thì gần như không thốt nên lời, cô bé cứ liên tục xoay tròn, ngón tay nhỏ xíu chỉ trỏ vào những quả bóng bay và những người bạn học đang mặc những bộ linh vật ngộ nghĩnh. Trong khi đó, Kaigou lại đi sát phía sau chúng tôi, ánh mắt cô sắc lạnh như thường lệ, liên tục quan sát xung quanh như một người vệ sĩ thực thụ đang bảo vệ một đoàn xe quan trọng.

``\vie{Đừng có căng thẳng quá như vậy, Kaigou-chan,}'' tôi khẽ nói, nụ cười thoáng hiện trên môi. ``\vie{Hôm nay chúng ta là một phần của lễ hội mà.}''

Cô nàng hừ một tiếng nhẹ, nhưng tôi thấy vai cô đã thả lỏng hơn đôi chút: ``\vie{Tôi chỉ đang đảm bảo không có kẻ nào phá hỏng tâm trạng của Nao thôi. Và cả của Suzu nữa.}''

Bỗng nhiên, hệ thống loa phóng thanh của trường phát ra những tiếng rè rè quen thuộc, rồi vang lên giọng nói đầy năng lượng và sôi nổi của bộ đôi dẫn chương trình quen thuộc - Kaoru và Kana.

``\textbf{Chào buổi sáng, trường Trung học Aogami! Đây là Kaoru và Kana, những người sẽ đồng hành cùng các bạn trong hai ngày hội rực rỡ nhất năm nay!}'' giọng của Kaoru vang lên đầy phấn khích.

``\textbf{Đúng vậy! Bầu không khí lúc này đang nóng hơn bao giờ hết!}'' Kana tiếp lời. ``\textbf{Và bây giờ, để chính thức bắt đầu cho Lễ hội Văn hóa Aogami lần này, xin mời thầy Hiệu trưởng lên có đôi lời phát biểu khai mạc!}''

Tiếng vỗ tay vang dội khắp sân trường. Thầy Hiệu trưởng bước lên bục cao, bài phát biểu của thầy ngắn gọn nhưng tràn đầy sự kỳ vọng: thầy nhắc nhở chúng tôi hãy tận hưởng hết mình, hãy tạo nên những kỷ niệm không bao giờ quên, và quan trọng nhất là hãy giữ gìn tinh thần đoàn kết của học sinh Aogami.

Khi lời tuyên bố ``Lễ hội chính thức bắt đầu!'' vang lên, một chùm bóng bay đủ màu sắc được thả lên bầu trời xanh ngắt, hòa cùng tiếng hò reo vang dội của hàng ngàn học sinh và khách tham quan.

Các cửa tiệm đồng loạt mở cửa. Tiếng mời chào, tiếng cười nói và tiếng nhạc bắt đầu đan xen tạo thành một bản giao hưởng náo nhiệt.

Lớp 1-C chúng tôi cũng chính thức đi vào hoạt động, nhưng để có được khoảnh khắc mở màn suôn sẻ này, cả lớp đã phải trải qua một buổi sáng chuẩn bị đầy hối hả.

Từ sáu giờ rưỡi sáng, tôi đã cùng nhóm hậu cần của Inari vận chuyển những thùng dâu tây tươi mọng và hàng chục lít sữa béo ngậy - món quà hỗ trợ từ tập đoàn của Nanase-san - vào khu vực bếp tạm thời. Touka và Kanade đã có mặt từ trước đó, kiểm tra độ sắc của dao và nhiệt độ của lò nướng với sự tập trung cao độ như những chuyên gia thực thụ.

Chúng tôi sắp xếp lại bàn ghế, phủ lên đó những tấm khăn trải bàn bằng ren trắng tinh khôi mà Hotaru đã tỉ mỉ lựa chọn. Trên mỗi bàn, một lọ hoa nhỏ với những nhành cúc họa mi trắng tạo nên điểm nhấn nhã nhặn. Tiếng lách cách của những chiếc tách sứ được xếp lên khay, tiếng máy pha cà phê bắt đầu khởi động, tất cả hòa quyện thành một nhịp điệu hối hả nhưng đầy hứa hẹn.

Trong khi đó, ở khu vực cuối hành lang, ``Con đường Linh hồn'' của Shion lại mang một bầu không khí hoàn toàn trái ngược. Tôi ghé mắt qua khe cửa và thấy một khoảng không đen đặc, chỉ có vài ánh đèn strobe xanh lạnh lẽo thỉnh thoảng chớp nháy. Shion đang thắp một loại nhang trầm có mùi hương kỳ quái để tạo không khí, trong khi chiếc máy tạo khói đang phun ra những làn sương mù dày đặc.

``Hệ thống âm thanh đã sẵn sàng chưa, Game?'' Shion hỏi vào khoảng không.

``Đừng hỏi những điều hiển nhiên,'' giọng Game vang lên, khô khốc và đầy vẻ đe dọa. ``Tôi đã hiệu chỉnh dải tần số để nó đánh thẳng vào nỗi sợ nguyên thủy của lũ người thường. Chỉ cần chúng bước vào, tôi sẽ đảm bảo linh hồn chúng được 'chăm sóc' kỹ lưỡng.''

Tôi khẽ rùng mình, thầm cầu nguyện cho những vị khách gan dạ nhất. Quay trở lại phòng thay đồ, tôi cũng phải chật vật không kém để cài hàng cúc áo Butler cứng nhắc và đeo chiếc găng tay trắng muốt dưới sự giám sát gắt gao của Hotaru.

Đây không còn là sự chuẩn bị nữa. Trận chiến thực sự - và cũng là kỷ niệm đẹp nhất của chúng tôi - đã chính thức bắt đầu.

``Ta cùng bắt đầu tận hưởng Lễ hội Văn hóa thôi nào.''

\pov{Hikawa Suzuki}

Tiếng chuông lanh lảnh treo ở phía trên cánh cửa lớp học 1-C vừa vang lên những nhịp đầu tiên, báo hiệu quán café ``Ma \& Butler'' chính thức mở cửa đón những vị khách đầu tiên của mùa lễ hội. Tôi đứng ở phía sau quầy thu ngân, hai bàn tay đan chặt vào nhau, cảm nhận được trái tim mình đang đập thình thịch liên hồi trong lồng ngực.

Thực sự, bộ trang phục hầu gái mà Hotaru và câu lạc bộ Thời trang đã dày công thiết kế cho tôi\dots\ là một sự kết hợp giữa niềm kiêu hãnh và cả sự ngượng ngùng đến tột độ. Nó có tông màu hồng phấn phối cùng trắng sữa tinh khôi, với tầng tầng lớp lớp ren bồng bềnh ở gấu váy khiến mỗi bước đi của tôi đều trở nên nhẹ bẫng. Trên đầu tôi là một chiếc băng đô gắn tai mèo đen tuyền cực kỳ tinh xảo, và ở hai cổ tay, những quả chuông nhỏ được thắt bằng ruy băng hồng cứ mỗi khi tôi cử động lại phát ra những tiếng ``leng keng'' vui tai. Hotaru nói rằng tôi vào vai ``em gái nhỏ mèo con'' là chuẩn nhất, nhưng thú thật, việc phải nói câu chào mừng khách hàng bằng chất giọng nũng nịu vẫn là một thử thách cực đại đối với tôi.

Nhưng chỉ cần nhìn sang phía Onii-sama, nỗi lo lắng của tôi dường như tan biến quá nửa.

Anh đứng ở ngay cửa ra vào, khoác trên mình bộ quản gia đen tuyền được cắt may vô cùng bảnh bao và lịch lãm. Dù giọng nói của anh ấy vẫn giữ nguyên tông đơn sắc, không chút biểu cảm như thể một cỗ máy đang vận hành, nhưng cách anh ấy cúi chào một góc chuẩn xác 45 độ và dẫn khách vào bàn lại toát lên một sự chuyên nghiệp đến kinh ngạc. Từng động tác bưng khay, rót trà hay sắp xếp khăn ăn đều cực kỳ nhịp nhàng và tinh tế. Dường như những kinh nghiệm từ hồi làm Tổng quản ở \hll\ (anh ấy nói như thế) cho đến những ca làm thêm tại cửa hàng tiện lợi đã rèn giũa cho anh ấy một phong thái phục vụ không chê vào đâu được.

``Chào mừng quý khách đã ghé thăm. Bàn số bốn đã sẵn sàng, xin mời đi theo tôi,'' Onii-sama nói, tay ra hiệu đầy lịch thiệp. Dù khuôn mặt không hề nở một nụ cười rạng rỡ, nhưng sự niềm nở trong từng cử chỉ nhỏ nhặt lại khiến khách hàng cảm thấy họ thực sự đang được phục vụ bởi một quản gia cao cấp. Chính vẻ mặt có chút lạnh lùng như thái độ nồng ấm đó đã ghi điểm được với nhiều người.

Trong khi đó, bầu không khí phía trong quán là một sự pha trộn kỳ diệu giữa sự náo nhiệt và những tình huống lúng túng đầy dễ thương.

``\textbf{Chào mừng chủ nhân đã trở về nhà ạ! Á--!}''

Tiếng hét nhỏ của Honoka vang lên kèm theo một cú vấp chân đầy ``nghệ thuật'' khi cô nàng đang bưng đĩa bánh kem dâu tây. Rất may, nhờ phản xạ nhanh nhạy của một người hoạt bát, cậu ấy đã giữ được đĩa bánh thăng bằng trên lòng bàn tay ngay trước khi nó kịp chạm đất. Cô nàng diện bộ đồ hầu gái có phần váy ngắn và xòe rộng, trang trí bằng rất nhiều dải duyên dáng giúp cô trông như một chú sóc nhỏ tinh nghịch. Cô nàng bờm Cáo sa mạc lè lưỡi cười hì hì: ``Hehe, là do sàn nhà quá yêu quý Honoka nên mới níu chân lại chút xíu thôi mà! Chủ nhân đừng giận nhé, để Honoka bù đắp bằng một nụ cười thật tươi nha!''

Sự nhí nhảnh và có chút hậu đậu đó lại vô tình trở thành điểm nhấn lôi cuốn, khiến các vị khách nam đều bật cười thích thú và chẳng nỡ lòng nào trách mắng.

Ở khu vực trung tâm, Hanabi và Miku đang tạo nên hai thái cực hoàn toàn khác biệt. Cô nàng năng nổ với bộ trang phục hầu gái được cách tân theo phong cách lôi cuốn mà tôi đã thấy hôm thử đồ, mỗi cử chỉ của cô ấy đều toát lên một sự duyên dáng khó cưỡng. Cô nàng khẽ nghiêng người, nụ cười rạng rỡ và ánh mắt lấp lánh khi đứng bên cạnh bàn khách: ``Chào mừng các vị khách quý. Các bạn có muốn thử món 'Bánh kem Ký ức' không? Nó mang hương vị ngọt ngào như một lời hứa giữa mùa hè đấy!'' Cách cô ấy khéo léo giới thiệu thực đơn khiến các vị khách gần như không thể từ chối.

Trong khi đó, cô nàng bờm Sói đen lại khoác lên mình bộ đồ hầu gái nhã nhặn với phần váy dài quá gối, toát lên vẻ trưởng thành và quý phái của một người quản lý thực thụ. Cô đứng thẳng lưng, cử chỉ rót trà vô cùng điềm tĩnh và chuẩn xác: ``Xin mời dùng trà. Nếu cần thêm bất cứ điều gì, xin đừng ngần ngại gọi tôi.''

Dù vậy, mỗi khi có vị khách nào thốt lên lời khen ngợi vẻ ngoài thanh tú của mình, tôi lại thấy đôi tai Miku đỏ ửng lên một cách rõ rệt. Cô ấy sẽ bối rối cúi đầu, đôi bàn tay hơi run nhẹ khi chỉnh lại chiếc tạp dề: ``Cảm ơn quý khách\dots\ tôi\dots\ tôi chỉ đang làm tốt bổn phận của mình thôi ạ\dots'' Vẻ mặt ngượng nghịu đó trái ngược hoàn toàn với dáng vẻ quý phái lúc nãy, tạo nên một sức hút rất riêng.

Nhưng tâm điểm của sự ``hỗn loạn'' thầm lặng lại nằm ở phía quầy nhận món, nơi Suzune đang đứng trực chiến.

Cậu ấy cũng vào vai em gái giống như tôi, nhưng dường như tâm trí cậu ấy hoàn toàn bị bỏ bùa mê bởi hương thơm của những chiếc bánh ngọt. ``Mmm, cái lớp kem này trông mượt mà quá đi\~{} Chắc là vị dâu tây này ngọt lịm lắm đây\~{}'' cậu ấy lẩm bẩm, ngón tay vừa định đưa về phía đĩa bánh mẫu thì một chiếc khay bạc đã nhẹ nhàng nhưng dứt khoát chặn lại ngay trước mũi.

``Suzune-san, đây là lần thứ năm trong buổi sáng nay rồi đấy,'' Touka xuất hiện từ phía sau bức màn bếp với một cái nhìn sắc lẹm. Touka và Kanade cũng diện trang phục hầu gái nhưng được thiết kế gọn gàng hơn với phần tay áo lửng để tiện cho việc nấu nướng. Bộ đồ của cô bạn Thần đồng có màu xanh coban đậm, tôn lên làn da trắng và vẻ nghiêm túc của cậu ấy.

Touka xoay chiếc khay trên tay, ra hiệu cho cô bạn thân quay trở lại bàn khách: ``Quay lại vị trí đi. Nếu cậu còn dám ăn vụng thêm một miếng nào nữa, mình sẽ bảo Onii-sama của chúng ta cắt phần điểm tâm chiều của cậu đấy.''

Suzune phồng má, phụng phịu lẩm bẩm: ``Tớ chỉ muốn kiểm tra xem kem có bị tan không thôi mà\dots\ tớ đi ngay đây!'' Nhìn cậu ấy lật đật chạy đi với khuôn mặt vẫn còn vương chút tiếc nuối, tôi không khỏi mỉm cười. Kanade lúc này cũng ló đầu ra, tay bưng theo hai ly latte nghệ thuật với hình vẽ chú mèo nhỏ rất sinh động, khẽ gật đầu chào khách hàng bằng nụ cười hiền hậu thường thấy.

Phía bên ngoài hành lang, tiếng xì xào bàn tán và tiếng cười đùa bắt đầu rộ lên ngày một lớn. Đó là nhờ công lao to lớn của ``linh vật'' lớp chúng tôi - Naomi.

Em ấy đang mặc một bộ váy ``Công chúa Hoa'' rực rỡ với gam màu vàng chanh và trắng, kết hợp cùng đôi cánh tiên lấp lánh sau lưng và một chiếc vương miện kết bằng hoa tươi trên đầu. Cô bé giống như một thỏi nam châm nhỏ, đứng ở ngay lối vào hành lang để chào mời khách tham quan. Với vẻ đáng yêu thuần khiết và giọng nói trong trẻo như tiếng chuông ngân, Naomi chỉ cần vẫy tay là đã có hàng tá người phải dừng chân lại.

``Mọi người ơi, mau vào quán của lớp 1-C đi ạ!'' Naomi vẫy chiếc gậy phép nhỏ có gắn ruy băng hồng, xoay một vòng khiến chiếc váy bồng bềnh tung bay. ``Ở đây có bánh ngon như phép màu và những anh chị quản gia siêu ngầu luôn! Có cả nhà ma đáng sợ dành cho những ai dũng cảm nữa đấy ạ!''

Cô bé tiến lại gần một nhóm học sinh trường khác đang phân vân, nắm lấy tay một chị gái và ngước đôi mắt to tròn, lấp lánh lên nhìn: ``Chị ơi, chị vào xem Naomi múa một điệu phép thuật nhé? Chỉ cần một miếng bánh thôi, chị sẽ thấy vui cả ngày luôn ạ!'' Không ai có thể cưỡng lại được sự tấn công bằng nét đáng yêu đến mức ``vô cực'' này. Em cười khúc khích khi thấy nhóm khách nọ vui vẻ đi vào quán, cô bé còn không quên nháy mắt tinh nghịch: ``Chúc mọi người có một hành trình kỳ diệu tại lớp 1-C nhé!''

Saya, Hotaru và Mitsuki luôn túc trực ở gần đó để trông chừng cô bé. Cô bạn Thiên văn với dáng vẻ chậm rãi, thỉnh thoảng lại đưa nước cho Naomi, trong khi đôi bạn Thiết kế thời trang thì nhanh nhẹn hỗ trợ phát tờ rơi quảng cáo và hướng dẫn khách hàng xếp hàng theo thứ tự. Họ không trực tiếp phục vụ trong quán nhưng những công việc vụn vặt đó lại là chất xúc tác quan trọng giúp mọi người vận hành trơn tru.

Cứ như vậy, bánh xe của quán café lớp 1-C đã bắt đầu quay những vòng đầu tiên trong sự phối hợp nhịp nhàng đến không ngờ. Mùi hương của hạt cà phê rang xay, vị ngọt ngào của kem tươi hòa quyện cùng tiếng cười nói rộn rã đã biến không gian lớp học cũ kỹ trở thành một thiên đường nhỏ bé.

Tôi hít một hơi thật sâu lần nữa, cảm thấy dòng máu nhiệt huyết đang chảy rần rần trong huyết quản. Không thể để mọi người thất vọng được!

Tôi cầm lấy chiếc khay nhỏ, bước về phía một bàn khách vừa mới ngồi xuống, môi vẽ nên một nụ cười rạng rỡ và nũng nịu nhất có thể: ``Chào mừng chủ nhân đã trở về nhà ạ! Mèo nhỏ Suzuki xin được phép phục vụ quý khách, chủ nhân muốn dùng món gì để nạp thêm năng lượng cho ngày hôm nay ạ?''

Nhưng nếu nói về người sở hữu kỹ năng thay đổi thái độ nhanh đến chóng mặt nhất, thì chắc chắn phải kể đến Onee-sama của tôi.

Đứng từ quầy thu ngân, tôi không khỏi há hốc mồm khi quan sát cách chị tôi phục vụ các nhóm khách khác nhau. Nó giống như thể cô ấy đang sở hữu hai linh hồn hoàn toàn tách biệt trong cùng một cơ thể vậy.

Khi một nhóm nữ sinh lớp dưới rụt rè bước vào, Onee-sama lập tức nở một nụ cười dịu dàng và ấm áp đến mức tôi cũng phải ngỡ ngàng. Chị khẽ cúi người, giọng nói trong trẻo như tiếng suối: ``Chào mừng các tiểu thư đã về nhà. Hôm nay là một ngày tuyệt vời để thưởng thức một tách trà hoa cúc và bánh ngọt, đúng không nào? Hãy để tôi dẫn các bạn đến bàn có góc nhìn đẹp nhất nhé.'' Nhìn cách chị ấy ân cần kéo ghế cho họ, tôi thầm nghĩ: ``Đây đúng là người chị mẫu mực mà mình luôn ngưỡng mộ rồi!''

Thế nhưng, ngay khi chị ấy quay lưng lại để đón một nhóm nam sinh đang tò mò nhìn ngó, bầu không khí xung quanh chị lập tức giảm xuống dưới độ không tuyệt đối. Đôi mắt sắc lẹm như dao cạo, nụ cười biến mất không dấu vết, thay vào đó là một vẻ mặt lạnh lùng và đầy sự khinh miệt.

``Bàn số bảy. Ngồi xuống đi và đừng có gây ồn ào,'' chị ấy buông một câu cụt ngủn, tay chỉ về phía góc quán mà không thèm nhìn lấy họ một cái.

Một cậu bạn run rẩy hỏi: ``Ơ\dots\ cho mình xem thực đơn được không ạ?''

Kaigou ném tập thực đơn lên bàn với một tiếng *CẠCH* khô khốc: ``Nó ở ngay trước mắt cậu đấy. Có mắt thì tự mà xem, đừng có làm lãng phí thời gian của tôi.''

Tôi đứng phía sau, đổ mồ hôi hột, vội vàng chạy lại thì thầm vào tai chị ấy: ``\vie{Onee-sama, ít nhất thì hãy giả vờ lịch sự đi một chút thôi mà\dots\ Chị dọa thế người ta sợ đấy\dots}''

Chị ấy chỉ khẽ nhướng mày, giọng nói vẫn mang theo sự kiêu ngạo của một vị nữ vương: ``\vie{Chị đang lịch sự theo cách của ta rồi, Suzuki. Đám này không xứng đáng để nhận được nụ cười của chị đâu. Chưa bị ăn đấm là may rồi đấy.}''

Kỳ lạ thay, những vị khách nam kia không hề tỏ ra giận dữ, ngược lại họ còn có vẻ phấn khích hơn hẳn, liên tục xì xào: ``Đúng là phong cách Nữ vương mà Touka-san đã nói! Tuyệt quá!'' Có vẻ như ý tưởng của cô bạn Thần đồng đã thực sự đánh trúng vào một thị hiếu kỳ lạ nào đó của người dân Aogami rồi.

Đến khi gặp một cặp đôi, sự ``hai mặt'' của Onee-sama lại càng được đẩy lên đỉnh điểm. Chị nhẹ nhàng đặt tách trà trước mặt bạn nữ: ``Mời tiểu thư dùng trà. Chúc bạn có một buổi hẹn hò thật lãng mạn nhé.'' Nhưng ngay lập tức, chị quay sang nhìn bạn nam với ánh mắt như muốn xé xác: ``Còn cậu, uống nhanh rồi tính tiền đi. Đừng có làm tốn chỗ của những người khách tiềm năng khác.''

Ngay cả với những vị khách lớn tuổi, Onee-sama cũng giữ một thái độ rất mực thước. Chị không niềm nở thái quá như với phái nữ, nhưng cũng không thô lỗ như với giới trẻ. Đó là một sự tôn trọng đúng mực của một người quản gia dành cho những bậc tiền bối, khiến họ cảm thấy rất hài lòng và khen ngợi phong thái của chị không ngớt.

Sự biến hóa khôn lường giữa vai ``Thân thiện'' và vai ``Nữ vương'' đã khiến chị trở thành một trong những nhân vật được săn đón nhất quán café. Tôi chỉ biết thở dài, thầm cảm ơn trời đất vì Onii-sama và Touka-san đã tìm ra cách để biến tính cách khó gần của chị thành một loại ``dịch vụ'' độc nhất vô nhị như thế này.

Mọi thứ trong quán đang vận hành cực kỳ trơn tru, cho đến khi tiếng chuông cửa lại vang lên một lần nữa, báo hiệu một đợt khách mới đang kéo đến. Lễ hội văn hóa thực sự là một cuộc đua sức bền về thể lực và tinh thần, nhưng nhìn thấy nụ cười của mọi người, tôi thấy mọi mệt mỏi dường như tan biến hết.

\ct

Tiếng chuông cửa ngân lên lần thứ bao nhiêu trong suốt buổi sáng mà chẳng ai đếm nổi, nhưng lần này không phải là một vị khách xa lạ. Là giáo viên chủ nhiệm của chúng tôi.

Yuko-sensei bước vào với nụ cười hiền hậu, cô dừng lại trước quầy của tôi và khẽ gật đầu tán thưởng: ``Tất cả các em đều làm rất tốt. Nhìn Suzuki-san thật giống một nàng mèo nhỏ xinh xắn, rất hợp với không khí này. Còn Kaigou-san\dots\ không ngờ em lại có thể hóa thân xuất sắc vào vai diễn Nữ vương như vậy đấy. Có vẻ như các em thực sự đã tìm thấy vị trí của mình ở đây rồi.''

Tôi ngượng nghịu cúi đầu chào cô: ``Cảm ơn Sensei ạ. Mọi người đều đang cố gắng hết sức để lớp chúng ta có một kỷ niệm thật đẹp.''

Yuko-sensei nhìn quanh quán một lượt, ánh mắt đong đầy sự tự hào rồi khẽ dặn dò chúng tôi phải chú ý nghỉ ngơi và uống nước đầy đủ trước khi rời đi để tiếp tục công việc giám sát của mình.

Đồng hồ trên tường điểm mười giờ rưỡi. Bầu không khí trong quán vẫn đang ở mức nhiệt độ cao nhất với lượng khách ra vào không ngừng nghỉ. Đúng lúc này, một bóng dáng quen thuộc xuất hiện ở cửa - là Inari. Cô nàng cũng diện bộ đồ hầu gái nhưng được phối thêm dải băng điều phối màu đỏ rực ở bắp tay.

``Inari! Bà tới rồi!'' Hanabi reo lên đầy phấn khích, chạy lại ôm chầm lấy cô bạn thân của mình.

Miku ngạc nhiên nhìn cô nàng: ``Bọn mình cứ ngỡ cậu phải bận rộn với việc điều phối chung cho toàn trường cùng Hội học sinh chứ?''

Inari mỉm cười, chỉnh lại dải băng: ``Đúng là mình đang phụ trách việc đó, nhưng Hội trưởng và Hội phó đang làm quá tốt ở khu vực cổng chính rồi. Mình không thể ngồi yên nhìn lớp 1-C chiến đấu một mình được, nên mình đã quyết định xin phép được tham gia vào đội hình hỗ trợ vào phút cuối. Dù sao thì mình cũng là thành viên của lớp mà!''

Sự xuất hiện của Lớp trưởng giống như một luồng gió mới tiếp thêm năng lượng cho cả nhóm đang bắt đầu thấm mệt. Ngay sau đó, ca trực đầu tiên chính thức kết thúc để chuyển giao.

``\vie{Kaigou-chan, Suzuki, hai người đi nghỉ đi,}'' Onii-sama lên tiếng, tay vẫn đang nhanh nhẹn sắp xếp lại các tách trà trên khay một cách điệu nghệ. ``\vie{Inari-san và Suzune-san sẽ lên thay vị trí của hai đứa ở sảnh và quầy thu ngân.}''

Tôi nhìn anh ấy, không khỏi lo lắng khi thấy những giọt mồ hôi lấm tấm trên thái dương anh: ``\vie{Còn anh thì sao, Onii-sama? Anh không nghỉ ca này sao?}''

Anh khẽ lắc đầu, giọng nói vẫn bình thản và chuyên nghiệp như một quản gia thực thụ: ``\vie{Anh ổn. Công việc này không quá nặng nhọc với anh. Vả lại, nếu anh rời đi thì hiện tại không có ai thay thế được. Cả lớp 1-C chỉ có mình anh là quản gia thôi, còn ai khác nữa đâu. Anh sẽ nghỉ sau.}''

Biết không thể lay chuyển được ý chí của anh, tôi và Onee-sama đành tháo tạp dề, chuẩn bị cho ca nghỉ của mình. ``\vie{Vậy bọn em sẽ mang cái gì đó thật ngon về cho anh nhé, Onii-sama!}'' tôi vẫy tay chào anh trước khi bị Onee-sama kéo đi.

``\vie{Đi thôi Suzu. Chị muốn xem cái 'Con đường Linh hồn' của Shion-san rốt cuộc là cái gì mà mọi người cứ kháo nhau kinh dị lắm,}'' chị nói, ánh mắt lộ rõ vẻ tò mò pha chút thách thức. ``\vie{Ít nhất thì chị cũng giải phóng được việc phải tiếp mấy gã đàn ông ghét mắt đó.}''

Ngay lúc chúng tôi vừa rời khỏi lớp, bất chợt\dots

``Suzuki-oneechan! Kaigou-oneechan!'' một giọng nói trong trẻo bất ngờ vang lên từ phía sau. Naomi chạy ào tới, bộ váy Công chúa Hoa tung bay, đôi cánh tiên lấp lánh sau lưng. ``\vie{Em cũng được Onii-chan cho nghỉ rồi! Em muốn đi chơi nhà ma với hai chị!}''

Tôi ngạc nhiên quay lại, cúi xuống nhìn Naomi: ``\vie{Nhưng nếu Onii-sama cho em nghỉ thì ai thay vị trí của em?}''

Naomi mỉm cười toe toét, đôi mắt híp lại thành hai vầng trăng khuyết: ``\vie{Các chị Saya-san, Hotaru-san và Mitsuki-san sẽ quán xuyến mà! Các chị nói em cứ đi chơi thoải mái đi, công việc ở đó đã có họ lo rồi!}''

Tôi nhìn Onee-sama, thấy chị cũng đang chau mày suy nghĩ. Có lẽ chị cũng đang cân nhắc xem có nên để Naomi đi cùng không, nhất là sau khi chúng tôi vừa nghe những lời đồn đại về ``Con đường Linh hồn'' của Shion - nơi được cho là đáng sợ nhất lễ hội. Dẫu rằng em ấy đã từng là một hồn ma, nhưng đâu chắc được em ấy sẽ không sợ hãi khi đối mặt với những trò hù dọa ghê rợn chứ?

Onee-sama khẽ nhíu mày: ``\vie{Nao, nhà ma đáng sợ lắm đấy. Em chắc chắn muốn vào chứ?}''

Naomi gật đầu lia lịa, đôi mắt to tròn sáng rực: ``\vie{Em không sợ đâu! Có hai chị ở đây mà!}''

Tôi mỉm cười, nắm lấy tay em. Và thế là, đội hình ba chị em chính thức bước vào ``Con đường Linh hồn.''

Chúng tôi bước về phía cuối hành lang tầng hai, nơi một tấm rèm đen dày đặc và cũ kỹ đang che phủ lối vào lớp học để trống được trưng dụng làm nhà ma. Một cảm giác lạnh lẽo bất thường thoát ra từ phía sau tấm rèm, mang theo mùi của nhang trầm và gỗ mục, khiến tôi rùng mình dù trời đang là giữa mùa hè rực nắng.

Vừa bước qua cửa, bóng tối lập tức bủa vây lấy chúng tôi một cách tuyệt đối. Không phải là cái bóng tối bình thường, mà là một sự đen đặc quánh lại, khiến người ta mất đi cảm giác về không gian và phương hướng ngay lập tức.

``Chào mừng các lữ khách đã đến với vực thẳm,'' giọng nói của Game vang lên từ hệ thống loa âm thanh vòm được giấu kín khắp căn phòng. ``Tại đây, những ký ức mà các ngươi muốn chôn giấu nhất sẽ được phơi bày dưới ánh sáng của sự tuyệt vọng. Hãy cẩn thận với từng bước chân, vì ranh giới giữa thực và ảo đang mỏng manh hơn bao giờ hết.''

Dù biết rõ đây chỉ là chiêu trò phối hợp giữa kỹ thuật của Game và sự dàn dựng của Shion, nhưng không hiểu sao tôi vẫn cảm thấy sống lưng mình lạnh toát. Một luồng gió buốt giá bất ngờ thổi qua từ phía dưới chân, mang theo những sợi tơ mảnh nhện giả quấn quýt lấy cổ chân tôi. Những tiếng thì thầm vô hình bắt đầu vang lên từ khắp phía, lúc gần lúc xa, gọi tên chúng tôi bằng những thanh âm méo mó.

``Chỉ là kỹ xảo thôi, Suzu. Nắm chặt tay chị, đừng có buông ra,'' Onee-sama nói, nhưng tôi có thể cảm nhận được bàn tay chị cũng hơi siết chặt lại.

Và rồi, thứ âm thanh đó vang lên.

Tiếng cười. Khe khẽ, đều đều, ngắt quãng. Vọng ra từ phía sau lưng chúng tôi, như thể có ai đó đang bám sát gót chân mà chúng tôi không hề hay biết.

Tim tôi hẫng một nhịp. Tiếng cười ấy\dots\ tôi đã từng nghe nó. Trong những hành lang tối tăm của trường Aogami cũ, nơi mà những bóng ma của các bạn học từng thoắt ẩn thoắt hiện, cười không thành tiếng, nhìn mà không thấy. Ký ức về lần đầu tiên về Reiko - hồi tôi còn chưa nhớ về bản thân - nghe thấy tiếng cười quái dị ấy vọng ra từ một lớp học trống, rồi chứng kiến năm bóng người đứng bất động bên trong, bỗng chốc ùa về như dòng thác lạnh.

``\vie{Chị có nghe thấy không\dots?}'' tôi thì thầm, giọng run rẩy.

Onee-sama không trả lời ngay. Bàn tay chị siết chặt đến mức các đốt ngón tay trắng bệch. Tôi hiểu: chị cũng đang nghe thấy. Và chị cũng đang nhớ. Nhớ về những lần chạy trối chết trên hành lang, khi những cánh tay đen ngòm vươn ra từ mặt sàn nứt toác, túm lấy chân Onii-sama và kéo anh ấy xuống vực. Nhớ về cảm giác tuyệt vọng khi nhìn những người bạn biến thành vỏ rỗng, lặp lại cùng một câu nói như con rối bị giật dây.

Bỗng nhiên, một luồng sáng strobe xanh loét chớp lên đột ngột, hé lộ những bóng hình dị dạng đang treo lơ lửng trên trần nhà. Đó là những con búp bê vải mà nhóm của Mari đã dày công làm, nhưng dưới ánh sáng lập lòe và hiệu ứng khói đặc, chúng trông không khác gì những cái xác khô đang oằn tà oại trong đau đớn.

``Gyaaah!'' Naomi hét lên, nhảy phắt lên ôm chầm lấy eo Onee-sama. Chị tôi giật mình, nhưng lập tức đưa tay che cho em bé, phản xạ bảo vệ của một người chị lớn lập tức trỗi dậy.

``\vie{Không sao, Nao. Chỉ là đồ giả thôi,}'' chị nói, giọng vẫn bình tĩnh, nhưng tôi thấy cơ hàm chị hơi cứng lại. Bởi vì trong khoảnh khắc ánh đèn chớp nháy, những bóng hình treo lơ lửng kia gợi nhớ đến cảnh tượng kinh hoàng mà chúng tôi từng chứng kiến ở hành lang ngôi trường cũ: trần nhà võng xuống, từng mảng gạch rơi phịch, bụi mù mịt, và sau đó là những bóng người hiện lên rồi tan biến trong nháy mắt.

Càng đi sâu vào trong, căn phòng càng trở nên chật hẹp và ẩm thấp. Lối đi thu hẹp lại thành một hành lang dài và tối, hai bên là những cánh cửa trượt bị che kín bằng giấy đen. Tiếng cào cấu vào gỗ vang lên kèn kẹt ngay sát bên tai, hòa cùng tiếng khóc thút thít vọng ra từ một chiếc tủ cũ nát. Cảm giác quen thuộc đến đau lòng: giống y hệt những hành lang lặp đi lặp lại, nơi mọi cánh cửa đều bị khóa, mọi lối đi đều dẫn về cùng một chỗ.

``Này, Suzuki-user, Kaigou-user\dots'' giọng Game bỗng hạ thấp, chỉ đủ cho hai chúng tôi nghe. ``Có quen không? Hành lang dài, cửa đóng, và cảm giác bị theo dõi\dots\ Ta chỉ tái hiện lại một phần nhỏ thôi, nhưng có vẻ hai cô đang phản ứng mạnh hơn dự kiến đấy.''

``Ông già im đi, `` Onee-sama gằn giọng. Nhưng tôi thấy bước chân chị nhanh hơn, gấp gáp hơn, đúng cái nhịp mà chúng tôi từng chạy khi ngôi trường bắt đầu sập, khi mặt đất nứt toác, khi cánh tay đen ngòm vươn ra từ bóng tối.

Một cánh tay trắng nhợt và lạnh ngắt từ bóng tối bất ngờ vươn ra, định chạm vào vai tôi. Tôi rít lên, né sang một bên và nép sát vào người chị. Tôi nhận ra đó là Aku đang hóa trang thành một hồn ma không đầu, động tác di chuyển của cô ấy dặt dẹo và khó đoán đến mức đáng kinh ngạc.

Nhưng trong thoáng chốc, não tôi không xử lý được đó là Aku. Nó chỉ thấy: cánh tay trắng, bóng tối, và cảm giác bị túm lấy, y hệt lúc Onii-sama bị kéo xuống hố đen trong ngôi trường cũ, y hệt lúc Onee-sama hét lạc giọng rồi lao theo anh xuống vực sâu. Tim tôi đập dồn dập đến mức tai ù đi.

``Sao thế?'' Game lại lên tiếng, giọng mỉa mai. ``Những kẻ đã từng đối đầu với Thần linh và thực thể tà ác mà lại sợ mấy cái bóng này sao? Thật là một sự mỉa mai đầy thú vị.''

``\vie{Không phải sợ,}'' tôi thì thầm, nắm chặt tay Naomi đang run lẩy bẩy. ``\vie{Chỉ là\dots\ nó nhắc em nhớ lại thôi.}''

Onee-sama im lặng. Nhưng cách chị kéo tôi và Naomi sát lại, che chắn bằng cả cơ thể mình, đã nói lên tất cả. Chị cũng đang nhớ những lần về lần mình đấm nát rễ cây để giải cứu Onii-sama trong khu rừng Aogami, nhớ về cảm giác bất lực khi nhìn Hanabi và Inari biến thành những con rối vô hồn lặp lại cùng một câu nói.

Rồi chúng tôi bước vào một căn phòng rộng hơn, và ở đó, năm bóng người đứng rải rác, bất động, khuôn mặt bị che khuất bởi bóng tối và khói. Một cô gái tóc dài sát khung cửa sổ. Một cô gái hai bím tóc nghiêng đầu nhìn ra ngoài. Một bóng nhỏ có đôi tai thú trên đầu. Một cô gái tóc cột lệch xoay lưng. Và một người có hai búi tóc nhỏ nhô lên.

Mọi thứ đông cứng lại. Tôi nghe thấy nhịp tim mình gõ từng tiếng rõ mồn một trong lồng ngực trống rỗng.

Năm cái bóng. Đứng im. Cười không thành tiếng.

Giống. Quá. Giống.

``\vie{Chị\dots}'' giọng tôi lạc đi. ``\vie{Chị có thấy không\dots?}''

Onee-sama nghiến răng. Tôi thấy vai chị run lên, không phải vì sợ, mà vì cơn giận bùng lên khi ký ức đau đớn bị khoáy vào. ``\vie{Có,}'' chị đáp khẽ. ``\vie{Nhưng đây không phải là `họ'. Đây chỉ là lễ hội. Bình tĩnh lại, Suzu.}''

Naomi ngước lên nhìn hai chị, đôi mắt to tròn đầy lo lắng: ``\vie{Kaigou-oneechan, Suzuki-oneechan\dots\ hai chị ổn không? Sao mặt hai chị tái mét vậy?}''

Tôi cố nặn ra một nụ cười: ``\vie{Ổn mà, Nao. Chị chỉ\dots\ hơi giật mình thôi.}''

*XOẸT!*

Đèn chớp. Năm bóng người biến mất. Chỉ còn căn phòng trống, búp bê vải nằm la liệt trên sàn. Tôi thở phào, nhưng đôi chân vẫn còn run.

Tôi phải thừa nhận, nếu so với những trận chiến sinh tử để khôi phục lại linh khí cho thị trấn trước đây, thì những thứ này về mặt logic chẳng thấm vào đâu. Nhưng cái cách mà Shion kết hợp giữa yếu tố tâm linh thực sự và kỹ thuật đánh lừa thị giác của Azuki thực sự khiến người ta phải ``rén''. Nó không tấn công vào thể xác, mà nó bóp nghẹt tâm lý bằng sự chờ đợi và những cái bẫy cảm giác không hồi kết. Và với những người như chúng tôi - những kẻ mang theo vết sẹo của vòng lặp - nó đánh thẳng vào chỗ đau nhất: ký ức.

Chúng tôi đi qua một đoạn hành lang đầy rẫy những chiếc gương méo mó, nơi hình ảnh phản chiếu của chúng tôi trông như những con quái vật đang nhe răng cười. Onee-sama khẽ cau mày, tay chị nắm tay tôi chặt hơn đến mức tôi có thể cảm nhận được mạch đập của chị. Có vẻ như ngay cả chị cũng đang cảm thấy không thoải mái với bầu không khí quỷ dị này.

``Các ngươi cũng khá lắm,'' giọng Shion bất ngờ vang lên từ đâu đó trong bóng tối, mang theo chất chuunibyou đặc trưng nhưng pha lẫn sự chân thành. ``Không nhiều kẻ đi được đến đây mà vẫn giữ được bình tĩnh đâu. Hầu hết đều bỏ chạy từ đoạn thứ hai rồi.''

``Bình tĩnh cái quái gì,'' Onee-sama hừ một tiếng. ``Tôi chỉ đang kìm không đấm vỡ mấy cái búp bê kia thôi.''

Ở đoạn cuối của cuộc hành trình, cô ấy xuất hiện như một nữ vương bóng tối thực thụ, đứng giữa một vòng tròn nến đen huyền bí. Cô ấy không cần nhảy ra dọa dẫm, chỉ cần đứng im với đôi mắt nhắm nghiền và lẩm bẩm những câu chú bằng tiếng Latin cổ trong làn khói tím nhạt, cũng đủ khiến bầu không khí đặc quánh lại như thể có một áp lực vô hình đang đè nặng lên ngực.

Naomi, lúc này đã quên sợ, ngước lên nhìn Shion bằng đôi mắt lấp lánh: ``Shion-san trông ngầu quá! Giống phù thủy trong truyện cổ tích luôn!''

Shion hé một mắt, khóe môi khẽ cong lên trước lời khen hồn nhiên của cô bé, nhưng lập tức lấy lại phong thái uy nghiêm: ``Hừ, đương nhiên rồi. Ta là chúa tể của bóng tối mà.''

Khi chúng tôi cuối cùng cũng thoát ra ngoài qua tấm rèm ở phía đối diện, ánh nắng rực rỡ và tiếng ồn ào náo nhiệt của sân trường khiến tôi phải nheo mắt lại vì chưa kịp thích nghi. Tôi hít một hơi thật sâu, cảm thấy phổi mình như vừa được lọc sạch bụi bặm của cõi âm ti.

``\vie{Phù\dots\ Shion-san và Game đúng là một sự kết hợp chết chóc,}'' tôi thở phào, đưa tay lau mồ hôi lạnh đang lấm tấm trên trán.

``\vie{Không thể ngờ là ông ta có thể tái tạo được những trải nghiệm kinh khủng đó một cách chân thực đến vậy,}'' chị tôi lẩm bẩm, ``\vie{Đúng là đáng sợ mà.}''

Naomi thì ngược lại hoàn toàn, cô bé nhảy chân sáo, vẻ phấn khích: ``\vie{Vui quá! Em muốn đi lại lần nữa!}''

``\vie{\textbf{Tuyệt đối không!}}'' tôi và Onee-sama đồng thanh.

Chị tôi đứng chỉnh lại mái tóc và bộ trang phục cho ngay ngắn, dù môi hơi mím lại nhưng chị vẫn phải lên tiếng công nhận: ``Hừ, cũng không tệ. Coi như là một trải nghiệm thú vị để tỉnh táo hơn trước khi đi tìm món gì đó lót dạ vậy.''

Tôi mỉm cười nhìn chị. Dù có thực sự cảm thấy sợ hãi một chút, nhưng đó là một cái sợ đầy phấn khích - cái sợ đặc trưng của một lễ hội trường học thực sự mà chúng tôi đang được tận hưởng trọn vẹn cùng nhau như những nữ sinh bình thường. Một cảm giác thật xa xỉ nhưng cũng thật tuyệt vời.

Nhưng sâu trong lòng, tôi biết cả tôi và Onee-sama đều đang nghĩ cùng một điều: những ký ức ấy - hành lang vô tận, tiếng cười quái dị, bóng ma bạn bè, cánh tay kéo xuống vực - chúng không bao giờ thực sự biến mất. Chúng chỉ ngủ yên, chờ một khoảnh khắc để thức dậy. Và hôm nay, trong ngôi nhà ma của Shion, chúng đã thức dậy thực sự, dù chỉ trong tích tắc.

Nhưng điều khác biệt là: lần này, chúng tôi không chạy trốn. Chúng tôi đã đi qua nó, tay nắm tay, với Naomi hồn nhiên ở giữa. Và điều đó, theo một cách nào đó, đã khiến những vết sẹo cũ bớt nhức nhối hơn một chút.

Chúng tôi bắt đầu rảo bước về phía sân trường, nơi những gian hàng ẩm thực đang tỏa hương thơm phưng phức, chuẩn bị cho một hành trình khám phá mới trong ca nghỉ của mình.

\pov{Hikawa Kaigou}

Ánh nắng rực rỡ của giữa trưa đổ xuống sân trường Aogami, làm bừng sáng những tấm áp phích đầy màu sắc và những dải cờ đuôi nheo giăng khắp các lối đi. Bầu không khí ngoài này khác hẳn với cái thứ tăm tối, ngột ngạt mà chúng tôi vừa nếm trải trong ``Con đường Linh hồn''.

Tôi vươn vai, hít một hơi căng tràn lồng ngực mùi thơm của bạch tuộc nướng, của bơ đường và của những thanh yakisoba xèo xèo trên chảo thép.

Sau khi rời khỏi căn nhà ma của Shion, Naomi đã chính thức trở thành ``cái đuôi'' của hai chị em tôi. Cô bé lanh lợi bảo rằng Saya, Hotaru và Mitsuki đã quay về câu lạc bộ của họ để chuẩn bị cho các hoạt động riêng, nên em ấy muốn được đi cùng chúng tôi trong lúc nghỉ. Đương nhiên là tôi chẳng có lý do gì để từ chối. Nhìn con bé mặc bộ váy Công chúa Hoa tung tăng đi giữa hai chị em, thi thoảng lại kéo tay Suzu chỉ trỏ vào mấy gian hàng, lòng tôi cũng bất giác thấy nhẹ nhõm.

*Ting tong, ping pong\~{}*

Tiếng chuông báo phát thanh vang lên, cắt ngang tiếng ồn ào náo nhiệt của sân trường. Từ hệ thống loa, giọng nói trong trẻo và đầy năng lượng của Kana vang vọng khắp mọi ngóc ngách:

``A lô, a lô! Xin chào tất cả các học sinh và quý khách tham quan lễ hội văn hóa trường Aogami! Mình là Tokiwa Kana đây\~{}! Các bạn đã ăn trưa chưa nào? Đừng quên ghé qua khu vực sân sau để thưởng thức các món ăn tuyệt hảo do chính tay các đầu bếp nghiệp dư của trường chế biến nhé! Gian hàng takoyaki của lớp 2-B đang có chương trình mua năm tặng một đó nha! Nhanh chân lên nào, số lượng có hạn!''

Tôi bật cười, quay sang Suzu: ``\vie{Nghe hấp dẫn đấy chứ. Chị em mình đi kiếm gì bỏ bụng thôi. Bữa sáng tiêu hóa hết sạch vào mấy cái màn hù dọa của ông già Game rồi.}''

``\vie{Vâng!}'' Suzu gật đầu, mắt cũng sáng rực lên. ``\vie{Em muốn ăn bánh crepe! Hồi sáng em thấy có lớp 1-A làm crepe dâu tây trông ngon lắm!}''

``Em cũng muốn ăn bánh crepe! Em muốn ăn kem nữa!'' Naomi giơ cao tay hưởng ứng, đôi cánh tiên sau lưng rung lên bần bật vì phấn khích.

Thế là ba chúng tôi bắt đầu hành trình càn quét khu ẩm thực.

Nói là ``càn quét'' thì hơi nhẹ, phải gọi là ``bão càn'' mới đúng. Chỉ trong chưa đầy mười lăm phút, Suzu đã đánh bay hai chiếc bánh crepe dâu tây cỡ bự, một xiên kẹo táo, và hiện tại đang nhóp nhép nhai nốt phần karaage nóng hổi. Con bé ăn với tốc độ ánh sáng mà hai má vẫn phồng lên đáng yêu, không hề có dấu hiệu ngán hay no.

``\vie{Từ từ thôi, Suzu!}'' tôi vuốt mồ hôi hột trên trán, tay xách lỉnh kỉnh ba túi giấy đựng đồ ăn vặt. Tiền thì chúng tôi không thiếu, thẻ đen của Kuon-kun tôi vẫn đang cầm đây, nhưng vấn đề là\dots\ ``\vie{Em mà mua hết sạch thế này thì những người xếp hàng phía sau lấy gì mà mua nữa? Lát nữa họ lại tưởng trường mình có quái vật phàm ăn giáng trần đấy!}''

``\vie{Nhưng mà ngon lắm Onee-sama!}'' Suzu chớp chớp mắt, đưa xiên gà rán còn lại ra trước mặt tôi. ``\vie{Lâu lắm rồi em mới được ăn đồ ăn tẩm ướp đậm đà thế này. Ở thế giới cũ làm gì có mấy thứ này đâu\dots\ Ăn bao nhiêu cũng không sợ béo, tội gì không ăn! Ưm!}''

``\vie{Thế giới cũ?}'' Naomi đang mút cây kẹo bông gòn khổng lồ màu xanh, ngơ ngác ngước lên nhìn. ``\vie{Đó là đâu vậy ạ?}''

Tôi vội giật lấy xiên gà từ tay Suzu, nhét vào miệng mình để đánh trống lảng: ``\vie{À, ý chị ấy là hồi bọn chị còn ở quê ấy mà. Quê bọn chị nghèo lắm, không có lễ hội to thế này đâu. Đúng không Suzu?}''

Em tôi gật đầu lia lịa, lại chạy tót sang gian hàng bán yakisoba. Tôi chỉ biết thở dài ngao ngán, điềm tĩnh xách thêm hai hộp takoyaki và yakisoba (tôi không quên mua cho Kuon-kun nữa) đi theo sau. Nhưng thật lòng mà nói\dots\ nhìn con bé vui vẻ như vậy, tôi cũng đang rạo rực không kém. Đây là lễ hội trường học đầu tiên trong cuộc đời cao trung của tôi cơ mà! Một lễ hội đúng nghĩa, không có quái vật, không có khói lửa tàn tích, cũng không có những nhiệm vụ giải cứu thế giới vắt kiệt sức lực. Chỉ có tiếng cười, đồ ăn ngon và bầu không khí thanh xuân rực rỡ.

Khi chúng tôi đang đứng trước quầy bán đá bào để chờ lấy phần cho Naomi, một giọng nói dịu dàng quen thuộc vang lên từ phía sau.

``Kaigou-san, Suzuki-san!''

Chúng tôi quay lại. Đứng đó là Shino. Cô nàng đang mặc bộ đồng phục mùa hè của trường, trên tay cầm một túi giấy nhỏ in hình linh vật lễ hội. Trông cô ấy toát lên một vẻ thanh tao nhưng cũng rất gần gũi, nét mặt tươi tắn hơn những ngày học bình thường kể từ khi chúng tôi bắt đầu sinh sống tại đây.

``Ồ, Shino-san,'' tôi mỉm cười đáp lại. ``Cậu đang đi nghỉ giữa giờ à? Công việc điều phối bên Hội Học sinh xong rồi sao?''

Shino gật đầu, bước tới gần chúng tôi với nụ cười nhẹ nhõm: ``Mình vừa bàn giao lại công việc cho các thành viên khác. Bây giờ là thời gian nghỉ của mình. Thật tình cờ lại gặp hai cậu ở đây.''

Suzu vẫy tay chào, tay kia đưa cốc đá bào si-rô việt quất cho Naomi: ``Shino-san! Cậu đi một mình sao? Nếu không phiền thì đi chung với bọn mình luôn nhé? Bọn mình đang định đi ăn xong rồi lượn qua xem mấy gian hàng trò chơi.''

Mắt cô nàng tóc nâu khẽ chớp, nụ cười trên môi cô ấy dường như rạng rỡ hơn một chút: ``Được sao? Nếu vậy thì\dots\ mình đi cùng nhé. Thực ra mình cũng đang đi dạo loanh quanh vì không biết nên bắt đầu vui chơi từ đâu nữa.''

Tôi vỗ vai cô ấy, nhếch môi tự tin: ``Cứ đi theo chị đại đây, đảm bảo vui nổ trời.''

Và thế là đội hình ba người chúng tôi giờ có thêm Shino. Bốn chị em chúng tôi kéo nhau đi tìm một chiếc ghế đá trống dưới gốc cây anh đào cổ thụ mát mẻ để xử lý đống đồ ăn vừa mua.

Suzu cắn một miếng bánh crepe dâu tây, lớp kem tươi dính một chút trên khóe môi, hai má phồng lên nhai nhóp nhép trông rõ yêu. Naomi thì bận rộn múc từng muỗng đá bào, thỉnh thoảng lại thè lưỡi ra khoe màu xanh lè khiến cả bọn phì cười. Tôi mở hộp takoyaki còn nóng hổi, dùng tăm ghim lấy một viên rồi đưa sang cho Shino.

``Cảm ơn cậu,'' cô ấy nhận lấy, thổi nhẹ lớp cá bào katsuobushi đang nhảy múa trên mặt bánh rồi cắn một miếng nhỏ. Biểu cảm trên khuôn mặt cô ấy dãn ra, ánh mắt ánh lên sự thích thú: ``Ngon thật đấy. Hương vị đậm đà hơn mình nghĩ.''

Tôi dựa lưng vào thành ghế, nhấp một ngụm trà lúa mạch lạnh ngắt: ``Tất nhiên rồi. Mùi vị của lễ hội mà. Cậu thấy sao, Shino-san? Bầu không khí của trường Aogami ngày hôm nay ấy?''

Shino khẽ dừng lại, ánh mắt cô ấy hướng về phía sân trường đầy người qua lại. Tiếng nhạc xập xình từ xa vọng tới, tiếng gọi nhau í ới của các học sinh đang phát tờ rơi, và cả những tràng cười giòn tan vang lên không ngớt từ các gian hàng.

``Tuyệt vời lắm.'' Cô khẽ nói, giọng cô ấy pha chút hoài niệm, nhưng không phải là sự hoài niệm đau thương. ``Từ trước đến nay, mình luôn phải sống trong cảm giác tủi thân, bị đè nén bởi những vòng lặp vô tận. Đã có những lúc, khi nhìn ra cửa sổ lớp học vắng lặng, mình từng nghĩ một cuộc sống học đường bình thường, nhộn nhịp như thế này chỉ là một giấc mơ xa xỉ không bao giờ có thể chạm tới.''

Tôi nhìn cô ấy, bàn tay đang cầm hộp bánh hơi khựng lại. Suzu cũng dừng ăn, chăm chú lắng nghe, trong mắt dâng lên một sự thấu cảm sâu sắc. Hơn ai hết, chúng tôi hiểu rõ cảm giác đó.

Shino quay sang nhìn chúng tôi, nụ cười của cô ấy giờ đây trong trẻo lạ thường: ``Nhưng bây giờ, được ngồi đây cùng các cậu, ăn đồ ăn vặt, nghe tiếng mọi người cười nói\dots\ Mình cảm thấy như mình đang thực sự sống. Một cuộc sống đúng nghĩa, không còn sợ hãi hay lo âu. Mình\dots\ thực sự rất vui, Kaigou-san, Suzuki-san.''

Suzu mỉm cười dịu dàng, khẽ gật đầu: ``Mình cũng vậy. Bọn mình đều đã đi một chặng đường rất dài và khó khăn để có được ngày hôm nay. Vậy nên, cậu cứ tận hưởng nó cho thỏa thích nhé, Shino-san.''

Tôi bật cười, chọc nhẹ vào eo Shino khiến cô nàng giật mình rụt người lại: ``Đúng rồi đó. Bỏ ba cái suy nghĩ sâu xa đi! Lát nữa tôi sẽ dẫn cậu đi phá đảo mấy cái gian hàng mini-game. Để xem năng lực của cậu khi cầm vòng ném chai có bách phát bách trúng như lúc dùng pháp thuật không!''

Shino bật cười thành tiếng, một nụ cười rạng rỡ không vướng bận: ``Được thôi. Mình không dễ thua đâu nhé.''

*XẸT XẸT*

Hệ thống loa lại bật mở. Lần này, thay thế cho chất giọng nhí nhảnh của Kana là chất giọng trầm ấm, điềm đạm và trưởng thành của Kaoru:

``Kính thưa quý vị và các bạn học sinh, đây là Ukai Kaoru. Hiện tại, Câu lạc bộ Nghiên cứu Khoa học ở tầng ba dãy nhà B đang có màn trình diễn phản ứng hóa học phát sáng cực kỳ đẹp mắt. Các bạn yêu thích sự kỳ diệu của khoa học đừng bỏ lỡ nhé. Bên cạnh đó, gian hàng bói toán `Vén Màn Tương Lai' của lớp 3-A tại dãy nhà C cũng đã bắt đầu mở cửa đón khách. Xin nhắc lại\dots''

Tôi đứng phắt dậy, phủi nhẹ bụi trên váy, hai tay chống hông đầy khí thế: ``Nghe thấy chưa? Lớp 3-A có bói toán kìa. Đám năm ba chắc có nhiều trò hay ho và lừa tình lắm. Có muốn đi để bắt bài họ không?''

Suzu nhanh chóng nhét cốc rỗng vào thùng rác, lau miệng bằng khăn giấy, rồi cầm lấy tay Naomi: ``Shino-san, mình có bội rồi nhé?'' Shino gật đầu, nhìn em ấy, ánh mắt dịu dàng: ``Cậu nhỏ cũng đi cùng phải không?''

``Có ạ! Em muốn xem liệu em có thể cao lên được không!'' cô nhóc hồn nhiên đáp khiến cả bọn cười vang.

Chúng tôi cùng nhau rời khỏi khu ẩm thực, tiến về phía dãy nhà C. Dọc hành lang, các lớp học đã được biến hóa một cách kỳ công, rực rỡ sắc màu. Chúng tôi ghé qua một lớp làm quán cà phê phong cách rừng rậm, nơi vài nam sinh mặc đồ thú nhồi bông cồng kềnh đi lại lạch bạch, vẫy tay chào mọi người.

Tiếp đó, đúng như lời thách thức, tôi kéo cả bọn vào gian hàng ném vòng vào chai của khối năm hai. Tôi tiêu tốn mất tận hai trăm yên, ném trượt mấy lần đầu vì căn lực sai, nhưng nhờ sự cổ vũ rộn ràng của Suzu và Naomi, cú ném cuối cùng của tôi đã trúng phóc vào cái chai mang phần thưởng.

Tôi hãnh diện vớt vát lại được một chiếc móc khóa hình gấu trúc béo ú để tặng cho Naomi, khiến cô bé vui sướng reo hò nhảy cẫng lên.

Shino cũng thử sức, và dù cú ném đầu tiên hơi lệch, cô ấy nhanh chóng điều chỉnh lại và ném trúng ngay cú thứ hai. Mắt cô ấy sáng rực: ``Mình\dots\ trúng thật sao?''

``Trúng rồi đó! Cậu giỏi hơn tôi đấy!'' tôi cười lớn, vỗ vai cô nàng. Suzu cũng vỗ tay reo hò. Cô ấy nhận phần thưởng là một chiếc kẹp tóc hình ngôi sao, và ôm nó như thể đó là chiến tích quý giá nhất đời mình. Nhìn biểu cảm hạnh phúc thuần khiết trên khuôn mặt cô ấy, tôi chợt nghĩ: Shino cũng giống chúng tôi, lần đầu tiên được trải nghiệm những niềm vui nhỏ bé mà lẽ ra đã phải có từ lâu.

Đến phòng học của lớp 3-A, không gian bên trong được làm mờ ảo với những tấm vải voan tím và tiếng nhạc mang âm hưởng Trung Đông huyền bí. Một chị năm ba ăn mặc như nữ thầy bói du mục ngồi sau quả cầu thủy tinh phát sáng, nghiêm trang lật từng lá bài tarot cho chúng tôi.

Khi bói cho Naomi, chị ấy mỉm cười hiền hậu: ``Em bé này có một tâm hồn rất trong sáng. Xung quanh em có nhiều người yêu thương và che chở em. Em sẽ được khỏe mạnh và lớn lên thật cao đẹp nhé.'' Cô bé reo lên, quay sang ôm chầm lấy Suzu: ``\vie{Suzuki-oneechan nghe thấy không? Em sẽ cao lên được!}'' Suzu bật cười, xoa đầu em: ``\vie{Ừ, chị nghe thấy rồi!}''

Khi bói cho Shino, chị ấy lật lá bài ``The Sun,'' ngẩng lên nhìn cô với ánh mắt trầm trồ: ``Sau một chặng đường dài đi trong sương mù, mặt trời cuối cùng cũng đã lên. Tương lai của em sẽ rực rỡ và đong đầy sự bình yên.'' Shino nghe xong chỉ khẽ mỉm cười, một nụ cười thấu hiểu. Cô ấy không nói gì, nhưng đôi mắt ánh lên thứ gì đó sâu lắng, có lẽ là sự biết ơn đối với một thế giới cuối cùng đã cho cô cơ hội sống bình thường.

Suzu bói xong, chị thầy bói cười tươi: ``Em là người rất giỏi quan sát và cảm nhận. Sự dịu dàng của em là thế mạnh, nhưng đừng quên rằng em cũng có quyền được chiều chuộng, không phải lúc nào cũng phải là người lo lắng cho người khác.'' Con bé hơi ngượng, gãi gãi má. ``\vie{Đúng là Onii-sama hay nói vậy\dots}'' cô bé lẩm bẩm.

Cuối cùng là đến lượt tôi. Bà chị lật bài, nhíu mày một lúc, rồi ngẩng đầu lên nhìn thẳng vào mắt tôi: ``Em có một bức tường rất dày bao quanh trái tim mình. Bức tường đó được xây lên bởi những vết thương cũ, và nó giúp em sống sót qua rất nhiều thứ. Nhưng\dots'' chị ấy dừng lại, chạm nhẹ vào lá bài, ``không phải ai cũng muốn làm em đau. Hãy cho phép bản thân tin rằng có những người thực sự đáng để mở lòng.''

Tim tôi nhói lên.

Không phải vì câu phán đó hay ho hay chính xác gì cả. Mà vì nó chạm đúng vào cái thứ mà tôi đã chôn sâu nhất: sự ghê tởm, không, đúng hơn là sự sợ hãi đối với đàn ông. Ở thế giới cũ, chúng là những kẻ đã đẩy tôi và Suzu vào vực thẳm. Từ những gã ngon ngọt đường mật chỉ để lợi dụng, cho tới những gã cặn bã không coi chúng tôi ra gì từ cái xã hội mục nát mà chúng tôi phải tự cắn răng chịu đựng. Tôi xây bức tường đó để sinh tồn, và nó đã bảo vệ chúng tôi tốt đến mức tôi gần như quên mất rằng nó cũng đang giam giữ chính tôi.

Nhưng rồi, Kuon-kun xuất hiện. Một cậu con trai cũng lạc lõng như chúng tôi, cũng bị quăng vào một thế giới xa lạ, nhưng thay vì lợi dụng hay bỏ rơi, cậu ta đã chọn đồng hành, chọn bảo vệ. Cậu ta đã hứa với tôi ở The Void, rằng cậu ta sẽ không bao giờ trở thành thứ mà tôi ghê sợ. Và cho đến giờ phút này, cậu ta vẫn đang giữ lời hứa đó.

``\dots Cảm ơn chị,'' tôi đáp khẽ, cố giữ giọng bình thản, nhưng Suzu bên cạnh đã nhận ra. Con bé khẽ nắm lấy tay tôi dưới bàn, bóp nhẹ một cái. Tôi không quay sang, nhưng tôi biết con bé đang mỉm cười, nụ cười của một đứa em gái hiểu chị hơn bất cứ ai trên đời.

Shino nhìn tôi với ánh mắt thấu hiểu, không hỏi thêm. Cô ấy chỉ khẽ nói: ``Đi tiếp thôi, Kaigou-san. Còn nhiều chỗ vui lắm.'' Tôi gật đầu, đứng dậy, hít một hơi thật sâu rồi phồng má hô lớn: ``Được rồi! Đi thôi! Lần này đổi đích, tôi muốn xem Câu lạc bộ Khoa học có gì hay!''

\ct

Khác với Kuon-kun đang ngập đầu trong vai trò quản gia dập khuôn, hay Ayame và Wata đang tất bật với bộ đàm của Hội Học sinh, chúng tôi đang tận hưởng những phút giây nhàn rỗi quý báu nhất của tuổi thanh xuân.

Điểm dừng đầu tiên sau gian bói toán là phòng triển lãm ảnh của Câu lạc bộ Nhiếp ảnh, nằm ở tầng hai dãy nhà A. Ngay từ cửa vào, những bức ảnh lớn được đóng khung gỗ sáng màu đã khiến cả bốn chúng tôi dừng bước.

``Aogami qua bốn mùa,'' Shino đọc tên triển lãm, giọng thoáng ngạc nhiên. ``Họ chụp đẹp thật\dots''

Bộ ảnh được trưng bày theo trình tự mùa: mùa xuân với hàng anh đào nở rợp sân trường, mùa hè với ánh nắng chiếu xuyên qua những tán lá xanh biếc, mùa thu với lá phong đỏ rực phủ kín lối đi, và mùa đông với lớp tuyết trắng xóa phủ lên mái ngói cũ kỹ của trường.

Đặc biệt, ở góc cuối phòng, có một khu vực dành riêng cho những bức ảnh chân dung chụp các học sinh trong những khoảnh khắc tự nhiên nhất: cười, ngáp, ngủ gật, hay đang mải mê đọc sách dưới gốc cây.

Naomi kéo tay tôi, chỉ vào một bức ảnh: ``Kaigou-oneechan, nhìn kìa! Có ảnh Suzuki-oneechan nè!''

Tôi nhìn theo hướng mà cô bé chỉ. Đúng thật: một bức ảnh chụp Suzu đang ngồi ở bậc thềm hiên trước lớp học, ánh chiều vàng rực chiếu nghiêng, tay cầm một cuốn sách mở rộng, gió nhẹ lùa mái tóc sang một bên. Biểu cảm trên mặt con bé thanh bình và dịu dàng đến mức tôi phải nín thở.

``A-Ai chụp cái này vậy!?'' Suzu đỏ bừng mặt, luống cuống nhìn quanh.

Từ phía sau quầy, Uzuki - cũng là cô nàng bờm Thỏ của lớp tôi - khẽ giơ tay lên: ``Mình chụp những bức đó đấy. Ánh sáng lúc đó rất đẹp. Các cậu thấy thế nào?'' Tôi không còn lạ gì cô bạn tóc xanh-trắng kia nữa, khi kỹ năng chụp ảnh của cậu ấy phải nói là thượng thừa; thậm chí còn đoạt được giải cao trong các cuộc thi nhiếp ảnh.

Akane đứng bên cạnh thì cười toe toét, giơ máy ảnh lên: ``Bức đó là một trong những bức đẹp nhất triển lãm đấy, Suzuki-san! Uzuki-chan chọn góc chuẩn lắm!''

Suzu ấp úng, không biết nên vui hay nên xấu hổ. Tôi bật cười, vỗ lưng con bé: ``Đẹp mà ngại gì. Lần sau để Onee-sama mình canh góc cho cậu chụp thêm mấy kiểu nữa, Uzuki-san!''

Uzuki khẽ gật đầu: ``Được thôi.'' Cô bạn bờm Chó lại cười khúc khích.

Rời khỏi phòng triển lãm, chúng tôi xuống sân trường, nơi Câu lạc bộ Bóng rổ đang tổ chức một gian hàng ném bóng vào rổ mini. Sân được chia thành ba mức độ: dễ, trung bình và khó, tùy theo khoảng cách ném.

Và ở đó, tôi nhìn thấy một bóng dáng quen thuộc đang đứng phía sau quầy phát thưởng - Shiina. Cô nàng khoác chiếc áo lớp 1-C bên ngoài bộ đồng phục thể thao, tóc buộc cao, toát lên vẻ lạnh lùng thường ngày. Cạnh cô ấy là Yae, đang hướng dẫn một nhóm học sinh lớp dưới cách ném bóng đúng tư thế.

``Ồ! Kaigou-chan! Suzuki-chan!'' cô bạn Thể lực vẫy tay gọi to khi thấy chúng tôi. ``Mọi người tới chơi đi! Ném ba quả trúng hai là được thưởng đồ uống miễn phí này!''

Cô bạn bờm Sư tử nhìn tôi, nhếch môi: ``Hay là thử sức luôn đi, Kaigou? Mức khó ấy?''

Tôi chống hông, cười khẩy: ``Tưởng gì, dễ ợt.''

Không ngờ tôi đã đánh giá thấp mức khó của CLB Bóng rổ. Khoảng cách ném xa hơn dự tính, và cái rổ bé tí tẹo chỉ vừa đủ lọt quả bóng. Cú ném đầu tiên trượt, cú thứ hai đập vành rổ rồi bật ra. Tôi nghiến răng, tập trung hết sức cho cú cuối cùng.

*PHỊCH!*

Trúng! Quả bóng xuyên gọn qua rổ, không chạm vành. Shiina huýt sáo: ``Không tệ đấy.''

``Mới trúng một thôi mà,'' Yae cười lớn, ``nhưng cũng đáng khen rồi. Để cậu ấy phát cho cậu một lon trà xanh nhé!''

Tôi nhận lon trà với vẻ hơi bất mãn, đáng ra phải trúng cả ba mới đúng bản lĩnh của tôi chứ. Shino đứng bên cạnh, cười khúc khích với dáng vẻ bực dọc trong bộ đồ hầu gái của tôi: ``Kaigou-san ném giỏi thật đấy. Mình thì chắc không ném nổi quả nào đâu.''

``Cậu phải thử mới biết!'' tôi đẩy quả bóng vào tay cô ấy.

Cô bạn tóc nâu ngập ngừng, cầm quả bóng lóng ngóng, rồi ném. Quả bóng bay một đường cong kỳ lạ, chệch khỏi rổ cả một mét. Cô ấy quay sang nhìn tôi, mặt hơi đỏ: ``Mình\dots\ mình nói rồi mà.'' Cả bọn cười ồ lên, nhưng không phải cười chê, mà là cái cười ấm áp của những người bạn đang cùng nhau tận hưởng những khoảnh khắc giản dị nhất.

*Ting tong, ping pong\~{}*

``A lô a lô! Kana đây nè! Các bạn ơi, tin nóng hổi vừa ra lò! CLB Thiên văn ở phòng 2-C tầng hai đang có chương trình quan sát Mặt trời qua kính thiên văn đặc biệt! An toàn tuyệt đối, các bạn yên tâm nhé! Ngoài ra, CLB Cây cảnh đã hoàn thiện góc chụp ảnh `Khu vườn bí mật' ở sân sau dãy nhà D. Rất đẹp, rất thơ, rất đáng xem! Đừng bỏ lỡ nha!''

``CLB Thiên văn!'' Suzu reo lên. ``Saya-san chắc đang ở đó!''

Chúng tôi leo lên tầng hai. Phòng 2-C được bài trí như một phòng quan sát vũ trụ thu nhỏ: trần nhà được dán giấy đen và rắc đầy những ngôi sao dạ quang, một mô hình hệ mặt trời treo lơ lửng ở chính giữa với các hành tinh quay chầm chậm nhờ dây cước và động cơ nhỏ. Ở ban công, một chiếc kính thiên văn được gắn bộ lọc đặc biệt hướng thẳng lên bầu trời.

Saya đang đứng cạnh chiếc kính, dịu dàng hướng dẫn một nhóm học sinh cách điều chỉnh ống kính. Khi thấy chúng tôi, cô ấy mỉm cười ấm áp: ``Các cậu tới rồi sao? Tuyệt quá\~{} Hôm nay trời trong lắm, nhìn rõ cả vết đen Mặt trời luôn đấy\~{}''

Naomi háo hức chen vào hàng đầu tiên, nhón chân nhìn vào ống kính. ``Oa! Em thấy Mặt trời rồi kìa! Nó to lắm! Mà sao có mấy cái chấm đen vậy, Saya-oneechan?''

Cô nàng tóc vàng cúi xuống, giải thích cho em bé bằng giọng chậm rãi, nhẹ nhàng thường ngày: ``Đó là vết đen Mặt trời đó, Naomi-chan. Nó giống như vết tàn nhang trên khuôn mặt Mặt trời vậy. Có vẻ dữ dội lắm, nhưng thực ra nó rất bình thường và vô hại đó\~{}''

``Vết tàn nhang!'' Naomi reo lên thích thú, ``Vậy Mặt trời cũng có tàn nhang giống mấy bạn học sinh ạ?''

Cả phòng phì cười trước sự hồn nhiên của cô bé. Tôi nhìn sang Shino, thấy cô ấy đang ngước lên nhìn mô hình hệ mặt trời, ánh mắt xa xăm nhưng bình yên.

``Cậu đang nghĩ gì thế?'' tôi khẽ hỏi.

``Mình đang nghĩ\dots\ vũ trụ thật rộng lớn,'' cô ấy đáp, giọng nhẹ bẫng. ``Nhưng ở góc nhỏ bé này, trên mảnh đất Aogami này, chúng ta đang được sống một cuộc đời trọn vẹn. Thế là quá đủ rồi.''

Tôi không đáp, chỉ gật đầu. Có những lúc, Shino nói ra những điều mà cả tôi cũng không biết diễn tả thành lời.

Sau CLB Thiên văn, chúng tôi kéo nhau xuống sân sau dãy nhà D, nơi CLB Cây cảnh đã biến một góc sân bê tông thành ``Khu vườn bí mật,'' một không gian xanh mát với những chậu hoa đủ sắc màu, dây leo phủ kín khung cổng gỗ giả cổ, và những chiếc ghế trắng nhỏ xinh được kê dưới giàn hoa hồng leo. Saki đang tất bật chỉnh sửa mấy chậu hoa ở lối vào, trong khi một vài học sinh lớp khác đang háo hức xếp hàng chờ chụp ảnh.

``Chào mọi người!'' Saki vẫy tay, nụ cười hiền hòa nở rộng khi thấy chúng tôi. ``Vào chụp ảnh đi! Hoa tươi mới cắm sáng nay, đẹp lắm luôn!''

Naomi lập tức lao vào khu vườn như một chú bướm nhỏ, xoay tròn giữa những chậu hoa cúc vàng và hoa oải hương tím. Suzu nhanh tay rút điện thoại ra, chụp liên tục. Shino đứng bên khung cổng hoa, tay chạm nhẹ vào những cánh hồng, nở nụ cười dịu dàng nhất mà tôi từng thấy ở cô ấy.

``Đứng yên nhé, Shino-san, để tôi chụp cho!'' tôi giơ điện thoại lên. Cô ấy trông hơi ngượng, nhưng vẫn đứng yên, ánh nắng chiều hắt qua giàn hoa tạo nên một khung hình đẹp đến nao lòng.

*TÁCH*

Tôi nhìn lại bức ảnh vừa chụp. Shino trong ảnh trông thật thanh tao, tĩnh lặng, nhưng cũng tràn đầy sức sống. Tôi lưu lại, tự hứa sẽ gửi cho cô ấy sau.

Rời CLB Cây cảnh, chúng tôi dừng chân trước phòng sinh hoạt của CLB Thời trang ở tầng một dãy nhà B. Bên trong, một buổi trình diễn thời trang mini đang diễn ra trên sàn catwalk được ghép tạm từ những tấm ván ép phủ vải đen bóng.

Mitsuki đang đứng bên cánh gà, tay cầm bảng ghi chú, chỉ đạo trình tự ra sân khấu. Khi thấy chúng tôi, cô ấy nháy mắt: ``Các cậu đến đúng lúc quá! Chương trình sắp bắt đầu rồi!''

``Ồ! Mitsuki-san kìa!'' Naomi chạy ào tới. ``Em nhớ chị lắm!''

Mitsuki cúi xuống xoa đầu cô bé: ``Chị cũng nhớ Nao-chan. Em mặc bộ váy Công chúa Hoa xinh quá na! À mà\dots'' cô ấy liếc sang tôi với nụ cười bí ẩn, ``chương trình hôm nay có một phần đặc biệt cho khán giả nhí đấy nhé.''

Trước khi tôi kịp hỏi ``phần đặc biệt'' là gì, tiếng nhạc đã nổi lên. Lần lượt các ``người mẫu'' bước ra học sinh từ nhiều lớp khác nhau, mặc những bộ trang phục do các thành viên CLB Thời trang tự thiết kế: từ phong cách thường nhật ngày hè đến váy dạ hội tái chế từ vải vụn. Tất cả đều toát lên sự sáng tạo và tâm huyết.

Hotaru ngồi ở hàng ghế phía trước, tay cầm bảng chấm điểm, mắt sáng rực khi nhìn từng bộ trang phục lướt qua. Cô nàng quay lại nhìn chúng tôi, vẫy tay hớn hở: ``Kaigou-san! Suzuki-san! Mấy bộ này mình cũng có góp tay vẽ mẫu đó nha!''

``Cậu giỏi quá, Hotaru-san!'' Suzu tấm tắc khen.

Rồi đến phần ``đặc biệt.'' Mitsuki bước lên sân khấu, micro trong tay, giọng trang trọng nhưng ấm áp: ``Và bây giờ, chúng tôi có một phần trình diễn đặc biệt! Xin mời vị khách nhí đặc biệt nhất hôm nay\dots\ \textbf{Naomi-chan!}''

Naomi trố mắt, chỉ tay vào mình: ``E-Em ạ!?''

Cô nàng tóc hồng-tím gật đầu, mỉm cười: ``Đúng rồi na! Chị đã chuẩn bị một bộ trang phục riêng cho em. Em thử lên sân khấu trình diễn nhé?''

Cô bé ngước lên nhìn tôi và Suzu, ánh mắt vừa ngại ngùng vừa háo hức. Tôi cười, đẩy nhẹ lưng em: ``Lên đi, Nao. Cho mọi người thấy công chúa của nhà Hikawa đáng yêu cỡ nào!''

Vài phút sau, Naomi bước ra trong bộ váy trắng phối ren hoa, đầu đội vòng hoa dại, tay cầm gậy phép thuật giả lấp lánh. Bước đầu rụt rè, nhưng khi nghe tiếng vỗ tay, cô bé dần tự tin, xoay một vòng, váy tung ra như cánh hoa nở.

``M-Mọi người thấy em thế nào?'' cô bé hỏi, giọng hơi run nhưng vẫn đủ lớn để mọi người nghe thấy.

Không cần một câu trả lời nào, cả phòng vỗ tay rào rào như thay cho những lời muốn nói. Suzu rơm rớm nước mắt vì xúc động, còn tôi thì mím môi cười, tự hào không diễn tả nổi.

Mitsuki quỳ xuống ôm Naomi khi cô bé bước về: ``Em làm tốt lắm, Nao-chan. Em là ngôi sao sáng nhất hôm nay đó\~{}''

Naomi nhảy chân sáo xuống từ sân khấu, mặt đỏ bừng vì phấn khích nhưng đôi mắt lấp lánh niềm vui. Cô bé lao vào lòng tôi, cười rạng rỡ: ``Kaigou-oneechan! Nao làm tốt không?''

Tôi xoa đầu cô bé, không giấu nổi vẻ tự hào: ``\vie{Tuyệt lắm, Công chúa của chúng ta là nhất rồi.}''

Shino đứng bên cạnh, ánh mắt cũng đầy trìu mến. Cô ấy khẽ vuốt lại mái tóc cho Naomi, rồi quay sang tôi và Suzu: ``\vie{Cảm ơn hai cậu. Nếu không đi cùng mọi người, có lẽ mình đã bỏ lỡ những khoảnh khắc đẹp đẽ này. Đây thực sự là lần đầu tiên mình cảm thấy mình được sống như một nữ sinh bình thường, thay vì là một người nắm giữ vận mệnh mệt mỏi.}''

Chúng tôi bắt đầu rảo bước ra khỏi dãy nhà B, cùng nhau đi qua những gian hàng cuối cùng của các lớp khối dưới. Bầu không khí vẫn náo nhiệt, nhưng nhịp độ bắt đầu chậm lại, chuẩn bị cho sự kiện lớn nhất cuối ngày.

\ct

*Ting tong, ping pong\~{}*

Giọng Kaoru vang lên trong loa phát thanh: ``Xin thông báo, chương trình biểu diễn văn nghệ trên sân khấu chính sẽ bắt đầu trong ba mươi phút nữa. Xin mời tất cả các bạn học sinh và quý khách di chuyển về phía sân trường để thưởng thức các tiết mục đặc sắc từ Câu lạc bộ Âm nhạc và nhiều nhóm biểu diễn khác. Xin nhắc lại\dots''

Tôi nhìn đồng hồ, lúc này đã chỉ về ba rưỡi chiều. Ca làm việc cũng sắp hết rồi. Tôi quay sang cả bọn: ``Sân khấu chính kìa. Mình ghé qua xem trước khi quay lại quán nhé?''

Shino gật đầu, nhưng trước khi đi, cô ấy dừng lại, quay sang nhìn tôi, Suzu và Naomi. Ánh mắt cô ấy dịu dàng và chân thành: ``Cảm ơn ba người. Hôm nay\dots\ là ngày vui nhất mà mình từng có kể từ khi được tái sinh.''

Tôi mỉm cười, vỗ vai cô ấy: ``Vẫn còn ngày mai đang chờ cậu phía trước, Shino-san. Tôi hứa đấy.''

Chúng tôi chia tay Shino tại sảnh chính để cô ấy quay về khu vực tập trung của Hội Học sinh. Tôi, Suzu và Nao rảo bước quay lại phòng 1-C.

Khi bước vào lớp, đập vào mắt tôi là cảnh tượng quán đã bắt đầu thưa khách. Ánh nắng chiều tà xuyên qua cửa sổ, hắt lên những chiếc bàn gỗ đã được thu dọn phân nửa. Kuon-kun lúc này vẫn đang trong bộ đồ quản gia, tay cầm chiếc khay, điềm đạm cúi chào hai vị khách cuối cùng vừa đứng dậy. Dù mái tóc hơi rối và mồ hôi lấm tấm trên trán, trông cậu ta vẫn giữ được cái khí chất chuyên nghiệp đến kỳ lạ.

Ở góc bếp, Touka và Kanade đang bận rộn rửa nốt những chiếc tách trà cuối cùng. Tôi nhìn quanh nhưng không thấy Honoka đâu.

``\vie{Honoka-san đâu rồi?}'' tôi hỏi khi bước tới quầy.

Kuon đặt khay xuống, thở phào một hơi: ``\vie{Cậu ấy được ưu tiên nghỉ sớm ba mươi phút để sang sân khấu chính chuẩn bị. CLB Âm nhạc có tiết mục mở màn mà.}'' Cậu ta lau mồ hôi, rồi nhìn chúng tôi, ánh mắt dịu lại khi thấy Nao vẫn đang cười toe toét. ``\vie{Chơi vui chứ?}''

``\vie{Trên cả vui ạ!}'' Suzu đáp. ``\vie{À, Onii-sama này, bọn em cũng mang về chút đồ ăn cho anh đây.}''

Con bé rút từ trong túi giấy ra một hộp takoyaki còn ấm và một phần yakisoba gói cẩn thận, đặt lên bàn trước mặt Kuon. Tôi cũng đặt thêm lon trà lúa mạch lạnh xuống bên cạnh, là phần thưởng mà tôi giành được ở gian hàng bóng rổ.

Kuon nhìn đống đồ ăn, chớp mắt vài giây, rồi một nụ cười nhỏ hiện lên trên môi cậu ta, cái kiểu cười hiếm hoi mà cậu ta chỉ dành cho gia đình: ``\vie{Cảm ơn hai đứa. Đúng lúc tôi cần.}''

``\vie{`Đúng lúc' là sao?}'' tôi nhíu mày, liếc nhìn cậu ta. ``\vie{Cậu chưa ăn gì từ sáng tới giờ à?}''

Kuon mở hộp takoyaki, gắp một viên, biểu cảm bình thản như thể việc nhịn ăn tám, chín tiếng đồng hồ là chuyện cơm bữa: ``\vie{Có gì đâu. Hồi còn làm tổng quản, tôi hay bỏ bữa sáng lẫn bữa trưa mà.}''

``\vie{Anh nói cái gì!?}'' Suzu trợn mắt. ``\vie{Vậy là từ sáng tới giờ anh chỉ uống nước suông thôi sao?}''

``\vie{Cũng không hẳn. Touka-san có pha cho tôi một cốc trà xanh lúc mười giờ,}'' cậu ta đáp tỉnh bơ.

Tôi chống hông, lắc đầu ngao ngán: ``\vie{Trà xanh. Cậu sống bằng trà xanh. Tuyệt thật.}'' Nhưng dù miệng nói cứng, trong lòng tôi lại thấy se thắt. Cái thói quen bỏ bữa đó không phải là sự lười biếng, nó là dấu vết của một đứa trẻ đã quen với việc nhường phần ăn cho người khác và tự gánh mọi thứ lên vai mình từ quá lâu. Giống tôi. Giống Suzu. Chúng tôi đều mang những vết thương tương tự, chỉ khác nhau ở cách biểu hiện.

``\vie{Ăn hết đi. Không được để thừa,}'' tôi ra lệnh, giọng cứng nhưng ánh mắt thì dịu. Kuon khẽ gật đầu, không phản bác.

Sau khi hai vị khách cuối cùng rời đi, quán chính thức đóng cửa. Cả lớp 1-C quây quần lại dọn dẹp, mỗi người một tay. Tôi và Suzu xắn tay áo, bắt đầu khiêng bàn ghế trả về vị trí cũ.

``Suzune-san! Cậu đang ăn vụng cái gì đấy!?'' tiếng Touka vang lên sắc gọn từ phía quầy bếp.

Suzune đứng đơ cứng, miệng còn ngậm nguyên nửa miếng bánh pancake mặt mèo, mắt tròn xoe. ``M-Mình không có ăn vụng! Cái này là\dots\ kiểm tra chất lượng!''

``Kiểm tra chất lượng bằng cách nhét hết vào bụng à?'' Touka khoanh tay, nheo mắt. Thế rồi, cô thở dài, miễn cưỡng khoát tay: ``Dù gì thì cũng hết khách rồi, ăn đi.''

``Ehe\dots\ xin lỗi mừ\~{}'' cô nàng ``Linh vạt'' le lưỡi, nhưng cái miệng vẫn nhóp nhép không ngừng. Touka thở dài bất lực, nhưng khóe môi cô ấy khẽ nhếch lên; một biểu hiện cực kỳ hiếm của sự khoan dung.

Inari đang lau bàn ở góc xa, khẽ mỉm cười khi nhìn cảnh tượng đó, rồi quay sang Hanabi bên cạnh: ``Hanabi-chan, cậu giúp mình khiêng cái bảng hiệu kia xuống được không?''

``Để Hanabi lo!'' cô nàng tóc đuôi ngựa bên xắn tay, tự tin nắm lấy một bên tấm bảng. ``Ừm\dots\ ê, nó nặng hơn tôi tưởng! Inari, đỡ phía kia với!''

Miku đang gấp lại từng chiếc tạp dề đã giặt sạch, xếp ngay ngắn thành từng chồng. Cô ấy nhìn đống tạp dề, rồi khẽ thở dài với nụ cười nhẹ: ``Vậy là hết rồi nhỉ\dots\ Nhanh thật.''

Kanade từ trong bếp bưng ra một chồng đĩa đã rửa xong, đặt nhẹ nhàng lên kệ: ``Miku-san, cậu có cần giúp gấp nốt mấy cái khăn trải bàn không?''

``Ừ, cảm ơn Kanade-san.'' Hai người lặng lẽ làm việc bên nhau, nhịp nhàng và ăn ý.

Căn phòng từ một quán café sang trọng dần trở lại thành một lớp học bình thường. Nhưng mùi thơm của trà hoa và bánh ngọt vẫn còn vương lại trong không khí, cùng với những ký ức ấm áp mà chúng tôi vừa tạo ra.

\ct

Khi mọi thứ đã được dọn dẹp gọn gàng, cả lớp kéo nhau ra sân trường. Lúc này đã gần năm giờ chiều, ánh nắng bắt đầu dịu đi, nhuộm một lớp vàng cam lên mọi thứ. Các nhóm từ khắp nơi bắt đầu tụ tập, và ở đó, tôi nhìn thấy nhóm nhà ma của Shion cũng đang tiến tới từ phía dãy nhà đối diện.

Shion dẫn đầu, bước đi với phong thái của một nữ vương vừa kết thúc ca trị vì. Phía sau cô ấy là Aku vẫn còn vương vài vệt sơn trắng nhợt trên cổ từ lớp hóa trang hồn ma, trông vừa buồn cười vừa rợn người. Yuka lững thững đi phía sau cùng, đôi mắt lim dim như thể cô nàng có thể ngủ gục bất cứ lúc nào.

``Ồ, nhóm nhà ma đây rồi!'' Suzune vẫy tay hào hứng.

Cô nàng Phù thủy khẽ hất mái tóc ra sau, khoanh tay, ánh mắt lướt qua tất cả chúng tôi với vẻ tự mãn đặc trưng: ``Hừ, cuối cùng các ngươi cũng tề tựu đông đủ. Ta có một bản báo cáo chiến tích cho `Con đường Linh hồn' đây.''

``Chiến tích?'' Kuon nhướng mày.

Cô rút từ trong túi áo ra một mẩu giấy, đọc to với giọng trang trọng như đang công bố kết quả thi: ``Tổng số lữ khách: một trăm hai mươi ba người. Số người hoàn thành trọn vẹn: bốn mươi bảy. Số người bỏ cuộc giữa chừng: bảy mươi sáu.''

``Hơn phân nửa bỏ cuộc!?'' Suzu trợn mắt.

``Chưa hết,'' Shion tiếp tục, giọng không giấu nổi sự khoái chí. ``Số người ngất xỉu: ba. Số người\dots'' cô ấy hắng giọng, ``\uline{dấm đài}: hai.''

Một khoảng im lặng ngắn ngủi. Rồi cả đám bật cười ầm ĩ.

``\textbf{Hai người dấm đài ư!?}'' Hanabi ôm bụng lăn ra cười. ``Shion-chan ghê gớm quá vậy!''

``Chết tiệt, vậy là có người sợ tới mức\dots'' tôi bịt miệng, cố không bật cười to hơn, nhưng thất bại thảm hại.

``Đúng như cậu ấy nói luôn đó\~{}'' Yuka ngân giọng, ``lúc bọn tôi bắt đầu dọn dẹp lại thì thấy có vũng nước ở đó mà\~{}''

Aku gãi đầu, vẻ mặt ngượng ngùng: ``Thật ra\dots\ một trong hai người đó ngất ngay trước mặt mình. Mình không biết phải làm gì nên chỉ đứng yên\dots\ rồi người ta tỉnh dậy, thấy mình vẫn đứng đó, lại ngất tiếp\dots''

``Cậu đứng yên mới là đáng sợ nhất đấy, Aku-san,'' Kuon lắc đầu, nửa cười nửa thở dài.

Game cũng lên tiếng từ chiếc đồng hồ trên cổ tay Kuon, giọng pha chút đắc ý: ``Tôi chỉ muốn bổ sung rằng hệ thống âm thanh vòm mà nhómm Azuki thiết kế đã hoạt động với hiệu suất 97,3\%. Đáng tiếc là 2,7\& phần trăm còn lại bị lỗi vì một ai đó tôi không nói tên đã vô tình giẫm lên dây loa.''

Cô nàng tóc xanh-tím rụt cổ: ``X-Xin lỗi\dots''

``Nhưng mà nói thật nha,'' Miku chống hông, nhìn Shion bằng ánh mắt vừa nể phục vừa kinh hãi, ``tôi nghe các chị lớp 2-A kể lại mà nổi hết da gà. Họ bảo sau khi đi xong thì ai cũng nhìn nhau im re, không dám nói gì cả. Shion-san, cậu có đang triệu hồi ma thật không đấy?''

Cô nàng phù thủy khẽ cười bí hiểm: ``Ta chỉ\dots\ tạo điều kiện cho chúng ghé thăm thôi.''

``Nghe mà muốn tè ra quần,'' Suzune rùng mình.

Tôi quay sang nhìn Kuon. Cậu ta đang lắng nghe mọi người, nhưng ánh mắt thoáng hướng về phía Naomi. Cô bé lúc này đang kéo ống tay áo Kuon, giọng hào hứng: ``Onii-chan! Onii-chan! Để em kể cho anh nghe này! Bọn em đi nhà ma của Shion-oneechan, sợ lắm luôn! Có mấy cái bóng ma nhảy ra, rồi Kaigou-oneechan la hét, rồi Suzuki-oneechan nắm tay em chặt lắm! Nhưng mà cuối cùng em không khóc đâu nhé!''

``Ê, tôi không có la hét!'' tôi phản đối ngay lập tức, mặt hơi nóng.

Kuon nhìn xuống Naomi, ánh mắt dịu dàng: ``Vậy à? Em giỏi lắm. Nhưng mà có đúng là Kaigou-chan la hét thật không?'' Cậu ta liếc sang tôi với nụ cười mỉm khó chịu.

``Im đi!'' tôi gầm gừ.

Suzu che miệng cười: ``Onee-sama có\dots\ hét một tiếng nhỏ thôi ạ.''

``Suzu!''

Cả bọn lại cười ồ lên. Kuon xoa đầu Naomi, giọng nhẹ nhàng: ``Cảm ơn đã kể cho anh nghe. Lần sau nếu có đi nhà ma nữa, anh sẽ đi cùng nhé.''

Cô bé reo lên, gật đầu lia lịa.

Giữa lúc tiếng cười còn chưa dứt, điện thoại của cậu ta rung lên. Cậu ta liếc màn hình, rồi đưa cho tôi xem. Một tin nhắn từ Honoka:

\txtbx{
  ``Sắp tới lượt bọn mình rồi! Mọi người ra sân khấu cổ vũ đi nhé!! \textbackslash(\textasciicircum\ O \textasciicircum)/''}

``Đi thôi, mọi người đang chờ rồi,'' cậu ta đứng dậy. ``Đám nhà ma dọn xong chưa, Shion-san?''

Shion phẩy tay: ``Dĩ nhiên. Ta đã thu hồi mọi bùa chú và hiệu ứng tâm linh rồi. Căn phòng giờ chỉ là phòng học trống thôi.''

``Thế thì ai dọn mấy cái búp bê và đạo cụ?'' Aku rụt rè hỏi.

``Ngươi chứ ai,'' Shion đáp tỉnh bơ. Cô bạn thân rũ người, thở dài não nề. ``Tại sao lúc nào cũng là tôi chứ\dots''

\ct

Sân khấu chính lúc này rực rỡ dưới ánh đèn pha. Hàng trăm học sinh và khách tham quan đã tập trung thành một biển người, tiếng ồn ào xen lẫn tiếng nhạc nền tạo thành một bản giao hưởng hỗn loạn nhưng đầy sức sống. Cả lớp 1-C tìm được một khoảng trống gần phía trước, quây quần lại với nhau.

Các tiết mục văn nghệ lần lượt diễn ra: một màn nhảy hiện đại sôi động của khối mười mở đầu, tiếp theo là vở kịch ngắn của CLB Văn học. Một câu chuyện cảm động về tình bạn khiến một vài người xung quanh rơm rớm nước mắt. Rồi đến tiết mục hợp xướng của CLB Giao hưởng khối năm ba, giọng ca hòa quyện du dương.

``\vie{Trường mình nhiều tài năng thật đấy,}'' Suzu thì thầm, mắt sáng long lanh.

``\vie{Ừ,}'' tôi gật đầu, ``\vie{nhưng tiết mục hay nhất chưa tới.}''

Kana chạy lên sân khấu, micro trong tay, giọng đầy năng lượng: ``Và bây giờ! Tiết mục mà mọi người chờ đợi nhất! Xin chào đón ban nhạc đến từ Câu lạc bộ Âm nhạc với ca khúc gốc mang tên `Ánh sao cuối chân trời'!''

Ánh đèn vụt tắt. Sân trường chìm vào bóng tối trong vài giây. Rồi một luồng ánh sáng xanh-tím chiếu thẳng vào trung tâm sân khấu.

Azuki đứng giữa, micro trong tay, phong thái của một ngôi sao thực thụ. Phía sau bên trái, Honoka cầm cây guitar, biểu cảm tập trung và rạng rỡ, hoàn toàn khác với cô nàng nhí nhảnh hậu đậu thường ngày. Bên phải, Yuki ngồi sau bàn keyboard, ngón tay đã sẵn sàng trên phím.

Cô bạn tóc đen-đỏ nhấn phím mở đầu. Một giai điệu piano nhẹ nhàng, mong manh, lan tỏa qua hệ thống loa. Rồi, cô bạn Cáo sa mạc bắt đầu đệm guitar, những nốt nhạc chậm rãi, ấm áp, như tiếng thì thầm của hoàng hôn. Và khi giọng hát của cô nàng Phiêu lưu cất lên, cả sân trường lặng phắc.

Giọng cô ấy cao vút, đầy nội lực nhưng cũng mang một nỗi buồn sâu thẳm. Ca từ kể về một ngôi sao nhỏ ở cuối chân trời, biết rằng bình minh sẽ đến và mình sẽ phải nhường chỗ cho ánh mặt trời, nhưng vẫn tỏa sáng hết mình trong khoảnh khắc cuối cùng của màn đêm.

``\vie{Bài hát này\dots}'' Suzu khẽ nắm tay tôi, giọng run run. ``\vie{Như đang nói về chúng ta vậy.}''

Tôi không đáp, nhưng tôi hiểu. Yuki đã viết bài hát này cho lớp 1-C, và dù có lẽ cô ấy không biết rõ chúng tôi sẽ rời đi, nhưng trực giác của một người nghệ sĩ đã cho cô ấy cảm nhận được điều gì đó. Giai điệu chuyển sang đoạn cao trào: Azuki hát vỡ òa, guitar điện của Honoka gầm lên mạnh mẽ, keyboard của Yuki dồn dập, rồi tất cả dừng lại đột ngột, chỉ còn giọng hát a cappella của Azuki vang vọng trong không gian tĩnh lặng, trước khi nốt cuối cùng tan dần vào gió chiều.

Im lặng. Một giây. Hai giây.

Rồi cả sân trường vỡ òa. Tiếng vỗ tay, tiếng hò reo, tiếng huýt sáo vang dội. Honoka trên sân khấu giơ guitar lên chào, nụ cười rạng rỡ. Azuki cúi đầu cảm ơn, mắt hơi đỏ. Yuki vẫy tay nhẹ, nụ cười dịu dàng.

``Honoka-san ngầu quá!'' Naomi nhảy cẫng lên.

``Bọn họ\dots\ tuyệt vời thật,'' Kuon khẽ nói, giọng đơn sắc như mọi khi, nhưng tôi thấy cậu ta đang vỗ tay rất chậm, rất chắc, như thể muốn khắc ghi từng nhịp vào ký ức.

Tôi đứng giữa đám đông, nhìn Honoka trên sân khấu. Cô nàng thường ngày hay nhí nhảnh, hay bày trò và lắm chuyện, giờ tỏa sáng đến mức không thể nhận ra. Rồi tôi nhìn sang Kuon và Suzu. Đây là lễ hội trường học đầu tiên của chúng tôi ở thế giới này, và có lẽ\dots\ cũng là cuối cùng. Một nỗi buồn man mác dấy lên trong lòng, nhưng tôi không muốn để nó lấn át khoảnh khắc này. Chúng tôi đã chiến đấu rất nhiều, đã trải qua bao vòng lặp chỉ để có được một ngày yên bình như thế này. Aogami này, ngôi trường này, những người bạn này\dots\ mọi thứ đều quá thật, đến mức việc phải rời đi khiến lồng ngực tôi thắt lại.

Chương trình văn nghệ kết thúc trong tiếng vỗ tay rền trời. Hội trưởng Ayame bước lên sân khấu, bộ đàm vẫn còn dắt ở hông. Cô nàng cầm micro, nụ cười tinh quái thường ngày giờ mang vẻ tự hào:

``Cảm ơn tất cả mọi người vì một ngày hội tuyệt vời! Nhưng hãy nhớ rằng, cuộc vui chỉ mới bắt đầu thôi! Ngày mai, Lễ hội thị trấn Natsu Matsuri sẽ chính thức diễn ra tại đền Aogami!'' Cô nàng giơ ngón tay lên, đếm từng mục: ``Sẽ có Yatai, đồ ăn ngon ơi là ngon! Sẽ có trò chơi truyền thống như vớt cá, bắn giải, câu bóng nước! Sẽ có Mikoshi, ai muốn khiêng kiệu thần thì đăng ký với Wata nhé! Sẽ có Bon Odori, cả trường cùng nhảy dưới trăng! Và quan trọng nhất\dots'' Cô nàng dừng lại, nháy mắt: ``\textbf{Màn pháo hoa hoành tráng nhất từ trước đến nay!} Được tài trợ bởi nhà Furukawa! Đừng có ai ngủ quên mà bỏ lỡ đấy nhé!''

Wata đứng phía dưới, giơ tay bổ sung với giọng điềm đạm: ``Bắt đầu từ tám giờ sáng. Nhớ mặc yukata đấy.''

Tiếng reo hò vang lên lần nữa. Tôi có thể cảm thấy sự nóng lòng của mọi người xung quanh.

Nhưng đối với bộ ba Hikawa chúng tôi, sự nóng lòng đó còn mang một ý nghĩa nặng nề hơn. Khi đám đông bắt đầu giải tán, Kuon khẽ kéo tôi và Suzu ra một góc yên tĩnh, giọng hạ thấp:

``\vie{Ngày mai, sau khi lễ hội thị trấn kết thúc, chúng ta sẽ tiến hành Nghi thức Thanh tẩy cho Naomi. Shion-san và Sakura-san đã xác nhận: trăng rằm tháng này là thời điểm thích hợp nhất.}''

Suzu siết chặt tay tôi. ``\vie{Vậy là\dots\ sau đêm mai\dots}''

``\vie{\dots Naomi sẽ được tự do,}'' cậu đáp, giọng vẫn bình thản, nhưng tôi thấy cậu ta nắm chặt tay lại. ``\vie{Và chúng ta cũng sẽ biết được\dots\ khi nào mình có thể trở về nhà.}''

Tôi nhìn Naomi đang chạy nhảy cùng Suzune ở phía xa, vô tư và rạng rỡ. ``\vie{Vậy thì ngày mai\dots\ chúng ta phải vui hết mình. Không được lãng phí một giây nào.}''

Cả hai người cùng gật đầu, mắt hơi đỏ nhưng nụ cười vẫn kiên định.

Tôi nhìn bầu trời Aogami đang chuyển sang sắc tím hoàng hôn. Ngày mai sẽ là ngày quan trọng nhất kể từ khi chúng tôi đặt chân tới thế giới này. Nhưng hôm nay, ít nhất là hôm nay, chúng tôi đã có một lễ hội trường học trọn vẹn, thứ mà cả ba chị em từng nghĩ chỉ là giấc mơ xa xỉ.

Và giấc mơ đó, dù chỉ kéo dài một ngày, cũng đã đủ để khắc sâu vào trái tim chúng tôi - những người được quay trở lại thời học sinh - quãng thời gian tuổi trẻ nhiệt huyết nhất.

\begin{center}
  \uline{Ngày hôm sau.}
  \pov{Hikawa Suzuki}
\end{center}

Ánh nắng buổi sáng của ngày thứ hai trong chuỗi lễ hội đổ xuống thị trấn Aogami với một sắc độ vàng ươm, trong trẻo như mật ong. Tôi thức dậy với cảm giác lồng ngực mình vẫn còn đập thình thịch những dư âm của đêm qua: tiếng nhạc của Azuki-san, ánh đèn sân khấu và cả lời nhắc nhở trầm mặc của Onii-sama về ``ngày cuối cùng''.

Nhưng khi mở cửa sổ nhìn xuống con phố dẫn về phía đền Aogami, mọi lo âu trong tôi dường như bị cuốn phăng bởi một luồng sinh khí mãnh liệt.

Thị trấn Aogami sáng nay đã lột xác hoàn toàn. Không còn là một thị trấn ngoại ô yên tĩnh, nó biến thành một mê cung rực rỡ sắc màu. Những dải đèn lồng chōchin màu đỏ và trắng đã được giăng kín các lối đi, đung đưa nhẹ nhàng trong gió sớm như những quả cầu lửa nhỏ. Các sạp hàng san sát nhau, bắt đầu đỏ lửa. Mùi thơm của sốt takoyaki đậm đà, vị ngọt lịm của kẹo táo và cả mùi khói than nướng thịt xiên yakitori hòa quyện vào nhau, tạo thành một thứ hương vị đặc trưng mà chỉ cần hít vào, người ta đã thấy lòng mình rộn ràng.

``\vie{Suzuki-oneechan! Nhanh lên nào! Mọi người đi hết bây giờ!}''

Tiếng gọi lanh lảnh của Naomi vang lên từ phòng khách. Tôi vội vàng chỉnh lại dải đai obi trên bộ yukata màu tím nhạt có họa tiết hoa linh lan của mình, rồi bước ra ngoài.

Onii-sama và Onee-sama đã đợi sẵn ở đó. Anh hôm nay trông lạ lẫm nhưng vô cùng nam tính trong bộ yukata màu đen tuyền giản dị, mái tóc được vuốt lại gọn gàng hơn thường ngày. Còn chị tôi chọn một bộ yukata màu xanh hải quân đậm với họa tiết sóng nước trắng xóa. Dù chị ấy cứ càu nhàu rằng ``mặc cái này khó di chuyển chết đi được'', nhưng tôi thấy Onee-sama thi thoảng lại lén soi gương với vẻ khá hài lòng.

Và nổi bật nhất là Naomi. Cô bé diện bộ yukata màu hồng đào với những bông hoa anh đào nhỏ xíu li ti, hai cái búi tóc được Mitsuki tạo kiểu từ tối qua đung đưa theo mỗi bước chạy.

Chúng tôi cùng nhau bước ra phố.

Dòng người bắt đầu nườm nượp đổ về phía trung tâm thị trấn. Những nhóm học sinh mặc thường phục, những gia đình dắt tay nhau, và cả những bậc cao niên trong những bộ kimono truyền thống sang trọng. Không gian tràn ngập tiếng cười nói, tiếng chuông gió leng keng và tiếng trống taiko thử loa từ phía xa vọng lại.

``\vie{Oa\~{} Chị ơi, nhìn kìa! Có cả trò câu cá vàng nữa!}'' Naomi kéo tay tôi, đôi mắt lục bảo mở to hết cỡ vì kinh ngạc. ``\vie{Lễ hội thị trấn to hơn ở trường mình nhiều chị nhỉ?}''

Tôi mỉm cười, nắm lấy bàn tay nhỏ nhắn của em: ``\vie{Tất nhiên rồi. Đây là hội lớn nhất năm của cả vùng mà. Em thấy thế nào?}''

Cô bé nghiêng đầu suy nghĩ một chút rồi đáp bằng giọng hồn nhiên: ``\vie{Ở trường hôm qua vui kiểu náo nhiệt, giống như một bữa tiệc lớn ấy. Còn ở đây\dots\ em thấy nó giống như trong mấy cuốn sách cổ tích. Mọi thứ cứ lung linh, trang trọng nhưng cũng ấm cúng lắm ạ.}''

Onee-sama nhếch môi, vòng tay qua vai Naomi: ``\vie{Đúng đó. Ở trường là để quậy phá, còn ở đây là để tận hưởng cái không khí truyền thống. Nhìn mấy cái đèn lồng kia đi, tối nay khi thắp sáng lên, cả thị trấn này sẽ trông như một vùng đất thần tiên thực sự đấy.}''

Onii-sama đi bên cạnh, dù vẫn giữ vẻ mặt điềm đạm thường ngày, nhưng tôi nhận ra ánh mắt anh đang quan sát mọi thứ rất kỹ. Không phải là sự quan sát cảnh giác của một tổng quản, mà là sự trân trọng của một người lữ khách đang cố gắng ghi lại mọi hình ảnh vào tâm khảm. Anh khẽ chạm vào một gian hàng bán bùa may mắn, ánh mắt thoáng chút dịu dàng.

Chúng tôi dọc theo con phố chính, đi qua những sạp hàng trò chơi ném vòng, bắn súng hơi và cả những quầy bán mặt nạ cáo, mặt nạ hyottoko đầy màu sắc. Càng tiến gần về phía đền Aogami, không khí càng trở nên trang nghiêm và nhộn nhịp hơn.

Đến cổng torii lớn ở lối vào khu vực đền, chúng tôi đã nhìn thấy những bóng dáng quen thuộc. Lớp 1-C đã tụ tập ở đó đông đủ. Những tà áo yukata đủ màu sắc khiến khu vực cổng torii trông như một vườn hoa rực rỡ.

Suzune đang nhảy cẫng lên vẫy tay gọi chúng tôi, trong khi Touka đứng bên cạnh cố gắng giữ cho cô bạn khỏi va vào người khác.

Phía trước cổng, nhóm Hội Học sinh đang bận rộn điều phối dòng người. Nanase và Inari đeo băng tay đỏ, tay cầm sơ đồ, chỉ dẫn vị trí cho các nhóm lữ khách một cách chuyên nghiệp. Ayame và Wata thì đang trao đổi gì đó với những người dân thị trấn phụ trách kiệu thần Mikoshi.

Thậm chí cả Shino cũng ở đó. Cô ấy mặc bộ yukata màu trắng tinh khôi với họa tiết lá phong xanh, trông thanh thoát như một nữ thần rừng xanh.

``Nhà Hikawa tới rồi sao?'' cô nàng mỉm cười chào chúng tôi, giọng cô ấy trong trẻo giữa đám đông. ``Mọi người mặc yukata hợp lắm.''

Hội trưởng thấy chúng tôi, lập tức chạy lại, chiếc quạt cầm tay phẩy phẩy liên tục: ``Cuối cùng cũng tới! Kuon-kun, Kaigou-chan, Suzuki-chan! Chuẩn bị tinh thần chưa? Hôm nay Hội Học sinh tụi này hơi bị bận rộn đó nha, nên mọi người cứ tự nhiên phá đảo cái lễ hội này đi!''

Tôi nhìn quanh, thấy những người bạn của mình đều đang rạng rỡ niềm vui. Nanase dù đang làm việc nhưng thi thoảng vẫn liếc mắt về phía chúng tôi mỉm cười. Inari thì có vẻ hơi run vì đám đông quá lớn, nhưng cô ấy vẫn cố gắng hoàn thành nhiệm vụ hướng dẫn.

Dưới cổng torii đỏ rực, giữa mùi hương trầm mặc của ngôi đền cổ và sự náo nhiệt của ngày hội, tôi cảm thấy một luồng cảm xúc trào dâng. Đây không chỉ là một lễ hội thông thường. Đối với chúng tôi, nó là sự cộng hưởng giữa quá khứ đau thương và một hiện tại rực rỡ mà chúng tôi đã phải đánh đổi rất nhiều mới có được.

Tôi siết nhẹ tay Naomi, nhìn Onii-sama và Onee-sama, thầm hứa trong lòng: Dù đêm nay có chuyện gì xảy ra, dù bình minh ngày mai chúng tôi có ở đâu, thì ngày hôm nay, chúng tôi sẽ sống một cuộc đời rực rỡ nhất có thể tại nơi này.

``Đi thôi nào, mọi người!'' tôi hô vang, dẫn đầu cả nhóm bước qua cổng torii, tiến vào trung tâm của lễ hội mùa hè Aogami.

\ct

Khác với Onee-sama và Onii-sama chọn đi nhóm trò chơi, tôi quyết định gia nhập nhóm càn quét ẩm thực. Thành viên nhóm gồm có tôi, Suzune-san, Honoka-san, Yae-san, Hanabi-san, Kana-san và Kanade-san. Một đội hình hoàn hảo cho những tâm hồn đam mê ăn uống.

Khu ẩm thực của Lễ hội thị trấn nằm dọc theo con đường rợp bóng cây dẫn thẳng tới điện chính của đền Aogami. Hàng chục chiếc xe đẩy bán hàng bằng gỗ được dựng san sát nhau, mỗi gian hàng lại treo những chiếc lồng đèn giấy đỏ rực đung đưa theo chiều gió. Không khí ở đây đặc quánh mùi thơm của mỡ hành, mùi khét nhẹ của nước tương nướng trên than hồng, và mùi ngòn ngọt của kẹo mạch nha.

Tiếng rao hàng của các cô chú tiểu thương vang lên không ngớt, hòa quyện với tiếng cười nói rôm rả của khách trẩy hội. Những đứa trẻ mặc yukata chạy lăng xăng trên con đường trải sỏi, tay cầm những chiếc chong chóng giấy đủ màu sắc.

``Mực nướng nguyên con đây! Nóng hổi vừa thổi vừa ăn!''

``Đá bào mát lạnh giải nhiệt mùa hè nào các cháu ơi!''

Chúng tôi vừa mới bước qua ba gian hàng đầu tiên, Suzune đã dừng lại trước một quầy bán takoyaki do một ông lão với chiếc khăn quấn ngang trán làm chủ. Ông lão thoăn thoắt lật những viên bánh tròn xoe trên chiếc chảo gang khổng lồ, tiếng xèo xèo vang lên vui tai mỗi khi ông rưới thêm một lớp bột mới.

``Ông ơi, cho cháu ba hộp ạ! Đổ thật nhiều sốt mayonnaise và cá bào cho cháu nhé!'' cô nàng hớn hở gọi, hai mắt sáng rực như sao.

Ông lão ngẩng lên, nụ cười hiền hậu làm hằn sâu những nếp nhăn trên khóe mắt: ``Có ngay, có ngay! Các cháu học ở trường Trung học Aogami phải không? Mặc yukata đẹp quá nhỉ. Hôm nay cứ ăn cho thỏa thích nhé, ông lấy giá ưu đãi cho học sinh!''

``Cháu cảm ơn ông ạ!''

Nhận lấy ba hộp takoyaki bốc khói nghi ngút, Suzune không chần chừ mà dùng tăm xiên ngay một viên nhét vào miệng. Do hôm nay không có Touka đi kèm quản lý, cô nàng ``Linh vật'' của lớp như cá gặp nước, thỏa sức ăn vô lo vô nghĩ.

``Cậu ăn thế không sợ đau bụng sao?'' tôi tròn mắt nhìn cô bạn. Hơi nóng từ viên bánh làm Suzune phải xuýt xoa, nước mắt trào ra một chút nơi khóe mi, nhưng khuôn mặt cô nàng lại ngập tràn hạnh phúc.

``Ngon nhắm\~{} Cậu thử hong Suzuki-chan?'' cô bạn nhai nhóp nhép, rồi ánh mắt cô nàng chợt lóe lên một tia sáng tinh nghịch. ``Hay là chúng mình thi ăn đi! Xem ai càn quét được nhiều gian hàng hơn!''

``Thi ăn sao? Nghe thú vị đấy!'' Honoka lập tức nhảy vào, giơ cao cây kẹo bông gòn màu hồng phấn thay cho chiếc micro. Cô nàng hắng giọng, cố bắt chước phong thái của một bình luận viên chuyên nghiệp. ``Vậy để tớ làm trọng tài cho! Đội bên trái là nhà vô địch bất bại đến từ lớp 1-C, Karamomo Suzune! Đội bên phải là kẻ thách đấu đầy bí ẩn, Hikawa Suzuki!''

Kana cũng hùa theo, vỗ tay rào rào như đang dẫn chương trình phát thanh trên đài: ``Oa! Một cuộc chiến ẩm thực không khoan nhượng! Ai sẽ là người gục ngã trước sức mạnh của carbohydrate đây!? Quý vị khán giả hãy cùng chờ xem!''

Kanade bên cạnh chỉ biết cười trừ, lấy từ trong túi áo ra một tờ giấy ăn nhỏ, nhẹ nhàng lau vết tương dính trên khóe miệng cô bạn thân: ``Kana-chan, cậu cũng đang dính đầy tương yakisoba kìa. Đừng nói lớn quá, mọi người đang nhìn kìa\dots''

Quả thực, một vài người dân thị trấn đi ngang qua đã dừng lại nhìn nhóm chúng tôi. Thế nhưng, thay vì khó chịu trước tiếng ồn ào, họ lại mỉm cười đầy thiện cảm trước sự hồn nhiên và năng lượng của những nữ sinh trung học. Bầu không khí của thị trấn Aogami này luôn mang lại một cảm giác thân thuộc và bao dung đến lạ.

Yae đang cầm một xiên mực nướng, khẽ thở dài với vẻ mặt đầy lo ngại của một người chơi thể thao chuyên nghiệp: ``Này mấy đứa, vui thì vui nhưng cũng phải giữ dáng chứ. Ăn kiểu đó ngày mai lại khóc ròng trên bàn cân cho xem. Lượng calo nạp vào phải cân bằng với lượng calo tiêu hao đấy. Nhất là khi mặc yukata, lỡ bụng to ra thì đai obi sẽ thít chặt lắm.''

``Lời nhắc nhở của Yae-san có lẽ đúng với người bình thường,'' Hanabi đang cầm một ly đá bào dâu tây, cười ranh mãnh chỉ vào tôi. ``Nhưng mà cậu ấy lại quên mất một điều tối quan trọng: Suzuki-chan có một cái dạ dày không đáy!''

``Làm gì có!'' tôi vội vàng phủ nhận, hai má nóng bừng lên.

Dù sự thật là ở thế giới cũ hay thế giới này, tôi có ăn bao nhiêu cũng không béo lên nổi. Ở thế giới u ám trước kia, việc tìm được một bữa ăn tử tế đôi khi là một thử thách. Có lẽ vì vậy mà khi đến đây, cơ thể tôi dường như lúc nào cũng ở trạng thái sẵn sàng hấp thụ mọi thứ, nhưng lạ kỳ thay, vóc dáng tôi vẫn không hề thay đổi. Một đặc điểm sinh học mà ngay cả ông Game đôi lúc cũng phải thắc mắc.

``Thế à?'' Hanabi tiến lại gần, vỗ nhẹ vào chiếc đai obi của tôi. ``Mình thắt obi cho cậu chặt thế này, mà nãy giờ cậu ăn bao nhiêu thứ rồi vẫn không thấy khó thở chút nào sao? Rõ ràng là cậu giấu một cái hố đen vũ trụ trong bụng rồi!''

Cả đám phá lên cười trước sự so sánh cường điệu của cô bạn năng nổ đó.

``Vậy thì chứng minh đi!'' Suzune đưa hộp takoyaki thứ hai ra trước mặt tôi như một lời khiêu chiến, ánh mắt rực lửa quyết tâm.

Tôi nhìn hộp bánh vàng rộm, thơm lừng mùi bạch tuộc và sốt mayonnaise, rồi chép miệng. ``Được thôi! Nếu cậu đã muốn thử sức. Nhưng nếu mình thắng, cậu phải bao mình món kakigori đá bào vị matcha đấy nhé!''

``Chơi luôn!''

Và thế là cuộc chiến ẩm thực chính thức bắt đầu.

Dưới sự cổ vũ cuồng nhiệt của Honoka và Kana, tôi và Suzune lướt qua các gian hàng như những cơn lốc nhỏ. Chúng tôi ``càn quét'' từ hàng ngô nướng bơ ngào ngạt, sang quầy kẹo hồ lô bọc đường đỏ rực. Tiếp đó, cả hai lại ghé vào sạp bán bánh xèo okonomiyaki siêu to khổng lồ, nơi một bà cô với đôi tay nhanh thoăn thoắt đang đảo bắp cải và thịt ba chỉ trên chảo phẳng.

Mỗi lần dừng lại, chúng tôi lại khiến các cô chú bán hàng phải tròn mắt ngạc nhiên trước sức ăn của hai nữ sinh nhỏ nhắn.

``Cháu gái, cháu ăn đến xiên thịt thứ năm rồi đấy! Có no quá không?'' một chú bán yakitori đội chiếc mũ rơm lo lắng hỏi tôi, tay vẫn đang quạt than phành phạch.

``Dạ không sao đâu chú, cháu vẫn còn để bụng cho món tráng miệng nữa ạ!'' tôi cười tươi đáp lại, trong khi tay nhận lấy xiên thịt gà nướng muối.

Yae và Kanade đi phía sau, vừa nhâm nhi phần ăn khiêm tốn của mình vừa cười bất lực. Hanabi thì liên tục lấy điện thoại ra quay video, miệng không ngừng bình luận: ``Các bạn đang chứng kiến kỷ lục Guiness thế giới mới tại Aogami! Hikawa Suzuki đang dẫn trước với khoảng cách nửa cái bánh xèo! Phải nói là sức chiến đấu của cô nàng quá khủng khiếp!''

Kanade thì thi thoảng lại giơ máy ảnh kỹ thuật số (mà cô mượn từ Uzuki) lên, chụp lại những khoảnh khắc phùng má trợn mắt của tôi và Suzune. ``Hai người trông dễ thương lắm,'' cô ấy khẽ cười.

Honoka vừa đi vừa cắn que kẹo bông gòn, bỗng nhiên trầm ngâm: ``Không biết nhóm của Kuon-kun bên khu trò chơi thế nào nhỉ? Với cái tính nghiêm túc của cậu ấy, chắc là đang tính toán góc độ và lực gió để vớt cá vàng cũng nên.''

``Còn Kaigou-san chắc đang đe dọa mấy ông chủ gian hàng bắn súng vì dám chỉnh lệch thước ngắm rồi!'' Kana hùa theo, và cả nhóm lại được một phen cười nghiêng ngả. Nghĩ đến Onee-sama đang xắn tay áo yukata lên cãi lý với người ta, tôi cũng không thể nhịn được cười.

Cuối cùng, sau khi đánh bay thêm hai phần taiyaki nhân đậu đỏ ở cuối con đường, Suzune đã phải giơ cờ trắng đầu hàng. Cô nàng ngồi phịch xuống một chiếc ghế đá gần đó, ôm cái bụng đã căng tròn, thở phì phò: ``M-Mình chịu thua\dots\ Bụng mình\dots\ sắp nổ tung rồi\dots''

Tôi ngồi xuống bên cạnh, đưa cho cậu ấy một chai nước khoáng lạnh mà tôi vừa mua. Thật lòng mà nói, tôi cũng đã thấy lưng lửng dạ dày, nhưng so với cảm giác no ứ của Suzune thì tôi vẫn thoải mái hơn nhiều.

``Đã bảo rồi mà, đừng bao giờ thách thức dạ dày của mình chứ,'' tôi cười khúc khích, mở nắp chai nước uống một ngụm.

``Suzuki-chan đỉnh thật đấy\dots'' Honoka vỗ tay tán thưởng. ``Vậy là giải thưởng kakigori thuộc về nhà vô địch Hikawa Suzuki!''

Chúng tôi ngồi nán lại trên chiếc ghế đá, vừa nghỉ ngơi vừa ngắm nhìn dòng người tấp nập qua lại. Gió đầu giờ chiều mùa hè thổi nhẹ, mang theo hơi nóng nhưng không hề bức bối, làm lung lay những dải đèn lồng giấy phía trên đầu. Tiếng cười đùa của trẻ con chạy nhảy quanh chân đền, tiếng trò chuyện rôm rả của các cụ già ngồi đánh cờ bên vệ đường, và tiếng xèo xèo của thức ăn trên bếp lửa. Mọi thứ đan xen vào nhau, tạo nên một bản nhạc nền êm đềm và thanh bình.

Tôi tựa lưng vào ghế, khẽ nhắm mắt lại để cảm nhận trọn vẹn từng nhịp đập của không gian xung quanh. Dù sắp phải đối mặt với một đêm đầy biến động, nhưng ở khoảnh khắc này, hòa mình vào không khí nhộn nhịp cùng những người bạn tuyệt vời này, tôi chỉ muốn để bản thân trôi theo những niềm vui bình dị nhất của tuổi thanh xuân. Một mùa hè rực rỡ và đong đầy kỷ niệm, một ngày lễ hội mà tôi biết chắc chắn rằng mình sẽ chẳng bao giờ có thể quên.

\pov{Hikawa Kaigou}

Trong khi con bé Suzu đang bận rộn càn quét khu ẩm thực thì tôi lại bị cuốn vào một chiến trường khác náo nhiệt không kém: khu trò chơi dân gian.

Trái ngược với mùi thơm nức mũi của những gian hàng đồ ăn, khu vực này tràn ngập những tiếng la hét phấn khích, tiếng súng hơi lốp bốp, tiếng còi tuýt tuýt và cả những tràng cười giòn tan của trẻ nhỏ. Những dải đèn lồng ở đây dường như cũng được thắp sáng rực rỡ hơn, đung đưa theo nhịp gió biển thổi vào từ vịnh Aogami, tạo nên một không gian huyền ảo đầy màu sắc.

Nhóm của tôi, ngoài tôi ra thì có Shiina, Aku, Azuki, Akane và Yuka. Lúc nãy Honoka bảo sẽ đi cùng Suzu ăn uống, nhưng chưa đầy mười lăm phút sau đã thấy cô nàng chạy thục mạng sang khu này, một tay cầm ly đá bào dâu tây ăn dở, tay kia lôi xềnh xệch Azuki nhập hội. Đúng là sức lực của tuổi trẻ, chẳng bù cho tôi, đi lại nãy giờ dưới thời tiết mùa hè này cộng thêm bộ yukata quấn chặt người cũng đã bắt đầu thấy oi bức.

``A! Bên kia có quầy vớt cá vàng kìa!'' cô bạn Cáo sa mạc reo lên, chỉ tay về phía một bể nước nông bằng nhựa xanh đang thu hút rất đông người.

Azuki liền hùa theo, đôi mắt sáng rực: ``Thử xem sao! Cậu biết đấy, mình luyện tập độ dẻo dai của cổ tay bằng việc gảy đàn guitar suốt ngày nên chắc chắn vớt cá sẽ rất cừ!''

Hai cô nàng hào hứng kéo nhau chạy lại gian hàng kingyo-sukui\footnote{Trò chơi vớt cá vàng bằng một cái vợt giấy mỏng manh.}. Akane cũng vội vã chạy theo, tay lăm lăm chiếc máy ảnh kỹ thuật số (mà cô bạn mượn được từ Trưởng CLB Nhiếp ảnh): ``Chờ mình với! Mình phải chụp lại cảnh Azuki-chan làm rách vợt giấy tơi tả mới được!''

Tôi lắc đầu cười trừ, thong thả đi dạo dọc theo dãy hành lang của khu trò chơi cùng với Shiina, Aku và Yuka. Cô nàng bờm Mèo thì vẫn giữ nguyên cái điệu bộ buồn ngủ thường ngày, hai mắt lim dim như thể có thể gục xuống đất bất cứ lúc nào. Nhưng kì lạ thay, khi chúng tôi dừng lại trước quầy yo-yo tsuri \footnote{Trò chơi câu bóng nước.}, cô nàng chỉ nhắm hờ mắt, tay cầm chiếc cần câu bằng giấy mỏng manh móc nhẹ một cái đã câu được ngay quả bóng màu tím có hoa văn sặc sỡ, trong khi mấy đứa trẻ bên cạnh làm rách lả tả đến cái móc thứ năm.

``Giỏi vậy?'' Aku tròn mắt ngạc nhiên, tiến lại gần nhìn chằm chằm vào quả bóng nước. ``Nhìn cậu như đang mộng du mà sao lại câu dính hay thế? Có bí quyết gì không?''

Yuka ngáp dài một cái, khẽ đung đưa quả bóng nước bằng sợi dây cao su: ``Không biết nữa\dots\ Chắc do quả bóng đó\dots\ nó tự dính vào móc đấy\dots''

Tôi đồ rằng năng lực ngoại cảm tiềm ẩn của cô nàng đã giúp một tay để ổn định mặt nước, nhưng thôi, ở lễ hội thì ai quan tâm đến mấy định luật vật lý cơ chứ.

Tiếp tục tiến sâu vào trong, chúng tôi đụng độ gian hàng Wanage - trò ném vòng quen thuộc mà tôi đã bắt gặp trong lễ hội trường học hôm qua. Những phần thưởng được xếp trên các bệ gỗ hình nón với độ xa gần khác nhau. Gần nhất là mấy bịch kẹo dẻo rẻ tiền, xa hơn là đồ chơi nhỏ, và ở vị trí xa nhất, khó ném trúng nhất là một chú gấu bông khổng lồ hình Aogami-kun - linh vật của thị trấn.

``Đến lượt mình thể hiện!'' Aku xắn tay áo yukata lên, móc tiền xu đưa cho người bán hàng để đổi lấy ba chiếc vòng nhựa. ``Dăm ba cái trò ném vòng tính toán quỹ đạo này sao làm khó được Mizuno Aku!''

Cô nàng lùi lại một bước, nhắm một mắt lại, ngắm nghía rất kỹ mục tiêu là hộp mô hình nhân vật anime ở hàng thứ hai. Khuôn mặt cô ấy toát lên vẻ tập trung cao độ, hệt như đang chuẩn bị tung tuyệt chiêu cuối trong game đối kháng.

*VÚT!* Chiếc vòng thứ nhất bay ra khỏi tay Aku\dots\ và rơi tõm xuống đất, cách mục tiêu cả mét.

*VÚT!* Chiếc vòng thứ hai sượt qua góc hộp mô hình rồi nảy tưng tưng ra ngoài.

*VÚT!* Chiếc vòng cuối cùng, với một lực ném quá đà, bay vèo qua đầu người bán hàng, khiến ông chú phải cúi rạp người né tránh, chiếc vòng an tọa gọn gàng trong thùng rác phía sau quầy.

``\textbf{Aaaa! Sao lại thế được!}'' Aku ôm đầu than vãn, gục ngã trước thực tại tàn khốc. ``Rõ ràng là mình đã tính toán quỹ đạo bay của nó theo phương trình ném xiên hoàn hảo rồi mà! Gió! Chắc chắn là do lực cản của gió!''

Shiina khoanh tay trước ngực, thở dài đầy ngao ngán, nhìn Aku với ánh mắt thương hại: ``Cậu tính toán kiểu gì mà cái vòng bay như boomerang thế hả? Trò này không có aim-assist (thước ngắm hỗ trợ) như trong game FPS đâu. Để đó cho tôi.''

Đội trưởng câu lạc bộ bóng rổ bước lên, chỉ cầm đúng một chiếc vòng duy nhất. Cô ấy không cần lấy đà, uốn cong đầu gối một cách nhịp nhàng, thực hiện một động tác ném ném rổ hoàn hảo với tư thế chuẩn mực. Chiếc vòng vẽ một đường cong tuyệt mỹ trên không trung rồi nhẹ nhàng hạ cánh trúng phóc vào phần thưởng ở hàng giữa — một bộ bài hoa Hanafuda cực xịn.

``Tuyệt vời!'' Tôi vỗ tay tán thưởng. Quả không hổ danh là át chủ bài của đội bóng rổ Aogami.

``Hừm. Chuyện nhỏ thôi,'' cô bạn bờm Sư tử đắc ý cười, nhận lấy phần thưởng từ ông chú bán hàng đang há hốc mồm. ``Aku, trò này cần sự dẻo dai của cổ tay và cảm giác không gian thực tế, không phải cứ bấm phím A liên tục trên tay cầm là vòng sẽ tự bay trúng đâu.''

Bị chạm tự ái, đôi má của cô nàng Hầu gái đỏ lên: ``Chỉ là do trò ném vòng không hợp với meta của mình thôi! Nếu là trò cần sự chính xác tuyệt đối như bắn súng thì mình chắc chắn sẽ không thua! Mình là rank (hạng) Thách Đấu của game bắn súng sinh tồn đấy!''

Nói là làm, Aku lôi tuột chúng tôi sang gian hàng shateki\footnote{Trò bắn súng hơi trúng mục tiêu, sử dụng đạn là nút cao su.} ngay bên cạnh. Gian hàng này trưng bày vô số phần thưởng hấp dẫn trên một chiếc kệ gỗ nhiều tầng, từ bánh kẹo cho đến những máy chơi game đắt tiền. Ông chủ quầy là một người đàn ông trung niên với bộ ria mép tỉa tót kỹ lưỡng, đang hờ hững lau chùi mấy khẩu súng hơi bằng gỗ, ánh mắt toát lên vẻ tự tin của một kẻ biết chắc mình không bao giờ lỗ.

Aku đập tiền xuống bàn, dõng dạc: ``Cho cháu năm viên đạn!''

Nhận lấy khẩu súng, phong thái của cô nàng bỗng nhiên thay đổi hoàn toàn. Khí chất hậu đậu, ngớ ngẩn thường ngày biến mất. Ánh mắt cô nàng trở nên sắc lạnh, tư thế đứng chụm hai chân, một tay đỡ báng súng, một tay kẹp cò cực kỳ chuyên nghiệp. Nhìn cô nàng lúc này chẳng khác gì một tay bắn tỉa thực thụ chuẩn bị thực hiện một phi vụ ám sát.

*ĐOÀNG!* Viên đầu tiên trúng ngay hộp kẹo dẻo ở hàng đầu, làm nó rơi phịch xuống bạt.

*ĐOÀNG! ĐOÀNG! ĐOÀNG!* Ba viên tiếp theo liên tiếp hạ gục ba hộp bánh quy ở hàng giữa một cách vô cùng chính xác. Tốc độ ngắm bắn và bóp cò của Aku nhanh đến mức người dân xung quanh bắt đầu ngừng chơi để quay sang ồ lên tán thưởng.

Ông chủ quầy vuốt ria mép, trán hơi rịn mồ hôi.

Tuy nhiên, ở viên đạn cuối cùng, khi Aku tự tin nhắm vào hộp mô hình đắt tiền ở hàng thứ ba.

*ĐOÀNG!* Viên đạn súng hơi bằng nút bần trúng ngay giữa hộp mô hình\dots\ nhưng nó chỉ khẽ nhúc nhích chứ không hề rơi xuống khỏi bệ đỡ.

``Cái gì!?'' Aku há hốc mồm, dụi mắt mấy lần. ``Trúng ngay hồng tâm mà không đổ sao? Cái định luật vật lý kỳ quặc gì thế này!?''

Ông chủ quầy cười đắc ý, lấy chiếc khăn lau mồ hôi trên trán: ``Hehe, xin lỗi nhé cháu gái. Nhắm chuẩn lắm, nhưng hộp đó hơi nặng, đạn súng hơi của chú lực không đủ để đẩy nó xuống đâu. Lần sau may mắn hơn nhé!''

``Rõ ràng là game bị lỗi mà\dots\ họ đã cấu kết với nhau từ trước rồi\dots'' Aku lẩm bẩm, mặt xị xuống trông rõ tội nghiệp. Dù đã có chiến lợi phẩm là một đống bánh kẹo, nhưng với lòng kiêu hãnh của một game thủ FPS, không hạ được trùm cuối là một sự sỉ nhục.

Tôi đứng khoanh tay quan sát nãy giờ, khẽ nhếch mép cười khẩy. Tôi bước lên phía trước, thong thả gõ ngón tay xuống mặt bàn gỗ.

``Chú à, cho cháu mười viên.'' Tôi đặt tờ tiền giấy thẳng thớm lên bàn.

``Ồ, lại một cô bé nữa muốn thử sức sao? Cứ tự nhiên! Để chú thay đạn cho.'' Ông chủ quầy vui vẻ nạp đạn, nghĩ bụng lại có thêm một con gà béo để vặt lông.

Tôi cầm lấy khẩu súng hơi, khẽ thở một hơi dài. Đưa nó lên ngang tầm mắt, tôi bắt đầu điều chỉnh tầm ngắm. Trước đó, tôi đã từng nghiên cứu cách vận hành của trò chơi này như thế nào, và cộng với am hiểu về vật lý, tôi cũng đã cho mình câu trả lời có đáp án đẹp nhất cho trò chơi này.

Tôi nhắm mắt lại một giây, điều chỉnh nhịp thở hòa làm một với nhịp đập của trái tim, rồi từ từ mở mắt ra. Đôi mắt tôi khóa chặt vào mục tiêu.

*ĐOÀNG!* Viên đạn đầu tiên bay trúng phần cạnh của chiếc hộp mô hình lúc nãy, làm nó xoay đi một góc nhỏ.

Ông chủ quầy bật cười ha hả: ``Chú đã bảo rồi, bắn kiểu đó không đổ đâu cháu ơi\dots''

*ĐOÀNG! ĐOÀNG! ĐOÀNG! ĐOÀNG!*

Tôi không đợi ông ta nói hết câu, bóp cò liên tục bốn phát với tốc độ kinh hồn mà không cần ngắm lại. Bốn viên đạn bay theo cùng một quỹ đạo tuyệt đối, găm trúng chính xác vào cùng một điểm ở mép hộp mô hình đang bị xoay lệch. Lực tác động liên tiếp vào cùng một điểm tạo ra mô-men xoắn, làm chiếc hộp mất thăng bằng hoàn toàn và rơi uỵch xuống đất.

Cả khu vực xung quanh bỗng nhiên im lặng phăng phắc. Tiếng cười của ông chủ tắt lịm. Aku tròn mắt nhìn tôi như thể đang nhìn một kẻ sử dụng ngắm tự động trong game. Shiina và Yuka thì ồ lên kinh ngạc. Ngay cả điếu thuốc lá ông chủ quầy đang ngậm trên môi cũng rơi tõm xuống đất.

``Nếu một viên không đủ lực đẩy thì dùng bốn viên liên tiếp. Định luật bảo toàn động lượng cơ bản thôi mà.'' Tôi thổi nhẹ nòng súng bằng gỗ một cách điệu nghệ, nhếch mép cười đắc thắng.

Chưa dừng lại ở đó, tôi lia nòng súng sang mục tiêu lớn nhất ở hàng cuối cùng: Một bộ máy console phiên bản giới hạn, thứ mà chắc chắn Onii-sama sẽ rất thích để giết thời gian. Nó được đặt nằm ngang, tiếp xúc diện rộng với bệ đỡ nên vô cùng chắc chắn, gần như không thể dùng đạn nút bần để đẩy ngã trực tiếp.

``Năm viên còn lại\dots\ kết thúc nhé.''

Lần này, tôi không bắn vào chiếc hộp, mà nhắm thẳng vào hai thanh gỗ nhỏ mỏng manh đang kê phía dưới nó để tạo độ nghiêng.

*ĐOÀNG! ĐOÀNG!* Thanh gỗ kê bên trái gãy gập làm đôi.

*ĐOÀNG! ĐOÀNG!* Thanh gỗ kê bên phải vỡ nát dưới áp lực.

Chiếc hộp console mất hoàn toàn điểm tựa, trượt khỏi bệ đỡ và lao thẳng xuống tấm bạt phía dưới trước sự ngỡ ngàng tột độ của tất cả mọi người.

``B\dots\ Bingo!'' ông chủ quầy lắp bắp vỗ tay một cách vô thức, mồ hôi hột tuôn ra như tắm. Lần này thì ông ta lỗ nặng thật rồi, gặp ngay phải dân chuyên nghiệp.

Tiếng hò reo bùng nổ dữ dội xung quanh gian hàng. Aku lao tới ôm chầm lấy tôi, nhảy cẫng lên sung sướng: ``Kaigou-san! Cậu đỉnh quá! Làm thế nào mà bắn chuẩn như hack thế!? Có nhận mình làm đệ tử không?''

``Chỉ là một chút quan sát với kiến thức vật lý cơ bản thôi,'' tôi nhún vai, thản nhiên nhận lấy chiếc máy console khổng lồ và hộp mô hình từ tay ông chủ đang mặt mày nhăn nhó, ruột đau như cắt.

Đúng lúc đó, Honoka, Azuki và Akane cũng vừa chạy tới. Tóc Honoka ướt đẫm mồ hôi, nhưng trên tay cô nàng đang cầm đến ba túi ni-lông chứa đầy cá vàng bơi tung tăng.

``Mọi người ơi! Tụi này vớt được cả đống cá luôn này!'' cô bạn Cáo sa mạc hớn hở khoe chiến tích. Nhưng rồi ánh mắt cô nàng khựng lại khi nhìn thấy chiến lợi phẩm khổng lồ trên tay tôi. ``K-Kaigou-san\dots\ Cậu vừa đi cướp ngân hàng đấy à? Sao lại có cả máy chơi game thế kia?''

Tôi bật cười, đưa hộp mô hình cho Aku đang nhảy nhót xung quanh và xách hộp máy console lên vai một cách nhẹ nhàng: ``Không, chỉ là lấy lại một chút công bằng cho các cậu ấy thôi.''

Nhìn những nụ cười rạng rỡ của mọi người xung quanh, nhìn Akane vui vẻ giơ máy ảnh lên bấm tanh tách để lưu lại khoảnh khắc này, và nghe tiếng pháo giấy nổ lẹt đẹt từ đằng xa, tôi bất giác cảm thấy một sự bình yên đến lạ thường. Khác với những chiến thắng ngập tràn mùi máu tanh và khói thuốc súng ở quá khứ, chiến thắng nhỏ bé và có phần ``ăn gian'' ở quầy bắn súng này mang lại cho tôi một niềm vui thuần khiết, đơn giản, và\dots\ rất con người.

Có lẽ, tôi thực sự bắt đầu yêu thích cái thế giới bình dị này rồi. Việc bảo vệ những nụ cười vô tư này, bảo vệ lễ hội này, dường như đã trở thành một phần mục đích sống mới của tôi.

Nhưng tận sâu trong thâm tâm, nhịp đập hối hả của một chiến binh vẫn nhắc nhở tôi rằng, sự bình yên này chỉ là phần nổi của tảng băng chìm. Thời gian của chúng tôi tại thị trấn Aogami đang đếm ngược từng giờ, và đêm nay, dưới ánh trăng rằm, một chương mới sẽ mở ra.

\pov{Naomi}

Đây là lần đầu tiên tôi được nhìn thấy một thế giới rực rỡ và tràn đầy màu sắc đến thế.

Trước kia, trong ký ức mờ nhạt của tôi, mọi thứ chỉ là bóng tối vô tận và tiếng gào thét câm lặng của những linh hồn bị lãng quên. Tôi là một đứa trẻ của quá khứ, một sinh mạng đã bị tước đoạt từ 1300 năm trước trong một nghi thức hiến tế tàn khốc để xoa dịu cơn giận dữ của thần linh và bảo vệ thị trấn Aogami này. Tôi đã bị chôn sống dưới gốc cây cổ thụ sâu trong rừng, linh hồn bị xích lại bởi những lời nguyền cổ xưa, vĩnh viễn không thể siêu thoát. Suốt hơn một thiên kỷ, tôi chỉ là một cái bóng vất vưởng trong ranh giới giữa sự sống và cái chết, cho đến khi Onii-chan và mọi người xuất hiện, cho tôi một cuộc sống mới và một gia đình để thuộc về.

Nhưng bây giờ, tôi đang ở đây, tại Lễ hội mùa hè của thị trấn Aogami.

``Naomi-chan, em đi cẩn thận kẻo vấp ngã đấy\~{}'' một giọng nói dịu dàng vang lên bên tai tôi, kéo tôi về với thực tại.

Đó là Saya-san. Chị ấy nắm lấy bàn tay nhỏ bé của tôi, bàn tay chị ấy thật ấm áp và mềm mại. Chị ấy lúc nào cũng toát lên một vẻ ân cần như một người mẹ hiền, cẩn thận dắt tôi đi qua dòng người đông đúc để tránh việc tôi bị xô đẩy.

``Naomi-chan mặc bộ yukata hoa anh đào này hợp lắm luôn ấy!'' chị Miku đi ngay bên cạnh, ánh mắt lấp lánh nhìn tôi đầy tự hào. ``Chị đã phải thức trắng hai đêm để tự tay thêu thêm mấy họa tiết kim tuyến dọc theo phần tà váy này đấy. Quả nhiên là mắt thẩm mỹ của chị không bao giờ sai! Em lộng lẫy hệt như một nàng công chúa phương Đông vậy!''

``Em cảm ơn chị Miku ạ,'' tôi bẽn lẽn đáp lại, dùng tay vuốt ve lớp vải lụa mềm mại của bộ trang phục. Lần đầu tiên khoác lên người một bộ trang phục truyền thống khác hẳn so với bộ kimono trước đây, tôi cảm thấy mình như biến thành một con người thực sự, một sinh mệnh đang hô hấp hòa nhịp với nhịp đập của thị trấn.

Đi phía trước chúng tôi là hai chị Yuki và Kaoru. Hai người họ vô cùng năng động, lúc nào cũng ríu rít nói chuyện. Chị Yuki thi thoảng lại cất giọng ngân nga một điệu nhạc dân gian, giai điệu vui tươi hòa nhịp cùng tiếng trống taiko vang dội từ phía đài yagura giữa sân đền. Chị Kaoru thì luôn miệng đóng vai trò như một hướng dẫn viên du lịch chuyên nghiệp, nhiệt tình giới thiệu cho tôi từng gian hàng một mà chúng tôi đi qua.

``Bên trái là gian hàng taiyaki siêu ngon của nhà Kamiyama, bên phải là quầy bán mặt nạ cáo! À, phía trước có bán wataame kìa! Naomi-chan có muốn ăn thử không? Ngọt lắm đó!''

Mắt tôi sáng rực lên khi nhìn thấy một chiếc máy inox nhỏ đang quay tít, biến những hạt đường tơi xốp thành một đám mây hồng khổng lồ bồng bềnh như có phép thuật.

``Wataame\dots\ là kẹo bông gòn ạ?'' tôi ngước nhìn đám mây hồng đó một cách say sưa. Làm sao người ta có thể ăn một đám mây được chứ? Khối lượng riêng của nó chắc chắn rất nhẹ, cấu trúc phân tử đường bị kéo giãn nhờ lực ly tâm và nhiệt độ\dots

``Đúng rồi! Ngon tuyệt cú mèo luôn!'' chị tóc ngắn gật đầu lia lịa.

Tôi đang định nhờ chị Saya mua giúp để nghiên cứu cấu trúc vật lý của nó, thì bỗng nhiên một cây kẹo bông gòn to bằng cả khuôn mặt tôi đã được chìa ra trước mũi.

Tôi ngơ ngác quay lại. Là chị Koishin. Chị khoanh tay trước ngực, quay mặt đi chỗ khác để tránh ánh mắt của tôi, hai má hơi ửng hồng dưới ánh sáng của những chiếc lồng đèn giấy: ``Đừng có hiểu lầm nhé! Chẳng qua là ông chú bán kẹo làm lố tay to quá, chị ăn không hết nên tiện tay đưa cho em thôi! Không ăn thì chị ném đi đấy!''

Chị Yuki bật cười khúc khích, chọc ghẹo: ``Rõ ràng là nãy giờ tớ thấy cậu đứng cãi cọ với chú ấy, bắt chú ấy phải làm một cây thật to màu hồng vì biết Naomi-chan thích màu hồng mà! Koishin-chan chẳng bao giờ thành thật với bản thân cả!''

``C-Cậu im đi, Yuki! Ai thèm quan tâm con bé thích màu gì chứ!'' chị ấy lắp bắp, lườm cô bạn thân một cái sắc lẹm nhưng không giấu nổi sự ngượng ngùng.

Tôi mỉm cười nhận lấy cây kẹo bông gòn khổng lồ từ tay cô chị bằng mặt không bằng lòng ấy: ``Em cảm ơn Koishin-san nhiều lắm ạ.'' Cắn một miếng nhỏ, vị ngọt lịm tan ngay trên đầu lưỡi, nhẹ bẫng và biến mất nhanh chóng như một ảo ảnh, nhưng dư vị ngọt ngào thì vẫn còn đọng lại mãi trong cổ họng. Tôi nheo mắt cười hạnh phúc. Chị Koishin dù ngoài miệng hay cằn nhằn, nhưng trái tim lại vô cùng dịu dàng và ấm áp.

Chúng tôi tiếp tục dạo bước qua các gian hàng. Sự náo nhiệt của lễ hội không làm tôi cảm thấy sợ hãi hay ngột ngạt như tôi từng phân tích dựa trên số liệu về mật độ đám đông. Ngược lại, những âm thanh ồn ào hỗn tạp ấy, hơi nóng hầm hập phả ra từ các lò nướng đồ ăn, tiếng cười đùa rôm rả của khách trẩy hội\dots\ tất cả đan xen vào nhau, tạo nên một bản giao hưởng tuyệt đẹp của sự sống.

``A! Quầy câu bóng nước yo-yo tsuri kìa!'' chị Kaoru chỉ tay về phía một chiếc bể cao su hình tròn chứa đầy nước, bên trong là hàng chục quả bóng bay nhỏ chứa nước đủ màu sắc đang nổi lềnh bềnh.

``Naomi-chan, em muốn thử chơi không?'' chị Saya cúi xuống, vén lọn tóc vương trên trán tôi rồi hỏi.

Tôi gật đầu háo hức. Tôi đã từng nhìn thấy thông tin về trò chơi này trong cơ sở dữ liệu của Game, cũng như trong ký ức của anh Kuon, nhưng tự mình điều khiển cơ thể để trải nghiệm thì chưa bao giờ.

Ông chú bán hàng thân thiện với chiếc tạp dề ướt sũng đưa cho tôi một chiếc cần câu nhỏ xíu làm bằng giấy bện: ``Nào, cô bé xinh xắn, thử xem có câu được quả nào không nhé. Giấy mỏng lắm, cẩn thận kẻo đứt đấy. Cứ thong thả thôi.''

Tôi cầm chiếc cần câu, căng thẳng nhìn xuống mặt nước dao động liên tục. Những quả bóng bay lắc lư theo từng gợn sóng. Lần đầu tiên tôi thử đưa móc kim loại vào phần thun của quả bóng màu xanh lục, nhưng do chưa quen kiểm soát lực tay, mặt nước gợn sóng đẩy quả bóng ra xa, khiến chiếc móc trượt đi và sức nặng của nước làm rách toạc phần dây giấy.

``A\dots'' tôi khẽ kêu lên, cảm thấy hơi hụt hẫng. Lần đầu tiên thử thất bại rồi.

``Đừng vội, Naomi-cha\~{},'' chị Saya dịu dàng vòng tay qua người tôi, nhẹ nhàng cầm lấy bàn tay đang giữ chiếc cần câu thứ hai của tôi. Hơi ấm từ cơ thể chị ấy truyền sang, bao bọc lấy tôi, giúp nhịp tim đang đập nhanh của tôi bình tĩnh lại. ``Trò này cần sự kiên nhẫn. Em phải chờ khi quả bóng trôi lại thật gần, giữ tay thật vững, và kéo lên một cách dứt khoát theo phương thẳng đứng. Nhịp nhàng thôi nhé, đừng dùng lực cổ tay nhiều quá nha\~{}''

Làm theo lời hướng dẫn của chị, tôi nhắm mắt lại hít một hơi thật sâu, rồi mở mắt ra, tập trung cao độ. Lần này, tôi nhắm vào một quả bóng màu hồng có điểm xuyết những hoa văn hình cánh hoa anh đào, giống hệt với họa tiết trên bộ yukata tôi đang mặc.

Chờ đợi. Quan sát sự chuyển động của mặt nước. Tính toán quỹ đạo trôi của quả bóng.

Khi quả bóng chầm chậm trôi ngang qua ngay dưới mũi cần câu, tôi từ từ hạ chiếc móc xuống. Móc câu bằng kim loại nhẹ nhàng lọt qua sợi dây thun đàn hồi của quả bóng.

``Bây giờ! Kéo lên đi em!'' Kaoru-san khích lệ.

Tôi khẽ nhấc cổ tay lên, dùng một lực vừa đủ nhưng rất dứt khoát. Quả bóng bay vọt lên khỏi mặt nước, kéo theo những giọt nước bắn tung tóe thành những hạt pha lê nhỏ xíu lấp lánh dưới ánh đèn lồng, và nằm gọn trong lòng bàn tay tôi một cách êm ái. Chiếc cần câu giấy vẫn còn nguyên vẹn!

``Oa! Câu được rồi! Giỏi quá Naomi-chan!'' Miku và Kaoru đồng thanh reo lên, vỗ tay rào rào cổ vũ.

Ông chú bán hàng cũng cười lớn, giơ ngón tay cái lên: ``Tuyệt quá! Cháu gái có năng khiếu đấy! Quả bóng này là của cháu!''

Tôi nắm chặt quả bóng trong tay, xỏ ngón tay vào sợi dây thun, cảm nhận sự đàn hồi của lớp cao su mỏng manh và tiếng sột soạt vui tai của dòng nước bên trong. Một cảm giác thành tựu nhỏ bé nhưng lại vô cùng lớn lao dâng trào trong lồng ngực. Tôi ngẩng mặt nhìn mọi người xung quanh, và nở một nụ cười rạng rỡ nhất từ trước đến nay. Những nụ cười phát xuất từ tận sâu thẳm tâm hồn, một nụ cười của một sinh mệnh sống thực sự đang tận hưởng niềm vui của cuộc đời, một thứ mà tôi đã hằng mong ước bấy lâu nay suốt một thiên niên kỷ.

Nhìn thấy tôi cười, cả nhóm cũng cười theo. Một vài người qua đường nhìn thấy cảnh tượng đó cũng mỉm cười thiện cảm.

Chị Miku vội vàng lấy điện thoại ra, chụp liên tục mấy chục tấm hình: ``Đáng yêu quá đi mất! Góc nghiêng này, nụ cười này! Chị muốn đem em về nhà nuôi luôn quá!''

``Không được đâu Miku-chan, Naomi-chan là báu vật vô giá của nhà Hikawa đấy! Cậu mà bắt cóc con bé là Kuon-kun sẽ đồ sát cậu đó!'' chị Yuki cười trêu chọc, khiến cả nhóm được một phen cười nghiêng ngả.

``Với cả, Miku-chan lúc bị hấp dẫn bởi Naomi-chan trông đáng yêu thật đó\~{}'' chị Saya khẽ cười.

Tôi khẽ ôm quả bóng yo-yo vào lòng, để mặc những giọt nước mát lạnh thấm vào lớp vải lụa, ánh mắt hướng về phía xa xa, nơi đài Yagura đang rộn ràng tiếng trống taiko chuẩn bị cho vũ điệu Bon Odori.

Báu vật của nhà Hikawa sao? Tôi vốn chỉ là một linh hồn vô định, một kẻ đáng lẽ ra đã bị phân mảnh và xóa sổ vĩnh viễn khỏi bờ cõi này nếu không có sự can thiệp của ba người họ. Đêm nay, dưới ánh trăng rằm tháng bảy, Nghi thức Thanh tẩy sẽ diễn ra. Anh Kuon, chị Kaigou, và chị Suzuki sẽ phải trở về thế giới thực của họ, rời xa thị trấn Aogami này mãi mãi.

Còn tôi? Lựa chọn của tôi sẽ là gì?

Trở thành một cư dân thực sự của Aogami, một con người bình thường sống một cuộc đời bình yên, được các chị yêu thương và chứng kiến gia đình Hikawa rời đi? Hay sẽ bước qua cánh cổng không gian cùng họ, dấn thân vào một tương lai vô định, đánh cược sự tồn tại của bản thân nhưng luôn có gia đình ở bên cạnh?

Tôi ngước nhìn bầu trời xanh ngắt, nắng vàng rực rỡ len qua từng tán cây cổ thụ trong sân đền. Mùi hương trầm thoang thoảng từ ngôi điện chính bay tới, hòa quyện cùng mùi thơm ngầy ngậy của lễ hội.

Dù tương lai có ra sao, dù bình minh ngày mai mang đến điều gì đi chăng nữa, thì những khoảnh khắc tuyệt diệu này, tình yêu thương vô bờ bến của chị Saya, sự nhiệt thành chân chất của các chị, và hơi ấm của thị trấn Aogami bao dung này\dots\ tôi sẽ khắc sâu tất cả vào linh hồn mình, mãi mãi không bao giờ quên.

\pov{Misora Shino}

``Thật kỳ diệu\dots'' tôi khẽ thì thầm, vươn tay ra đón lấy một cánh hoa anh đào lạc mùa vừa bị cơn gió biển thổi bay ngang qua.

Trước đây, trong những vòng lặp dường như vô tận của thị trấn Aogami, tôi chỉ là một bóng hình mắc kẹt giữa sự sống và cái chết. Tôi đã chết trong trận sạt lở đất nhiều thập kỷ trước, nhưng linh hồn tôi không thể siêu thoát, cứ mãi bị giam cầm trong vòng xoáy luân hồi của ngôi trường này. Mỗi vòng lặp, tôi lại thức dậy, lại nhìn thấy những khuôn mặt quen thuộc nhưng không ai nhớ tôi, lại chứng kiến bi kịch tái diễn, rồi lại chìm vào bóng tối. Lặp đi lặp lại. Lặp đi lặp lại. Cho đến khi tôi gần như quên mất mình là ai, cho đến khi ``Sora'' - mảnh linh hồn trốn chạy khỏi thực tại - suýt nữa nuốt chửng lấy nhận thức cuối cùng của tôi.

Khát vọng duy nhất mà tôi ôm ấp suốt những năm tháng cô độc ấy thật giản dị: được sống một cuộc đời học đường bình thường. Được có bạn bè, được cười đùa, được tham gia lễ hội cùng mọi người. Một ước nguyện nhỏ bé đến mức tàn nhẫn, bởi nó luôn nằm ngoài tầm với của một linh hồn bị giam trong vòng lặp.

Và lúc này đây, tôi đang khoác trên mình bộ yukata lụa trắng điểm xuyết họa tiết lá phong xanh biếc. Dưới chân tôi là đôi guốc mộc geta vang lên từng tiếng lách cách vui tai trên mặt sỏi. Xung quanh tôi là tiếng cười nói, là âm nhạc, là sự sống đang cuộn trào mạnh mẽ.

Tất cả những điều tuyệt diệu này, sinh mệnh mới này, tôi có được là nhờ gia đình Hikawa. Ba con người đến từ một thế giới khác đã phá vỡ vòng lặp, đối đầu với thực thể hắc ám, và kéo tôi ra khỏi vực thẳm. Kuon-san đã hứa sẽ cho tôi thấy một thế giới tươi đẹp. Và cậu ấy đã giữ đúng lời hứa.

``Shino-san! Cậu đứng đực ra đó làm gì thế? Lại đây giúp một tay đi!'' giọng nói lanh lảnh của Hội trưởng Ayame cắt ngang dòng suy tưởng của tôi.

Tôi quay đầu lại. Cậu ấy đang đứng chống nạnh trước quầy thông tin di động của Hội Học sinh, trán lấm tấm mồ hôi. Bên cạnh là Wata vẫn giữ phong thái điềm tĩnh thường ngày, tay thoăn thoắt phân loại các tờ rơi giới thiệu về các tiết mục trong lễ hội.

Tôi bước tới, khẽ mỉm cười: ``Có chuyện gì khiến Hội trưởng phải cuống cuồng lên vậy? Bọn trẻ đi lạc sao?''

``Còn hơn cả đi lạc ấy chứ!'' Ayame nhăn mặt, đưa tay phẩy phẩy chiếc quạt giấy. ``Gian hàng bắn súng shateki đằng kia vừa gọi bộ đàm báo là có một học sinh năm nhất vừa `càn quét' sạch sành sanh phần thưởng của họ, đến cả cái máy chơi game đắt tiền nhất cũng bị nẫng mất. Lão chủ quán đang khóc ròng đòi Hội Học sinh xuống can thiệp phân xử kia kìa!''

Tôi che miệng cười khẽ. Khỏi cần nói cũng biết thủ phạm là ai. Ở thị trấn này, người có khả năng làm như vậy với phong thái lạnh lùng như sát thủ đó chỉ có thể là Kaigou-san mà thôi.

``Việc này nằm ngoài phạm vi quản lý của Hội Học sinh chúng ta rồi, Ayame,'' Wata cười gượng, cất giọng đều đều. ``Giao dịch diễn ra hoàn toàn hợp lệ. Người bán đưa ra luật chơi, người chơi tuân thủ luật và giành chiến thắng. Lão chủ quán không có cớ gì để phàn nàn cả. Nếu lão ta tiếp tục làm ồn, chúng ta có quyền thu hồi giấy phép gian hàng vì tội gây rối trật tự lễ hội đó.''

Hội trưởng thở dài thườn thượt, úp mặt xuống bàn: ``Biết là vậy, nhưng quản lý cái lễ hội quy mô cỡ này đau đầu thật đấy. Sao lúc trước cậu lại xúi tôi lên làm Hội trưởng cơ chứ, Wata?''

``Vì năng lực lãnh đạo xuất chúng của cậu, và vì mình muốn xem cậu vất vả bù đầu ra sao,'' cậu ấy đáp tỉnh rụi, khiến Ayame tức nghẹn họng, mặt đỏ bừng nhưng không phản bác được câu nào.

Đúng lúc đó, Nanase và Inari cũng vừa hớt hải chạy tới từ phía cổng torii. Hai cô nàng đeo băng tay đỏ chót có in chữ ``Ban Kỷ Luật'', trông vô cùng ra dáng và nghiêm túc.

``Báo cáo Hội trưởng!'' Nanase đứng nghiêm, giơ tay chào kiểu quân đội. ``Khu vực cổng chính và đường hầm đền đã được phân luồng xong. Dòng người đang di chuyển rất trật tự, không xảy ra hiện tượng chen lấn xô đẩy. Ban tổ chức của gia đình Furukawa cũng đã cử thêm mười vệ sĩ mặc thường phục hỗ trợ duy trì an ninh!''

Cô bạn tóc trắng đứng cạnh cô Tiểu thư, chống hai tay lên đầu gối thở dốc, hai má đỏ bừng: ``M-Mọi người\dots\ đông quá\dots\ Mình vừa phải chỉ đường cho ba nhóm khách du lịch liên tiếp\dots\ C-Cổ họng mình khô hết cả rồi\dots''

Tôi bước tới, đưa cho Inari một chai nước khoáng mát lạnh mà tôi đã cẩn thận ướp sẵn từ trong thùng đá của Hội. ``Làm tốt lắm, Inari-san. Cậu uống chút nước đi cho đỡ mệt.''

Inari nhận lấy chai nước, ánh mắt rưng rưng cảm động nhìn tôi: ``C-Cảm ơn Shino-san nhiều lắm! Cậu lúc nào cũng chu đáo như một người chị cả vậy\dots''

Cô bạn tóc xanh lau mồ hôi trên trán, nhìn quanh khu phố lễ hội đang vào độ náo nhiệt nhất, ánh mắt ánh lên sự tự hào rạng rỡ: ``Công nhận là lễ hội mùa hè năm nay đông hơn hẳn mọi năm. Mọi người ai cũng vui vẻ cả. Cảm giác bao nhiêu công sức thức đêm thức hôm chuẩn bị của lớp mình và Hội Học sinh đều hoàn toàn xứng đáng.''

``Đúng vậy!'' Ayame đập tay xuống bàn cái rầm, đứng phắt dậy. ``Công việc điều phối tạm thời coi như đã đi vào quỹ đạo ổn định. Cả năm người chúng ta đều có bộ đàm hết rồi đúng không? Bây giờ, với tư cách là Hội trưởng Hội Học sinh trường Trung học Aogami, tôi dõng dạc tuyên bố: Tổ công tác của chúng ta chính thức\dots\ chuyển sang chế độ giải lao!''

``Hả? Bây giờ sao?'' Nanase tròn mắt ngạc nhiên, vẻ mặt nghiêm túc thường ngày thoáng bối rối. ``Nhưng Ayame-san, lỡ có sự cố đột xuất\dots''

``Khéo lo!'' Ayame khoác tay lên vai cô, cười rạng rỡ kéo cô nàng đi. ``Lễ hội là để vui chơi mà! Chúng ta là học sinh, không phải robot bảo vệ hay camera an ninh. Đi thôi nào mấy đứa! Tôi nghe nói gian hàng crepe đằng kia mới nhập loại nhân kem matcha mới ngon tuyệt cú mèo! Lần này tôi bao chót!''

Thế là dưới sự dẫn dắt của ``Hội trưởng ham chơi'' Ayame, cả nhóm chúng tôi tạm gác lại sổ sách và hòa vào dòng người đông đúc. Wata dù miệng lẩm bẩm chê bai sự thiếu trách nhiệm của Hội trưởng nhưng vẫn bước đi sát ngay bên cạnh cô bạn thuở nhỏ. Nanase thì kéo tay Inari, liên tục chỉ trỏ những gian hàng lưu niệm lấp lánh.

Tôi đi chậm lại phía sau một chút, để gió vờn qua mái tóc, thu vào tầm mắt toàn bộ khung cảnh ấm áp này. Khi ngang qua khu vực ``Con đường Linh hồn'' - nhà ma phiên bản lễ hội thị trấn của Shion - tôi bất giác khựng lại.

Tiếng la hét, tiếng cười sợ hãi xen lẫn phấn khích vang ra từ bên trong. Nhưng có một âm thanh khiến tim tôi thắt lại: tiếng chuông trường. Tiếng chuông báo hết giờ quen thuộc đến ám ảnh, giống hệt tiếng chuông trong những vòng lặp mà tôi từng bị mắc kẹt. Shion đã phối hợp với Game để thiết kế nhà ma theo chủ đề ``trường học bị nguyền rủa'', và ông ta đã không ngần ngại xoáy sâu vào nỗi ám ảnh thật sự của thị trấn Aogami: những vòng lặp thời gian, những hành lang lặp đi lặp lại không có lối thoát, và tiếng chuông trường vang mãi không ngừng.

Hôm qua, khi được ba chị em Hikawa kể lại cuộc hành trình vào trong đó, tôi đã phải cắn chặt môi để không bật khóc. Bởi vì thế giới bên trong nhà ma ấy giống đến rợn người với ký ức của chính tôi: những hành lang dài hun hút, ánh đèn nhấp nháy, và cảm giác dù đi mãi vẫn quay về đúng nơi bắt đầu. Game biết. Ông ta luôn biết. Và ông ta đã cố tình tái hiện lại điều đó, không phải để hành hạ tôi, mà để nhắc nhở: ``Đó đã là quá khứ, không phải hiện tại.''

Đúng vậy. Đây là thế giới thật. Tôi đã thoát ra rồi.

Tôi hít một hơi thật sâu, bước qua khu vực nhà ma mà không ngoái lại.

``Shino-san! Lại đây nhanh lên! Cậu có muốn thử món này không?'' Ayame vẫy tay gọi lớn, trên tay đang cầm một chai thủy tinh có thiết kế kỳ lạ với một viên bi ve lấp lánh chặn ở cổ chai.

Tôi bước tới, tò mò nhìn món đồ uống có màu xanh lam sóng sánh bên trong: ``Đây là thứ gì vậy? Trông giống nước phép thuật ấy nhỉ.''

``Haha, Shino-san đùa buồn cười thật đấy! Là ramune đó! Nước chanh có ga siêu ngon luôn!'' Hội trưởng hớn hở dùng một nắp nhựa ấn mạnh viên bi xuống. Một tiếng *XỊT* vang lên, bọt khí sủi lên sùng sục trông vô cùng kỳ ảo. Cô ấy đưa chai nước cho tôi. ``Cậu uống thử đi! Trải nghiệm mùa hè đích thực đấy!''

Tôi đón lấy chai nước. Hơi lạnh từ lớp thủy tinh sần sùi truyền qua lớp da, khiến tôi khẽ rùng mình thích thú. Tôi đưa lên miệng nhấp một ngụm. Ngay lập tức, những bọt khí li ti nổ lách tách trên đầu lưỡi, mang theo vị chua ngọt thanh mát lan tỏa khắp khoang miệng, xua tan đi hoàn toàn cái nóng nực và oi ả của buổi trưa hè.

``Ngon chứ?'' Nanase đang cầm một chiếc bánh crepe dâu tây kem tươi, tò mò hỏi.

``Thật kỳ diệu\dots'' tôi khẽ cảm thán, ánh mắt không rời khỏi viên bi thủy tinh đang lăn lóc bên trong cổ chai tạo ra tiếng lanh canh mỗi khi tôi nghiêng nó. ``Trước đây, mình chưa từng nếm thử thứ gì mang hương vị sống động và bùng nổ đến nhường này. Dường như cả một mùa hè rực rỡ được gói gọn trong một chiếc bình thủy tinh nhỏ bé vậy.''

Inari bẽn lẽn mỉm cười, đôi mắt cong lên thành hình bán nguyệt: ``Nếu Shino-san thích món này, lát nữa mình sẽ mua thêm cho cậu một lốc đem về ký túc xá nhé.''

``Cảm ơn cậu, Inari-san. Mình ghi nhận tấm lòng của cậu.'' Tôi đáp lại bằng một nụ cười hiền hậu, đưa tay vuốt nhẹ mái tóc màu trắng muốt mềm mại của cô bạn nhút nhát ấy. Cậu ấy đỏ mặt, khẽ cúi đầu bối rối.

Chúng tôi tiếp tục dạo quanh khu lễ hội. Mọi mệt mỏi và áp lực của công việc điều phối an ninh dường như tan biến hết khi chúng tôi thực sự hòa mình vào các trò chơi. Ayame và Nanase thi nhau ném phi tiêu trúng bóng bay, kết quả là hai cô nàng cãi nhau giữa thanh thiên bạch nhật vì cùng ném trượt mục tiêu liên tiếp ba lần. Wata thì đứng bên cạnh, bình thản lấy máy tính bỏ túi ra tính toán xác suất phi tiêu bị lệch khiến ông chủ gian hàng đổ mồ hôi hột. Cũng may là Sakura không có ở đây, chứ không thì chắc lại thành ``cái chợ'' mất.

Trong khi đó, Inari thì mua được một chiếc mặt nạ hồ ly trắng tinh xảo, đội lệch sang một bên đầu trông vô cùng đáng yêu.

Khi mặt trời đã lên cao nhất, cái nóng của mùa hè bắt đầu ngột ngạt, nhóm chúng tôi dừng lại nghỉ chân bên dưới bóng râm của cổng torii khổng lồ.

Ayame dựa lưng vào cột gỗ đỏ chót, thở hắt ra một hơi dài mãn nguyện, đôi mắt ngắm nhìn dòng người bất tận: ``Vui thật đấy. Ước gì ngày nào chúng ta cũng có lễ hội như thế này để bên nhau.''

``Tham lam vừa thôi chứ,'' Wata cười khúc khích, uống nốt ngụm nước cuối cùng trong ly. ``Nhưng mình công nhận, đây là lễ hội tuyệt vời nhất, rực rỡ nhất từ trước đến nay của Aogami. Cảm giác như toàn bộ thị trấn, từng con đường, từng hàng cây đều đang được thanh lọc và bừng sáng vậy.''

Từ ``thanh lọc'' phát ra từ miệng Hội phó một cách vô tình khiến tôi khẽ giật mình. Ánh mắt tôi bất giác hướng về phía điện chính của đền Aogami, nơi những bóng cây cổ thụ đang rủ xuống tĩnh mịch, chuẩn bị cho một sự kiện trọng đại.

Chỉ còn vài tiếng đồng hồ nữa thôi.

Khi trăng rằm tháng bảy lên đến đỉnh đầu, năng lượng tâm linh của toàn bộ thị trấn này sẽ hội tụ lại ở mức cao nhất. Đó sẽ là khoảnh khắc Nghi thức Thanh tẩy diễn ra. Khác với những vòng lặp tàn khốc và vô tận trong quá khứ mà tôi từng bị giam cầm, nghi thức đêm nay là để ban phát sự sống, để neo giữ linh hồn mỏng manh của cô bé Naomi lại với thế giới vật lý này.

Và đồng thời\dots\ đó cũng là cánh cửa không gian mở ra để những ân nhân lớn nhất của đời tôi - gia đình Hikawa - rời đi vĩnh viễn. Hoặc ít nhất, chưa hẹn ngày tái ngộ.

Tôi đưa tay phải lên ngực trái, áp lòng bàn tay vào vị trí trái tim đang đập từng nhịp vững vàng và mạnh mẽ. Một sinh mệnh mới, một trái tim bằng xương bằng thịt mà Kuon-san đã trao lại cho tôi sau khi đánh bại thực thể hắc ám và giải phóng tôi khỏi xiềng xích của Định mệnh ngàn năm.

``Cậu đã hứa sẽ cho mình thấy một thế giới tươi đẹp, Kuon-kun. Và cậu đã giữ đúng lời hứa của mình. Mình sẽ sống trọn vẹn sinh mệnh này để ghi nhớ ân tình đó.''``''

Tôi hít một hơi thật sâu, thu lấy hương vị mằn mặn của gió biển, sự thanh tịnh của trầm hương, và nhiệt huyết của sự sống vào căng lồng ngực. Tôi quay sang nhìn những người bạn của mình - Nanase, Inari, Ayame, Wata. Họ không chỉ là những thành viên Hội Học sinh ưu tú, mà còn là những người bạn thực sự đầu tiên của tôi trong cuộc đời tự do này.

``Mọi người này,'' tôi cất giọng, dịu dàng nhưng đủ rõ ràng để thu hút sự chú ý của cả bốn người. ``Đêm nay sẽ là một đêm rất dài, và cũng rất quan trọng đối với lớp 1-C của chúng ta. Mọi thứ chúng ta đã chuẩn bị bằng mồ hôi và nước mắt, mọi nỗ lực không mệt mỏi của chúng ta, đều sẽ đơm hoa kết trái dưới ánh trăng rằm kia.''

Nanase gật đầu kiên định, ánh mắt cô nàng ánh lên ngọn lửa quyết tâm rực rỡ: ``Tất nhiên rồi. Bọn mình đã thề với lớp 1-C, hứa với nhà Hikawa rồi mà. Dù có chuyện gì xảy ra, chúng ta cũng sẽ bảo vệ lễ hội này, bảo vệ những nụ cười này đến giây phút cuối cùng!''

``Chuẩn luôn! Hội Học sinh trường Trung học Aogami không bao giờ chùn bước trước bất kỳ thử thách nào!'' Ayame đấm hai tay vào nhau đầy khí thế.

Tôi mỉm cười rạng rỡ, một nụ cười trong trẻo và tràn đầy sức sống, nụ cười của một cô gái đã thoát ra khỏi vòng lặp và đang sống trọn vẹn từng khoảnh khắc của cuộc đời mới.

``Vậy thì, mọi người ơi,'' tôi dang rộng hai tay, để cơn gió đầu hạ hắt lên tà áo yukata trắng muốt, ``hãy cùng nhau tận hưởng trọn vẹn từng khoảnh khắc của đêm nay. Vì thị trấn Aogami, vì sự tái sinh, và vì những người mà chúng ta trân quý sâu sắc nhất!''

\pov{Hikawa Kuon}

Trái ngược hoàn toàn với sự ồn ào và náo nhiệt ở khu vực trung tâm lễ hội, con đường dẫn lên đền Aogami từ phía sườn đồi phía Đông lại tĩnh lặng và rợp bóng cây xanh mát. Dưới cái nắng oi ả của buổi trưa hè, những tán lá phong và bách tùng đan vào nhau tạo thành một mái vòm tự nhiên, lọc lấy những tia nắng chói chang thành những đốm sáng nhảy múa trên bậc thang đá rêu phong.

Từ sáng đến giờ, trong khi bốn nhóm còn lại đang quậy tung từng ngóc ngách của thị trấn, tôi lại chọn cho mình một lịch trình có phần\dots\ ``dưỡng lão'' hơn: đi dạo quanh các gian hàng thủ công truyền thống và vãn cảnh cùng nhóm của Yuko-sensei.

Đi cùng chúng tôi còn có Hotaru và Mitsuki. Hai cô nàng vốn thích sự yên tĩnh, nên việc họ chọn đi cùng giáo viên chủ nhiệm thay vì hòa vào đám đông ồn ào ngoài kia cũng là điều dễ hiểu. Lúc này, cô nàng da ngăm đang đứng ngẩn ngơ trước một quầy bán chuông gió furin làm bằng thủy tinh, những ngón tay thon dài khẽ chạm vào tờ giấy tanzaku đung đưa theo chiều gió. Cô nàng tóc tím-hồng đứng cạnh, mỉm cười dịu dàng chỉ cho cô bạn thân xem những họa tiết hoa bỉ ngạn được vẽ tay tỉ mỉ trên chiếc chuông.

``Này, Kuon-san.''

Giọng nói trưởng thành nhưng mang theo nét trêu chọc quen thuộc vang lên bên cạnh tôi. Tôi quay sang. Yuko-sensei đang đứng đó, khoác trên mình bộ yukata màu tím than với họa tiết pháo hoa ánh vàng, trên tay cầm một chiếc quạt lụa phe phẩy nhè nhẹ. Khác với bộ trang phục công sở cứng nhắc thường ngày, cô giáo chủ nhiệm lớp 1-C hôm nay trông đằm thắm và có sức hút của một người phụ nữ trưởng thành.

``Vâng, thưa cô?'' tôi đáp, khẽ gật đầu.

Yuko-sensei hơi nghiêng đầu, đôi mắt màu thạch anh khẽ híp lại đằng sau tròng kính: ``Từ sáng tới giờ tôi cứ thắc mắc mãi. Tại sao em lại chọn đi cùng nhóm với tôi và hai em ấy? Tôi cứ nghĩ ở cái tuổi rực rỡ nhất của thanh xuân này, em phải đang chạy bùng nổ cùng với mấy cô bạn ồn ào ngoài kia chứ?''

Tôi bật cười, đưa tay gãi gáy: ``Cô đánh giá cao thể lực và sức chịu đựng của em quá rồi đấy, Yuko-sensei.''

Cô gấp chiếc quạt lụa lại cái *CẠCH*, vỗ nhẹ lên lòng bàn tay: ``Ồ? Vậy em thử giải thích cho tôi nghe xem nào? Tại sao lại bỏ mặc dàn hậu cung\dots\ à nhầm, dàn bạn bè năng nổ của mình để đi vãn cảnh với một 'bà cô già' như tôi?''

``Cô mà già thì chắc trên đời này chẳng có ai trẻ cả,'' tôi buột miệng đùa lại một câu, khiến Yuko-sensei phải bật cười thích thú. Sau đó, tôi hắng giọng, bắt đầu đếm trên đầu ngón tay những lý do rất đỗi hợp lý của mình.

``Đầu tiên, nếu em đi cùng với nhóm của Suzuki. Cô cũng biết dạ dày của em ấy giống như một cái hố đen vũ trụ mà. Thêm cả Suzune-san với cái tính hiếu thắng nữa. Nếu em ở đó, em không chỉ phải đóng vai 'ví tiền di động' để thanh toán hóa đơn cho hàng núi đồ ăn, mà rất có thể em còn bị kéo vào làm trọng tài cho một cuộc thi ăn thịt nướng không hồi kết. Em không muốn chết vì kiệt sức đâu ạ.''

Yuko-sensei gật gù, khóe môi giật giật: ``Có lý. Tôi có thể tưởng tượng cảnh em đứng khóc bên đống hóa đơn.''

``Thứ hai,'' tôi giơ ngón tay tiếp theo, ``nhóm của Kaigou-chan. Em biết cô ấy từng là một xạ thủ và là người chuyên về mảng chiến đấu.Vào khu trò chơi, nhất là gian hàng bắn súng, Kaigou-chan sẽ không chơi để cho vui, mà cô ấy sẽ chơi như một sát thủ đi làm nhiệm vụ. Nếu em đi cùng, nhiệm vụ duy nhất của em sẽ là đi theo xin lỗi các chủ gian hàng vì lỡ quét sạch giải thưởng của họ, hoặc ngăn không cho cô ấy dùng kỹ thuật quân sự để\dots\ ném vòng cổ chai.''

``Cái này thì tôi phải công nhận,'' Yuko-sensei thở dài thườn thượt. ``Sáng nay tôi vừa nghe loáng thoáng trên bộ đàm của Ban Kỷ luật báo về vụ lùm xùm ở quầy shateki. Em chọn tránh xa là một quyết định sinh tồn rất sáng suốt.''

``Thứ ba, nhóm của Naomi,'' tôi nhìn về phía tháp chuông xa xa. ``Naomi đang trong giai đoạn tìm hiểu thế giới, nên cũng chẳng khác gì trẻ con. Đi cùng em ấy sẽ giống như một ông bố trẻ đi trông đứa con gái ở một khu vui chơi đông đúc vậy. Em sẽ phải trả lời một vạn câu hỏi 'Vì sao?' của Naomi và đôi lúc sẽ khiến em muốn phát mệt. Sự kiên nhẫn của em cũng có giới hạn mà. Với cả\dots'' tôi ngập ngừng một lúc lâu.

``Với cả?'' Yuko-sensei nghiêng đầu.

``\dots Không giấu gì cô, em cũng không thích trẻ con đâu ạ. Đừng hiểu nhầm em, Naomi là một cô bé ngoan và không nghịch ngợm, nhưng đôi lúc\dots\ em ấy cũng khiến em phải mệt bở hơi tai vì sự tò mò đấy ạ.''

Cô giáo chủ nhiệm bật cười thành tiếng, đưa tay che miệng: ``Ôi trời, em miêu tả cứ y như người từng trải vậy. Thế còn nhóm Hội Học sinh của Shino-san thì sao?''

``Lại càng không,'' tôi rùng mình ngay lập tức. ``Cô nghĩ Ayame-san sẽ để yên cho em thảnh thơi đi chơi sao? Chắc chắn cậu ấy sẽ nhét vào tay em một cái băng đeo tay, một cái bộ đàm, và lôi em đi phân luồng giao thông hoặc đứng gác ở cổng đền. Em đến đây để tham gia lễ hội, chứ không phải đi làm thêm không công cho Hội Học sinh đâu ạ.''

Nói xong, tôi thở hắt ra một hơi, dựa lưng vào lan can gỗ bên đường. ``Nên cuối cùng, sau khi dùng phương pháp loại trừ, em thấy đi cùng cô, Hotaru và Mitsuki là lựa chọn an toàn và yên bình nhất. Em chỉ muốn có một khoảng thời gian trầm lắng để thực sự ngắm nhìn thị trấn này, ngắm nhìn những thứ mang đậm dấu ấn của thời gian và văn hóa, thay vì cứ chạy quanh như một con quay.''

Yuko-sensei im lặng một chút, ánh mắt cô dịu đi rất nhiều. Cô mở quạt, phe phẩy tạo ra những luồng gió mát mẻ xua tan cái oi bức.

``Em luôn như vậy nhỉ, Kuon-san,'' cô khẽ nói, giọng mang theo một sự thấu hiểu sâu sắc. ``Luôn biết cách quan sát mọi thứ từ xa, luôn chọn đứng ở một vị trí có thể bao quát toàn bộ, và luôn mang trong mình những suy nghĩ trưởng thành hơn tuổi.''

``Chắc do em già trước tuổi thôi ạ,'' tôi cười trừ. Dù gì tôi cũng hơn các bạn học ở đây tận mười tuổi mà, nhưng tôi quyết định không thốt ra.

``Không đâu,'' cô lắc đầu. ``Đó là sự từng trải. Đôi khi tôi quên mất rằng em và hai người em gái của mình đến từ một thế giới khác, một nơi có lẽ không bình yên như thị trấn nhỏ bé này.''

Cô bước lên một bậc thang, đứng ngang hàng với tôi, ánh mắt hướng về phía đường chân trời phía xa, nơi mặt biển xanh ngắt đang giao hòa với bầu trời.

``Tôi không biết quá khứ của các em ra sao, cũng không biết những khó khăn mà các em đã phải đối mặt ở thế giới kia. Nhưng với tư cách là giáo viên chủ nhiệm của em ở thế giới này\dots\ tôi thực sự rất vui vì các em đã đến đây.'' Yuko-sensei mỉm cười, một nụ cười hiền hậu và ấm áp, không còn vẻ nghiêm khắc hay trêu đùa nữa. ``Aogami vốn là một thị trấn khép kín và mang đầy những nỗi buồn từ quá khứ. Nhưng chính sự xuất hiện của các em đã mang lại một luồng sinh khí mới, đã phá vỡ những xiềng xích vô hình, và cứu rỗi những linh hồn vốn tưởng chừng như đã rơi vào tuyệt vọng.''

Tôi khẽ cúi đầu: ``Chúng em chỉ làm những gì cần phải làm thôi ạ. Bọn em cũng nợ thị trấn này, nợ lớp 1-C rất nhiều.''

``Kuon-kun này,'' Hotaru đột nhiên lên tiếng. Cô nàng da ngăm đã bước tới từ lúc nào, trên tay cầm hai chiếc chuông gió furin bằng thủy tinh. Một chiếc vẽ hình hoa anh đào, chiếc còn lại vẽ hoa bỉ ngạn trắng.

Mitsuki đứng ngay sau lưng cô bạn, mỉm cười dịu dàng: ``Bọn mình định mua cái này để làm kỷ niệm. Kuon-kun thấy sao?''

``Đẹp lắm,'' tôi thật lòng khen ngợi. ``Tiếng chuông thủy tinh sẽ giúp xua đi cái nóng của mùa hè. Rất hợp với hai cậu đấy.''

Hotaru ngập ngừng một giây, rồi chìa chiếc chuông gió họa tiết anh đào về phía tôi. Nét mặt của một người thích tám chuyện sôi nổi thoáng xuất hiện một vành đỏ nhạt vì ngượng: ``Cái này\dots\ tôi tặng cậu.''

``Tặng tôi ư?'' tôi ngạc nhiên.

``Ừm,'' cô gật đầu nhẹ, mắt tránh nhìn thẳng vào tôi. ``Coi như\dots\ một lời cảm ơn. Vì cậu đã luôn ở bên cạnh, bảo vệ mọi người, và\dots\ mang đến cho lớp 1-C một luồng gió mới.''

Mitsuki cũng hùa theo, cười khúc khích: ``Nhận đi Kuon-kun, Hotaru-chan đã đứng chọn suốt nửa tiếng đồng hồ chỉ để tìm được chiếc ưng ý nhất cho cậu đấy na.''

``M-Mitsuki-chan!'' Hotaru giật mình, lườm cô bạn thân, nhưng vành tai đã đỏ bừng lên.

Tôi bật cười, đưa tay nhận lấy chiếc chuông gió. Sợi chỉ đỏ mềm mại trượt qua kẽ tay, viên bi thủy tinh chạm vào thành chuông tạo nên một tiếng lanh canh trong trẻo, vang vọng giữa không gian tĩnh lặng của khu rừng đền.

``Cảm ơn hai cậu. Mình sẽ treo nó ở một nơi thật đẹp,'' tôi nói, cảm nhận được một luồng hơi ấm nhẹ nhàng lan tỏa trong lồng ngực.

Yuko-sensei đứng cạnh, khẽ lấy quạt che đi nụ cười đầy ẩn ý: ``Ôi chà chà, thanh xuân vườn trường rực rỡ quá đi mất. Đứng cạnh hai cô gái ngoan hiền thế này, bảo sao cậu nhóc nhà ta lại chẳng muốn đi theo mấy 'bà chằn' kia.''

``Cô lại bắt đầu rồi đấy\dots'' tôi thở dài bất lực.

Nhưng đúng lúc này, tôi ngước nhìn lên cao. Bóng của cổng torii phía trước đền đã bắt đầu ngả dài trên những bậc thang đá. Mặt trời đang dần ngả về Tây, báo hiệu buổi chiều tà đang đến gần.

Không khí náo nhiệt của buổi sáng và đầu giờ chiều vẫn tiếp diễn, nhưng trong sâu thẳm, tôi có thể cảm nhận được một sự thay đổi tinh tế trong bầu không khí. Năng lượng tâm linh của ngọn núi đang dần thức tỉnh, dày đặc hơn, tĩnh lặng hơn.

Trái với sự rực rỡ của lễ hội mùa hè, phần tiếp theo của ngày hôm nay sẽ mang một màu sắc hoàn toàn khác. Khi mặt trời lặn xuống, nghi thức sẽ bắt đầu. Những lời hứa sẽ được thực hiện, và những cánh cửa không gian sẽ mở ra.

Tôi siết chặt chiếc chuông gió trong tay. Chỉ còn vài tiếng nữa thôi.

``Thưa cô,'' tôi cất giọng, quay sang nhìn Yuko-sensei. ``Chúng ta cũng nên chuẩn bị quay lại khu trung tâm thôi. Lát nữa, phần triển lãm các truyền thuyết thần thoại và nghi thức viết lời ước sẽ bắt đầu rồi.''

Yuko-sensei gật đầu, nét mặt trở lại vẻ nghiêm túc nhưng vẫn giữ nét dịu dàng: ``Em nói đúng. Mọi người đang đợi chúng ta đấy. Đi thôi.''

Dưới ánh nắng vàng rực rỡ của buổi chiều tà, bốn người chúng tôi cùng nhau bước xuống những bậc thang đá, hướng về phía trung tâm thị trấn đang vang lên những điệu nhạc dân gian rộn ràng. Nơi đó, những người bạn của tôi, những thành viên của lớp 1-C, đang chờ đợi để cùng nhau trải qua những khoảnh khắc cuối cùng của một lễ hội mùa hè không bao giờ quên.

\ct

\begin{center}
  \uline{Khu Triển lãm Văn hóa Thị trấn Aogami.\\16:30.}
\end{center}

Khi mặt trời bắt đầu ngả bóng dài trên những con dốc lát đá, nhóm của tôi cũng vừa vặn gặp lại những người khác tại sảnh đền phụ. Trái với vẻ tơi tả mà tôi dự đoán, cả Kaigou-chan và Suzuki đều trông khá tràn trề năng lượng. Cô chị lớn vác theo một túi to oành chứa đầy những chiến lợi phẩm từ quầy bắn súng, còn cô em gái thì trên tay vẫn cầm một xiên dango đang cắn dở. Nhóm của Naomi cũng đã đến nơi, cô bé nhỏ nhắn đang tò mò nhìn quanh, đôi mắt nhấp nháy háo hức khi nhìn thấy đám đông. Thậm chí cả hội của Ayame cũng đã có mặt đông đủ, có lẽ họ đã kịp bàn giao lại công việc cho nhóm ca chiều. Touka cũng đã có mặt ở đây, có lẽ cũng muốn tìm hiểu về nguồn cội về nơi mình đang sống.

Điểm đến tiếp theo của chúng tôi là Khu Triển lãm Văn hóa, được dựng lên ngay tại khu vực sảnh lớn phía trước đền Aogami. Đây là một hoạt động mới được thêm vào năm nay, nhằm tôn vinh những giá trị lịch sử và vẻ đẹp của thị trấn.

Vừa bước qua cánh cửa gỗ lớn, chúng tôi đã ngay lập tức bị choáng ngợp bởi một không gian được bài trí vô cùng trang trọng. Những bức rèm vải noren thêu họa tiết cổ xưa rủ xuống từ trần nhà, mùi hương trầm thoang thoảng tạo cảm giác linh thiêng và huyền bí.

``Chào mừng mọi người đến với Triển lãm Truyền thuyết Thần thoại Aogami!''

Một giọng nói lanh lảnh vang lên. Sakura đứng đó, khoác trên mình bộ trang phục vu nữ truyền thống với áo trắng quần đỏ, mái tóc hồng đào được búi gọn gàng. Hôm nay, cô ấy không chỉ đến tham gia lễ hội, mà còn đảm nhận vai trò 'hướng dẫn viên' cho khu vực này.

``Oa\~{} Sakura-san, cậu mặc bộ này trông hợp lắm đấy!'' Naomi trầm trồ, hai mắt sáng rực lên.

``Đương nhiên rồi! Đây là trang phục chuẩn của một Vu nữ ưu tú mà lị'' Sakura ưỡn ngực tự hào, rồi cầm một cây quạt giấy chỉ về phía những bức bình phong lớn được xếp dọc theo lối đi. ``Nào, hãy để tôi kể cho mọi người nghe về nguồn gốc của thị trấn này.''

``Vu nữ ưu tú'' à? Nghe thoang thoảng có gì đó tự cao ở đây\dots\ nhưng ít ra thì Sakura còn xứng đáng với cái danh ấy đó. Chẳng bù cho Miko-chi ở thế giới thực, nơi cô nàng tự nhận mình cái danh ``Ưu tú'' nhưng đi tới đâu thì bung bét tới đó.

Chúng tôi đi theo bước chân của cô ấy. Trên những bức bình phong là những bức họa cổ, miêu tả lại cảnh một ngọn núi phun trào, một trận đại hồng thủy, và hình ảnh một con rồng nước uốn lượn bảo vệ đất đai.

``Theo truyền thuyết, Aogami từng là nơi giao thoa giữa cõi người và cõi thần. Ngọn núi thiêng phía sau đền chính là nơi trú ngụ của những linh hồn cổ xưa,'' Sakura hắng giọng, bắt đầu thuyết minh với phong thái chuyên nghiệp đến bất ngờ, giọng nói vang lên dõng dạc. ``Hàng trăm năm trước, một trận đại tai ương đã giáng xuống. Khi đó, một vị Vu nữ đã hy sinh thân mình, lập ra khế ước với Thần linh để tạo ra kết giới bảo vệ thị trấn. Từ đó, Nghi thức Thanh tẩy được ra đời để duy trì sức mạnh của kết giới ấy.''

Touka chớp chớp mắt, ngước lên nhìn bức tranh, ngón tay khẽ chạm cằm: ``Dữ liệu lịch sử của Aogami không có ghi chép cụ thể về 'trận đại tai ương' này. Phải chăng đây là sự kiện sạt lở đất từng được giấu kín, hay là một hiện tượng siêu nhiên chưa được phân tích theo góc độ khoa học?''

Kaigou-chan khoanh tay, khẽ nhếch mép cười nhạt: ``Truyền thuyết luôn được viết lại để che đậy một sự thật nào đó mờ ám hơn. Những thảm họa tự nhiên thường được thần thánh hóa để con người có một chỗ dựa tinh thần. Nhưng không thể phủ nhận, cách người ta vẽ lại nó trông khá hoành tráng đấy.''

``Hai người không thể lãng mạn hóa nó lên một chút được à? Người ta đang cố gắng tạo không khí linh thiêng mà!'' Sakura phồng má dỗi, vung vẩy cây quạt giấy, rồi lại chỉ sang bức tranh tiếp theo. ``Nghi thức Thanh tẩy đêm nay chính là sự mô phỏng lại khế ước năm xưa. Mọi người có thấy những chiếc lồng đèn không? Chúng đại diện cho linh hồn của những người đã khuất, được thắp sáng để dẫn đường cho họ về với suối vàng, đồng thời mang lại sự bình yên cho người sống.''

Tôi đứng lặng yên, nhìn chằm chằm vào bức họa vẽ vị Vu nữ cô độc đứng giữa vòng tròn phép. Ánh mắt tôi vô tình chạm phải Shino-san, người cũng đang nhìn bức tranh ấy với một nụ cười thoáng buồn nhưng thanh thản. Khác với những người dân thị trấn chỉ xem đây là một câu chuyện cổ tích, chúng tôi biết rõ sự thật. Chúng tôi biết vị Vu nữ ấy thực sự là ai, những mảnh vỡ linh hồn đã trôi dạt ra sao, và chúng tôi biết vòng lặp chết chóc ấy đã tàn nhẫn đến mức nào. Nhưng giờ đây, những điều đó chỉ còn là 'truyền thuyết' được trưng bày trên giấy, không còn sức mạnh để làm hại bất kỳ ai nữa.

``Thật kỳ lạ khi thấy những thứ mình từng trải qua giờ lại trở thành một câu chuyện để kể cho du khách nghe,'' tôi thì thầm, giọng nói trầm xuống, đủ để Shino-san và hai người đồng hành nghe thấy.

``Nhưng như thế lại tốt,'' Kaigou-chan đáp khẽ, tay vuốt nhẹ một nếp gấp trên áo. ``Ký ức đau buồn nên được đóng khung lại, nhường chỗ cho tương lai xán lạn hơn.''

``Đúng vậy,'' Shino-san khẽ mỉm cười, đôi mắt xanh lấp lánh phản chiếu ánh đèn lồng dịu nhẹ. ``Truyền thuyết là để kể, còn cuộc đời là để sống. Và chúng ta\dots\ đang sống rất tốt ở hiện tại rồi. Mình cảm thấy ổn với điều đó.''

\ct

Sau khi đi hết khu vực Thần thoại, chúng tôi bước qua một vòm lá phong giả để tiến vào nửa sau của sảnh lớn: Khu vực Triển lãm Vẻ đẹp Thị trấn. Khác với không khí huyền bí ban nãy, nơi này tràn ngập ánh sáng rực rỡ và những màu sắc tươi tắn của hiện tại. Nơi đây không kể về những vị thần, mà kể về con người.

Những tấm ván lớn được xếp ngay ngắn, trên đó trưng bày hàng chục bức ảnh nhiếp ảnh và các tác phẩm hội họa đủ mọi kích cỡ.

``A, Kuon-san! Mọi người tới rồi!''

Uzuki bước ra từ sau một tấm pano lớn, trên cổ vẫn đeo chiếc máy ảnh DSLR quen thuộc. Đi cùng cô là Mari, cô bạn với đôi mắt hai màu đỏ-vàng nổi bật, tay đang cầm một bảng màu nhỏ xíu như một thói quen khó bỏ.

``Chà, hoành tráng thật đấy,'' Suzuki tiến lại gần một bức ảnh lớn, trầm trồ. Bức ảnh chụp lại khoảnh khắc bình minh trên biển Aogami, khi những tia nắng đầu tiên xuyên qua làn sương mù mờ ảo, nhuộm đỏ cả những mái nhà ngói cổ và những con thuyền đánh cá đang neo đậu phía xa.

``Đây đều là những bức ảnh cậu chụp sao, Uzuki-san?'' tôi hỏi, đảo mắt qua một vòng các tác phẩm. Không chỉ có phong cảnh, mà còn có những bức ảnh chụp nụ cười của các cụ già ở chợ sáng, những đứa trẻ đang chơi đùa trước sân đền, và cả khoảnh khắc lớp 1-C chúng tôi đang hì hục dựng gian hàng ngày hôm qua. Mọi thứ hiện lên sống động đến mức tôi có thể nghe thấy cả tiếng cười qua khung ảnh.

Cô bạn bờm Thỏ gật đầu, nở một nụ cười rạng rỡ và tự hào, hai tay chống nạnh: ``Đúng vậy! Bọn mình không chỉ muốn lưu giữ những kỷ niệm trong trường học, mà còn muốn cho mọi người thấy Aogami đẹp đến nhường nào. Mỗi góc phố, mỗi con người ở đây đều mang một câu chuyện riêng.''

``Oa\~{} Hình đẹp quá!'' Naomi thích thú thốt lên, kiễng chân lên và dùng ngón tay nhỏ xíu chỉ vào một bức ảnh chụp những dải ruy băng ngũ sắc tung bay trong gió. Cô bé mở to đôi mắt tròn xoe, ngước nhìn nhiếp ảnh gia của lớp: ``Uzuki-oneesan như có phép thuật ấy! Dù chỉ là tờ giấy thôi, nhưng em cảm giác như những dải ruy băng này sắp bay ra ngoài và chạm vào mũi em vậy!''

``Đó chính là 'hồn' của nhiếp ảnh đấy, Naomi-chan,'' Uzuki cười tít mắt, vỗ nhẹ lên mái tóc tết gọn gàng của cô bé tám tuổi.

``Và nhiếp ảnh không phải là thứ duy nhất có hồn đâu nhé!'' Mari lên tiếng, kiêu hãnh kéo chúng tôi về phía góc trưng bày của cô ấy.

Trái ngược với sự chân thực của những bức ảnh, các tác phẩm của cô ấy mang một vẻ đẹp mềm mại và bay bổng hơn của màu nước và sơn dầu. Bức tranh lớn nhất đặt ở chính giữa thu hút sự chú ý của tất cả chúng tôi. Đó là một bức tranh sơn dầu vẽ lại cảnh lớp 1-C đang đứng dưới gốc cây anh đào Shinomiya. Màu sắc được phối vô cùng ấm áp, từng cánh hoa anh đào rơi xuống đều mang theo một sắc thái riêng biệt.

``Đẹp quá\dots'' Inari, người vừa bước tới từ nhóm Hội Học sinh, thốt lên, đôi mắt mở to kinh ngạc.

``Cậu vẽ mất bao lâu vậy, Mari-san?'' Kaigou-chan cũng phải nheo mắt nhìn những chi tiết đánh bóng tỉ mỉ trên bức tranh. ``Cách phối màu này không phải là việc có thể làm trong ngày một ngày hai được.''

Mari vuốt lại bím tóc, cười tự mãn nhưng không giấu được sự ngượng ngùng, hai vành tai cô hơi đỏ lên: ``Cũng mất gần một tháng đấy. Mình muốn ghi lại khoảnh khắc chúng ta sát cánh bên nhau. Aogami không chỉ đẹp ở cảnh vật, mà vẻ đẹp lớn nhất của nó chính là con người\dots\ là tất cả chúng ta.''

Tôi nhìn bức tranh, lại nhìn sang những người bạn đang đứng quanh mình. Uzuki với nụ cười tỏa nắng, Mari với sự tinh tế của một nghệ sĩ, Sakura đang chống cằm suy tư, Shino với ánh mắt dịu dàng, nhóm Hội Học sinh vẫn đang chí chóe cãi nhau về một bức ảnh, và cả Kaigou, Suzuki, Naomi. Mọi người đều đang mỉm cười, đều đang hiện diện ở đây, sống động và chân thực hơn bất kỳ truyền thuyết hay tác phẩm nghệ thuật nào.

``Phải,'' tôi khẽ mỉm cười, cảm nhận được sự ấm áp lan tỏa trong lồng ngực. ``Aogami thực sự rất đẹp. Một vẻ đẹp mà dù có đi đâu xa, chúng ta cũng sẽ không bao giờ quên được.''


\ct

Sau khi rời khỏi Khu Triển lãm, cả nhóm di chuyển đến khoảng sân phía trước đền, nơi đặt cây trúc lớn dành cho phần Viết lời ước. Đây là một truyền thống không thể thiếu trong các lễ hội mùa hè, nhưng năm nay, dưới sự đề xuất của Sakura, quy mô của nó đã được mở rộng để tất cả học sinh và người dân thị trấn đều có thể tham gia.

Vừa bước tới nơi, tôi đã thấy một bóng người mặc áo choàng đen - trang phục hoàn toàn lạc quẻ giữa một rừng yukata rực rỡ - đang đứng đợi sẵn dưới gốc cây. Là Shion.

Thấy tôi, cô nàng bước tới, khuôn mặt giấu dưới vành mũ rộng tỏ vẻ hậm hực. Cô đưa tay ra, chìa trả lại cho tôi chiếc đồng hồ đeo tay quen thuộc.

``Trả cho ngươi này,'' Shion lầm bầm, giọng điệu nghe như một đứa trẻ bị bắt phải trả lại đồ chơi yêu thích. ``Nhà ma của ta đã đạt đủ chỉ tiêu doanh thu rồi. Dù ta thừa sức dọa cho bọn họ khóc thét thêm một tuần nữa, nhưng ông chú hắc ám này bảo đến giờ phải về với chủ nhân rồi.''

Tôi bật cười, nhận lấy chiếc đồng hồ đeo vào tay. Ngay lập tức, một giọng nói trầm khàn vang lên trong chiếc đồng hồ: ``Cô ta bóc lột sức lao động của tôi đến mức vắt kiệt từng chút năng lượng thiết kế ảo ảnh. May mà tôi không có thực thể vật lý, nếu không chắc tôi gãy lưng mất, Kuon-user ạ.''

``Ông vất vả rồi, Game,'' tôi thì thầm đáp lại.

Shion hừ lạnh một tiếng, khoanh tay quay mặt đi: ``Ta có nghe thấy đấy nhé. Lần sau có tổ chức sự kiện gì thì nhớ cho ta mượn lại ổng đấy.''

Tôi chỉ biết cười trừ. Lúc này, Sakura đã mang đến một xấp giấy tanzaku ngũ sắc và một hộp bút lông, phân phát cho từng người.

``Nào mọi người! Khấn nguyện với Thần bảo hộ đi!'' Sakura vỗ tay hớn hở. ``Viết điều ước của mình lên đây rồi treo lên cành trúc nhé! Đảm bảo vô cùng linh nghiệm!''

``Cậu tự phong mình là người nhận điều ước luôn à?'' Kaigou nhướng mày.

``Thì tôi là Thần bảo hộ mà lị!'' Sakura ưỡn ngực tự hào.

Mọi người tản ra xung quanh, tìm cho mình một góc để nắn nót viết những dòng tâm sự. Tôi đứng tựa lưng vào gốc trúc, lặng lẽ quan sát. Chỉ một lát sau, những dải giấy đủ màu sắc đã bay phấp phới trên cành. Tôi bước đến gần, đưa mắt đọc lướt qua những dòng chữ trên đó.

Dễ dàng nhận thấy, những lời ước được chia thành ba luồng cảm xúc rõ rệt.

Đầu tiên là những người chưa hề hay biết về bí mật của gia đình Hikawa - những người tin rằng ngày mai mặt trời vẫn sẽ mọc và chúng tôi vẫn sẽ đến trường cùng nhau.

\txtbx{
  Mình muốn chụp thêm thật nhiều, thật nhiều ảnh của lớp 1-C! Mong rằng nụ cười của mọi người sẽ mãi rạng rỡ như ngày hôm nay!
  \flushright - Uzuki
}

\txtbx{
  Nghệ thuật là ánh trăng lừa dối, nhưng tình bạn của chúng ta là nét cọ chân thực nhất. Ước gì tớ có đủ thời gian để vẽ chân dung cho từng người một.
  \flushright - Mari
}

\txtbx{
  Mình hy vọng sẽ tiếp tục phấn đấu để luôn đứng đầu trong năm học tới!
  \flushright - Touka
}

\txtbx{
  Nhà ma sang năm nhất định phải mở rộng gấp ba lần! Và phải có thêm hiệu ứng 4D để mọi người sợ khóc thét luôn!
  \flushright - Shion
}

Tôi mỉm cười nhẹ. Thật hồn nhiên và tràn đầy hy vọng về một tương lai bình dị. Tôi tiếp tục bước dọc theo cành trúc, lần lượt đọc từng điều ước của những người bạn cùng lớp còn lại.

\txtbx{
  Năm sau nhất định phải thuyết phục ban tổ chức bắn pháo hoa to gấp đôi, xịn gấp ba năm nay!
  \flushright - Hanabi
}

\txtbx{
  Sẽ phục thù và đánh bại Suzuki trong cuộc thi ăn thịt nướng lần tới! Quyết không bỏ cuộc!
  \flushright - Suzune
}

\txtbx{
  Mong sao những đóa hoa mà mình chăm sẽ đua nở rực rỡ như thế này vào năm sau.
  \flushright - Saki
}

\txtbx{
  Mong ngày nào cũng được chọc phá mọi người thoả thích mà không bị phạt!
  \flushright - Honoka
}

\txtbx{
  Cầu chúc cả lớp 1-C luôn mạnh khỏe, bình an và ai cũng giữ được vóc dáng chuẩn mực.
  \flushright - Yae
}

Tôi khẽ bật cười. Hanabi vẫn luôn đam mê pháo hoa như vậy, còn Suzune thì có vẻ sự ganh đua một chiều với cái "hố đen vũ trụ" Suzuki đã trở thành mục tiêu sống mất rồi. Sự năng động và đầy sức sống của họ khiến tôi cảm thấy những mệt mỏi của một ngày dài dường như tan biến. Lại bước thêm vài bước, tôi đọc tiếp những dải giấy khác.

\txtbx{
  Mong là không phải đối đầu thêm mấy thằng thích nộp mạng cho Shiina này nữa.
  \flushright - Shiina
}

\txtbx{
  Mình sẽ rèn luyện để trở thành thần tượng tỏa sáng nhất Aogami!
  \flushright - Yuka
}

\txtbx{
  Sẽ luôn rực cháy hết mình trong mọi trận đấu! Lớp 1-C cố lên!
  \flushright - Akane
}

\txtbx{
  Mong Aogami ngày càng thu hút nhiều du khách để mình có cơ hội làm hướng dẫn viên nhiều hơn!
  \flushright - Kaoru
}

\txtbx{
  Mình cũng ước giống Kaoru-chan, và mong cả lớp sẽ mãi thân thiết thế này.
  \flushright - Saya
}

Tuổi trẻ thật tốt. Họ có những đam mê, những định hướng rõ ràng và dồn hết nhiệt huyết vào đó. Bất kể là thể thao, nghệ thuật, hay chỉ đơn giản là mong muốn làm đẹp cho quê hương, những ước mơ này đều lấp lánh như những vì sao nhỏ vớt vát lại ánh sáng cho thị trấn này.

\txtbx{
  Mình mong sau này sẽ trở thành một người trưởng thành dày dặn.
  \flushright - Miku
}

\txtbx{
  Mong rằng tập thể lớp 1-C sẽ tiếp tục phát huy!
  \flushright - Kana
}

\txtbx{
  Buổi biểu diễn của CLB Âm nhạc tối nay sẽ là màn trình diễn tuyệt vời nhất từ trước đến nay!
  \flushright - Kanade
}

\txtbx{
  Mong tiếng chuông gió sẽ xua tan đi mọi muộn phiền của những người mình trân quý.
  \flushright - Hotaru
}

\txtbx{
  Ước gì Hotaru-chan sẽ cười nhiều hơn nữa. Nụ cười của cậu ấy rất đẹp.
  \flushright - Mitsuki
}

Nhưng rồi sự ấm áp lập tức lấp đầy lồng ngực khi tôi đọc được những dòng chữ nắn nót của Hotaru và Mitsuki. Hai cô bạn trầm tính này luôn biết cách quan tâm đến người khác bằng những điều tinh tế nhất.

\txtbx{
  Mong Koishin bớt tính khí nắng mưa lại để thành thật với bản thân hơn một chút! (P/s: Chúc cậu sớm có người rước).
  \flushright - Yuki
}

\txtbx{
  Luôn tràn trề năng lượng và vui vẻ mỗi ngày! Không bao giờ biết mệt!
  \flushright - Azuki
}

\txtbx{
  Chỉ mong năm nay không gặp xui xẻo hay rước họa vào thân nữa. Và\dots\ ừm, lớp 1-C là tuyệt nhất. Mọi người đừng có mà đọc trộm!
  \flushright - Koishin
}

\txtbx{
  Sẽ phá đảo mọi kỷ lục game bắn súng shateki ở tất cả các lễ hội mùa hè Aogami!
  \flushright - Aku
}

Quả nhiên, ở đâu có Yuki là ở đó Koishin lại bị chọc ghẹo. Nhưng dù cô nàng tóc tém có cố tỏ ra cằn nhằn hay khó chịu đến đâu, dòng chữ nhỏ xíu ``lớp 1-C là tuyệt nhất'' cũng đã phản bội lại vẻ ngoài gai góc ấy. Còn Aku, với tinh thần của một game thủ thực thụ, cậu ấy có vẻ đã chọn nhầm đối thủ khi định đối đầu với Kaigou ở mảng bắn súng rồi.

Nhưng rồi, ánh mắt tôi chạm đến những dải giấy của nhóm người thứ hai - những người đã biết rằng đêm nay là giới hạn cuối cùng. Lời ước của họ không cầu xin vật chất hay danh vọng, mà chất chứa một sự níu kéo dịu dàng và lòng biết ơn vô hạn.

\txtbx{
  Tuy không thể cùng nhau đi đến cuối con đường, nhưng mình sẽ luôn ghi nhớ những bài học mà các cậu đã mang lại. Chúc các cậu Thượng lộ bình an.
  \flushright - Wata
}

\txtbx{
  Dù ở thế giới nào, Hội Học sinh trường Aogami vẫn luôn dành sẵn ba chiếc ghế cho các cậu. Đừng quên chúng tôi nhé!
  \flushright - Ayame
}

\txtbx{
  Mong rằng dưới sự bảo vệ của bọn mình, trật tự và kỷ luật của thị trấn này sẽ mãi vững bền để không phụ công sức của ba người.
  \flushright - Nanase \& Inari
}

\txtbx{
  Mình đã được tái sinh nhờ các cậu. Lần này, mình sẽ sống một cuộc đời thật rực rỡ, để nếu có duyên gặp lại, các cậu sẽ thấy tự hào về Misora Shino.
  \flushright - Shino
}

\txtbx{
  Thần linh không thể thay đổi dòng chảy không gian, nhưng Thần linh sẽ luôn ban phước cho những người anh hùng. Cảm ơn vì đã cứu rỗi thị trấn này.
  \flushright - Sakura
}

\txtbx{
  Với tư cách là giáo viên chủ nhiệm, tôi cầu chúc cho tương lai của các em ở thế giới bên kia sẽ luôn rạng rỡ. Lớp 1-C sẽ luôn nhớ về các em.
  \flushright - Yuko
}

Sống mũi tôi hơi cay. Tôi hít một hơi thật sâu để ngăn cho tầm nhìn không bị nhòe đi. Kể từ khi đến thế giới này, gia đình Hikawa chưa bao giờ mong đợi sự đền đáp. Nhưng những tấm chân tình này, chúng tôi sẽ mang theo như một hành trang quý giá nhất.

Và cuối cùng, là những lời ước của gia đình Hikawa và Naomi. Khác với sự nuối tiếc, lời ước của chúng tôi mang đậm sự quyết tâm và chấp nhận.

\txtbx{
  Thịt nướng! Thịt nướng! Và thêm nhiều dango nữa! Mong là thế giới bên kia đồ ăn cũng ngon như ở Aogami!
  \flushright - Suzuki
}

\txtbx{
  Trận chiến ở đây đã kết thúc. Hy vọng rằng nơi chúng ta trở về sẽ không cần phải phòng thủ quá trớn thêm lần nào nữa. Bình yên là đủ.
  \flushright - Kaigou
}

\txtbx{
  Naomi sẽ lớn lên thật khỏe mạnh! Cho dù anh chị có đi đâu thì Naomi vẫn sẽ luôn bên cạnh mọi người! Đừng quên em nhé!
  \flushright - Naomi
}

Đọc xong tờ giấy của Naomi, tôi bật cười thành tiếng. Một đứa trẻ hồn nhiên được ban một cuộc sống mới, lúc nào cũng tò mò với mọi thứ xung quanh ở thế giới hiện đại, lúc nào cũng bám theo chúng tôi chỉ để hỏi những câu ngây ngô. Nhưng đó mới chính là Naomi mà chúng tôi yêu quý.

Tôi cầm dải giấy tanzaku màu thiên thanh của mình lên, nét mực đen vẫn còn ướt. Tôi nhón chân, tự tay buộc nó lên cành trúc cao nhất mà mình có thể với tới, để nó đón lấy ngọn gió biển lồng lộng đang thổi vào từ phía Đông.

\txtbx{
  Aogami đã dạy cho tôi biết thế nào là sống chứ không phải chỉ tồn tại. Mong rằng những nụ cười ngày hôm nay sẽ trở thành vĩnh cửu. Hẹn một ngày tái ngộ.
  \flushright - Kuon
}

Tôi lùi lại vài bước, ngắm nhìn hàng trăm dải giấy đang nhảy múa dưới nền trời đã chuyển hẳn sang màu tím thẫm của hoàng hôn. Tiếng trống taiko từ trung tâm lễ hội bắt đầu gõ những nhịp dồn dập hơn, báo hiệu phần cao trào nhất của đêm Natsu Matsuri chuẩn bị bắt đầu.

Mọi người quay sang nhìn tôi, ánh mắt ai nấy đều lấp lánh phản chiếu ánh đèn lồng rực rỡ.

``Đi thôi,'' tôi khẽ nói, giấu hai tay vào ống tay áo yukata. ``Phần lễ hội chính sắp bắt đầu rồi.''

\ct

Khi mặt trời hoàn toàn khuất bóng sau dãy núi phía Tây, nhường chỗ cho mảng màu chạng vạng pha lẫn sắc xanh đen và tím thẫm, cũng là lúc thị trấn Aogami thực sự khoác lên mình tấm áo hội hè rực rỡ nhất. Tầm sáu giờ chiều, hàng loạt những dãy đèn lồng chōchin dọc các con phố cổ đồng loạt bừng sáng, hắt thứ ánh sáng đỏ cam ấm áp lên những mái ngói rêu phong.

Sau khi hoàn thành phần viết lời ước, tôi cùng nhóm bạn lớp 1-C tiếp tục hòa mình vào dòng người trẩy hội. Chúng tôi đi dạo qua hàng chục sạp hàng yatai đang tỏa khói nghi ngút. Kaigou kéo tôi vào một gian hàng bán mặt nạ cáo, nhất quyết ép tôi phải đội thử một chiếc rồi bật cười sảng khoái. Suzuki và nhóm "thực thần" thì dường như vẫn chưa thỏa mãn cơn đói, họ tiếp tục càn quét qua quầy mực nướng và bánh xèo okonomiyaki. Cứ thế, chúng tôi lang thang hết con dốc này đến ngã rẽ khác, thu nhặt từng chút niềm vui vụn vặt của đêm hội Natsu Matsuri.

Mãi cho đến khi một âm thanh trầm hùng vang vọng khắp không gian, thu hút sự chú ý của tất cả mọi người.

Tiếng trống Taiko bắt đầu vang lên dồn dập từ phía đài yagura cao vút giữa sân đền. Những nhịp đập ``thình thịch'' vang dội như đang đánh thức từng tế bào trong cơ thể của hàng ngàn người tham gia lễ hội. Ánh sáng từ hàng trăm chiếc lồng đèn lồng giấy bọc quanh đài hắt xuống một màu đỏ cam rực rỡ, xua tan đi cái lạnh lẽo của màn đêm đang dần buông xuống.

``Nào mọi người! Vòng tròn đang hình thành kìa! Nhanh chân lên!'' Ayame giơ cao tay, vẫy gọi cả lớp 1-C tiến về phía trung tâm.

Điệu nhảy Bon Odori - linh hồn của đêm lễ hội mùa hè - đã chính thức bắt đầu. Từng nhóm người, từ người già đến trẻ nhỏ, đều bắt đầu xếp thành những vòng tròn đồng tâm rộng lớn bao quanh đài Yagura. Tiếng sáo trúc réo rắt bắt đầu hòa nhịp cùng tiếng trống, tạo nên một bản nhạc dân gian vừa rộn rã vừa mang chút u hoài đặc trưng của mùa hè.

Tôi lùi lại vài bước, đứng tựa lưng vào một thân cây cổ thụ gần đó, định bụng sẽ chỉ làm một khán giả thầm lặng quan sát mọi người. Quả thực, chỉ cần đứng nhìn mọi người nhảy múa cũng đủ để trông như một chương trình hài kịch xuất sắc rồi.

Ở phía vòng tròn ngoài cùng, Suzuki đang vung tay múa chân với một nguồn năng lượng thừa mứa như thể em chưa từng trải qua một ngày dài mệt mỏi làm việc ở quán café. Từng động tác của con bé đều mạnh bạo và dứt khoát đến mức có cảm giác như đang tập võ thay vì múa truyền thống.

Ngay bên cạnh, Naomi đang cố gắng bắt chước từng động tác của  em ấy. Cô bé nhỏ nhắn với bộ yukata hoa anh đào nhíu mày tập trung cao độ, vung cái tay nhỏ xíu lên rồi lại xoay vòng. Nhưng nhịp điệu của một bài múa dân gian dường như quá phức tạp so với một đứa trẻ tám tuổi, khiến cô bé hết giẫm phải chân mình lại suýt va vào người bên cạnh.

``Không sao đâu Naomi-chan! Sai nhịp cũng là một loại nghệ thuật đấy! Cứ nhảy nhiệt tình lên!'' Suzuki cười phá lên, kéo tay Naomi xoay một vòng khiến cô bé lảo đảo nhưng rồi cũng bật cười khúc khích.

Cách đó vài bước, Miku đang vừa vỗ tay nhịp nhàng vừa hướng dẫn động tác múa. Khác với vẻ lúng túng thường ngày khi đứng trước ống kính hay những người mẫu chuyên nghiệp, cô nàng lúc này trông vô cùng tự nhiên và rạng rỡ.

``Miku-chan múa dẻo thật đó\~{} Nhìn cậu lúc này trông rất đáng yêu lắm\~{}'' Saya đi bên cạnh khẽ nghiêng đầu mỉm cười, đôi mắt hiền hòa nheo lại.

Nghe lời khen bất ngờ từ cô bạn thân, vành tai cô nàng bờm Sói đen lập tức đỏ lựng lên. Cô nàng luống cuống dừng lại, hai tay ôm lấy má bối rối: ``C-Cậu nói gì vậy Saya-san! T-Tôi múa đại theo điệu nhạc thôi\dots\ Đừng có chọc tôi chứ!''

Saya chỉ che miệng cười khúc khích, kéo tay cô nàng bờm Tóc tết tiếp tục hòa vào vòng tròn.

Đưa mắt sang một góc khác, tôi không khỏi thở dài ngao ngán. Shion, thay vì múa theo điệu Bon Odori, lại đang thực hiện những động tác kỳ quái: hai tay chắp lại trước ngực rồi vung mạnh lên trời, miệng lầm bầm những câu thần chú nghe như đang triệu hồi quỷ dữ từ địa ngục.

``Này thuộc hạ, ngươi phải làm thế này này! Đây là điệu múa giao tiếp với các linh hồn hắc ám!'' Shion quay sang hướng dẫn một cách đầy tự hào.

Aku đi ngay phía sau, khuôn mặt đỏ bừng vì ngượng. Cô nàng vụng về thử nâng tay lên theo hướng dẫn của Shion, nhưng động tác cứng đờ như một cỗ máy hết pin. ``S-Shion-san\dots\ mọi người đang nhìn chúng ta kìa\dots\ T-Tôi thấy xấu hổ quá\dots''

``Đừng để tâm đến ánh mắt của phàm nhân! Khí chất của bóng tối không cho phép chúng ta lùi bước! Hãy giải phóng phong ấn trên cánh tay ngươi đi!''

Tôi lắc đầu ngao ngán, định bước tới giải vây cho cô bạn nhút nhát thì một tổ hợp trái ngược hoàn toàn khác lại thu hút sự chú ý của tôi.

Akane, với mái tóc tết gọn gàng năng động, đang nhảy cực kỳ sung sức, nụ cười rạng rỡ tỏa sáng cả một vùng. Trái ngược lại, cô bạn thân Yuka của cậu ấy thì đi theo sau với đôi mắt lim dim buồn ngủ, tay vung vẩy hờ hững như một bóng ma vất vưởng thiếu sức sống.

``Yuka! Mạnh tay lên nào! Lễ hội mà, phải bùng nổ lên chứ!'' cô nàng bờm Chó hò reo, vỗ một cái rõ kêu vào lưng Yuka.

``Đau\dots\ Tôi đang múa điệu 'mèo lười giỡn nắng' mà\dots\ Cho tôi ngủ đi\dots'' cô bạn bờm Mèo ngáp dài một cái, giọng nhừa nhựa, bước chân lảo đảo suýt va vào hàng người bên cạnh.

May mắn thay, Mitsuki đi ngay phía sau đã kịp thời đưa tay đỡ lấy Yuka. Với phong thái thanh lịch và dịu dàng, cô nàg không chỉ múa chuẩn xác từng nhịp điệu mà còn khéo léo điều hướng để cô bạn tóc tím không bị vấp ngã.

``Cậu mệt quá thì vào trong nghỉ ngơi một lát đi, Yuka-san,'' Mitsuki nhẹ nhàng khuyên bảo. Trông Mitsuki lúc này lộng lẫy và trưởng thành hệt như một nàng tiểu thư thực thụ bước ra từ những câu chuyện cổ tích.

Tôi khẽ mỉm cười, cảm nhận một sự bình yên kỳ lạ dâng lên trong lồng ngực. Một tập thể với đủ mọi cá tính kỳ quặc, luôn ồn ào và rắc rối, nhưng lại gắn kết với nhau bằng một thứ tình cảm chân thành đến lạ lùng. Tôi thầm nghĩ, liệu ở một nơi nào đó trên thế giới này, hay ở một chiều không gian khác, có một nơi nào mang lại cho tôi cảm giác ấm áp như thế này không?

Bỗng nhiên, một bàn tay mềm mại nắm lấy cổ tay tôi, kéo mạnh khiến tôi chới với bước lên phía trước.

Mùi hương hoa bách hợp nhè nhẹ thoảng qua mũi. Là Nanase.

Cô gái với bộ yukata họa tiết thanh nhã đang mỉm cười nhìn tôi, đôi mắt xanh thẫm lấp lánh phản chiếu ánh đèn lồng rực rỡ của đài yagura. Mái tóc dài được búi cao để lộ phần gáy trắng ngần, những giọt mồ hôi li ti lấm tấm trên trán chứng tỏ cô ấy đã nhảy múa từ nãy đến giờ.

``Cậu định đứng đó làm một ông già khó tính đến bao giờ hả, Kuon-san?'' Nanase nghiêng đầu, giọng điệu vừa trách móc vừa mang chút nũng nịu hiếm thấy.

``Tôi không có khiếu nhảy múa đâu,'' tôi lảng tránh ánh mắt của cô ấy, khẽ gãi má. ``Vào đó chỉ tổ làm rối đội hình của mọi người thôi. Hơn nữa, để một kẻ cứng nhắc như tôi múa hát trông dị lắm.''

``Không có lý do từ chối nào được chấp nhận đêm nay đâu,'' cô ấy siết chặt lấy tay tôi, nụ cười càng thêm rạng rỡ. ``Đã đến lễ hội thì phải nhảy. Cậu là thành viên của lớp 1-C cơ mà. Hơn nữa\dots\ đây là cơ hội để cậu thực sự cảm nhận được nhịp đập của Aogami đấy.''

Chữ ``cơ hội'' được Nanase nhấn mạnh, như một lời nhắc nhở ngầm về khoảng thời gian ít ỏi còn lại của chúng tôi. Lòng tôi bỗng chùng xuống một nhịp, nhưng rồi ngay lập tức được xoa dịu bởi hơi ấm từ bàn tay cô ấy.

Cô nàng kéo tôi vào vòng tròn, chen vào khoảng trống giữa cô ấy và Inari. Nhịp trống Taiko dội vào lồng ngực tôi từng hồi mạnh mẽ.

Tôi bắt đầu vung tay, bước chân những bước lóng ngóng đầu tiên. Quả thật, tôi hoàn toàn trật nhịp so với mọi người. Tay trái giơ lên thì tay phải lại quên hạ xuống, bước lên thì lại dẫm nhầm nhịp của Inari phía trước. Nhưng điều kỳ lạ là chẳng có ai cười nhạo hay chê bai.

Ayame từ phía trước quay lại nháy mắt tinh nghịch: ``Thả lỏng ra đi Kuon-san! Cậu đang đi hành quân đấy à?''

Wata khẽ gật đầu mỉm cười khích lệ.

Và Nanase bên cạnh thì kiên nhẫn hướng dẫn tôi từng động tác một.

``Chân phải bước trước, sau đó lùi lại một bước, hai tay vung sang ngang\dots\ Đúng rồi, nhịp nhàng theo tiếng trống thôi, cậu làm tốt lắm Kuon-kun!''

Tiếng cười nói rôm rả, tiếng nhạc xập xình, cái nóng hầm hập phả ra từ biển người xung quanh, tất cả hòa quyện thành một luồng sinh khí khổng lồ cuộn trào trong huyết quản. Cảm giác xa lạ và khoảng cách vô hình mà tôi luôn tự dựng lên với thế giới này bỗng chốc vỡ vụn.

Trong khoảnh khắc cùng mọi người nhấc chân theo điệu nhạc, tôi chợt nhận ra mình đang thực sự tồn tại ở hiện tại này. Không còn những toan tính về vòng lặp, không còn nỗi lo về các thế lực bám đuổi, chỉ còn lại một Hikawa Kuon đang tận hưởng niềm vui giản dị của một nam sinh trung học bình thường.

Gió biển lồng lộng thổi tới mang theo chút mằn mặn đặc trưng, làm những dải ruy băng ngũ sắc treo trên đài Yagura tung bay phấp phới. Tôi ngước nhìn lên bầu trời đêm không gợn bóng mây, nơi vầng trăng rằm đang tỏa ra thứ ánh sáng dịu dàng như dát bạc lên vạn vật. Tiếng cười của Nanase bên cạnh nghe trong trẻo như tiếng chuông gió.

``\textit{Nếu thời gian có thể ngừng lại, tôi thực sự muốn khoảnh khắc này kéo dài mãi mãi\dots}'' tôi thầm nghĩ, để mặc cho những âm điệu rộn rã của lễ hội cuốn đi mọi muộn phiền.

\ct

Đồng hồ điểm bảy giờ ba mươi tối.

Điệu múa Bon Odori dần thưa nhịp khi đám đông bắt đầu dạt sang hai bên, tạo thành một khoảng trống lớn dẫn từ cổng đền vào tận sân chính. Bầu không khí bỗng nhiên chuyển sang một trạng thái căng tràn năng lượng và đầy tính chiến đấu.

Từ phía xa, một chiếc kiệu mikoshi được sơn son thếp vàng lộng lẫy, trang hoàng bằng những dải lụa trắng và dây thừng bện thô, đang từ từ tiến vào. Bao quanh chiếc kiệu là hàng chục thanh niên trai tráng của thị trấn, mồ hôi nhễ nhại, vai kề vai gánh vác sức nặng của vị thần bảo hộ.

Nhưng điều khiến tôi bất ngờ nhất là ở ngay phía trước chiếc kiệu ấy, một nhóm các cô gái lớp 1-C cũng đang hừng hực khí thế chuẩn bị tham gia vào đội ngũ rước kiệu.

``Đến lượt chúng ta rồi! \textbf{Mọi người chuẩn bị tinh thần chưa nào!?}'' Hanabi đứng trên một bục đá nhỏ, giơ cao chiếc quạt giấy, vung mạnh xuống như một vị tướng quân ra lệnh xuất chinh.

``\textbf{Sẵn sàng!}'' Kaoru, Kana và Honoka đồng thanh hét lớn, xắn cao ống tay áo yukata. Thậm chí cả cô nàng thường ngày trầm tính như Kanade cũng đang mỉm cười rạng rỡ, sẵn sàng hòa vào dòng người náo nhiệt.

``\textbf{Kuon-kun! Lại đây nhanh lên! Cậu định trốn việc rước kiệu đấy à?}'' Hanabi chỉ thẳng chiếc quạt về phía tôi, ánh mắt rực lửa không cho phép từ chối.

Kaigou từ lúc nào đã đứng sẵn trong hàng ngũ, vai tựa vào thanh xà bằng gỗ lim của chiếc kiệu, hếch cằm nhìn tôi với vẻ thách thức: ``Hay là cậu yếu đuối quá không vác nổi kiệu thế?''

``\textit{Thì đúng thế mà\dots}'' tôi thầm nghĩ. Nhưng rồi, chưa kịp nói lên suy nghĩ của mình, tôi đã bị mọi người dồn tới, xắn tay áo bước vào đội hình, đứng ngay phía trước Kaigou. Vừa ghé vai vào thanh xà gỗ, sức nặng khủng khiếp của chiếc kiệu lập tức đè nghiến xuống khiến tôi hơi lảo đảo.
`'
``Nào, tất cả cùng hô vang nhé! Một, hai, ba\dots'' Hanabi bắt nhịp.

``\textbf{Wasshoi! Wasshoi!}''

Tiếng hô đồng thanh vang dội rúng động cả một góc trời.

``\textbf{Wasshoi! Wasshoi!}''

Nhịp điệu dồn dập, mạnh mẽ và đầy tính nguyên thủy. Mỗi một tiếng ``Wasshoi'' vang lên, chiếc kiệu lại được nâng lên rồi hạ xuống theo nhịp bước chân của hàng chục người.

Nghe tiếng hô vang dội của cô nàng năng nổ ngay bên cạnh, tâm trí tôi bỗng chốc mờ đi, một đoạn ký ức từ ở thế giới gốc bất chợt ùa về.

Đó là vào khoảng thời gian lúc tôi mới nhận chức tổng quản không lâu. Khi đó, tôi chỉ mới bắt đầu chuỗi ngày làm quần quật tháng ngày qua năm khác để kiếm tiền trang trải cuộc sống để đỡ đần bố mẹ. Trong một ngày mà tôi làm tới mức kiệt sức, Mat-chan - tức Matsuri - đã từng nói với tôi thế này: ``Kuon-senpai! Khi nào anh thấy mệt mỏi, hãy cứ hét thật to lên! Năng lượng của lễ hội sẽ cuốn trôi mọi buồn phiền!''

Giọng nói từ quá khứ đan xen với hiện tại. Tôi quay sang nhìn Hanabi. Tên có thể khác, nhưng vẫn là cậu ấy, vẫn là ngọn lửa sưởi ấm và truyền năng lượng cho tất cả mọi người xung quanh. Dù là ở thế giới nào, bản chất của những con người này vẫn luôn rực rỡ và chân thành như vậy.

``\textbf{Wasshoi! Wasshoi!}''

Lần này, tôi chủ động mở miệng hô lớn, âm thanh thoát ra từ tận sâu trong lồng ngực. Mồ hôi bắt đầu tuôn ra như tắm, ướt đẫm cả mảng lưng áo yukata, nhưng sự mệt nhọc dường như tan biến hết. Sức mạnh của cộng đồng, sự gắn kết truyền từ vai người này sang vai người khác qua thanh xà gỗ, tạo nên một cảm giác thiêng liêng khó tả.

``\textbf{Mọi người tuyệt lắm! Cứ giữ nét mặt đó nhé!}''

*TÁCH!* *TÁCH!*

Tiếng màn trập máy ảnh vang lên liên tục. Uzuki đang chạy lăng xăng quanh đoàn rước kiệu, liên tục thay đổi góc máy để bắt trọn những khoảnh khắc mồ hôi nhễ nhại nhưng nụ cười rạng rỡ của chúng tôi. Chiếc máy ảnh trên tay cô bạn bờm Thỏ như một phép màu, cố gắng đóng băng lại khoảng thời gian quý giá này.

``\textbf{Uzuki-chan! Chụp cho tớ một tấm thật ngầu vào nhé!}'' Honoka gào lên qua tiếng ồn.

``Yên tâm! Tớ sẽ bắt được khoảnh khắc cậu đang cắn môi vì mỏi vai!'' Uzuki cười tinh nghịch đáp trả.

Chiếc kiệu tiếp tục được diễu hành quanh sân đền thêm ba vòng nữa trước khi được hạ xuống vị trí trung tâm trong tràng pháo tay rầm rộ của mọi người. Bờ vai tôi tê rần, hơi thở gấp gáp, nhưng khi nhìn sang Kaigou, Hanabi, và những người dân của thị trấn, ai nấy đều đang mỉm cười đầy thỏa mãn.

\ct

Sau khi chiếc kiệu mikoshi yên vị ở chính giữa sân đền, sự cuồng nhiệt của đám đông cũng dần lắng xuống, nhường chỗ cho một bầu không khí trầm mặc và sâu lắng hơn. Những ngọn đèn lồng đỏ xung quanh sân đồng loạt được hạ độ sáng, chỉ để lại ánh trăng rằm vằng vặc chiếu rọi xuống đài yagura ở vị trí trung tâm.

Một vài người bắt đầu tản ra các góc để thưởng thức những món ăn nhẹ cuối cùng của lễ hội, nhưng phần lớn học sinh lớp 1-C và người dân thị trấn đều hướng ánh mắt lên sân khấu gỗ.

``Cảm ơn mọi người đã cháy hết mình cùng chúng mình từ hôm qua đến giờ,'' giọng trầm ấm của cô nàng mạo hiểm Azuki vang lên qua hệ thống loa, kéo sự chú ý của toàn bộ sân đền. Cô nàng ca sĩ chính của câu lạc bộ Âm nhạc đang ngồi trên một chiếc ghế đẩu, ôm cây đàn guitar acoustic trước ngực. Bên trái cô là Honoka với nụ cười rạng rỡ thường ngày nhưng có phần e ấp hơn khi cầm micro, và bên phải là Yuki đang nhẹ nhàng lướt những ngón tay trên phím đàn keyboard điện tử.

``Bài hát sau đây\dots\ cũng là bài hát cuối cùng của đêm hội hôm nay,'' Yuki cất lời, ánh mắt lướt qua toàn bộ những gương mặt quen thuộc đang đứng bên dưới. ``Đây là một ca khúc do chính tớ sáng tác, dành tặng riêng cho lớp 1-C của chúng ta, và cho cả thị trấn Aogami này. Ca khúc được mang tên, ``Mùa hè năm ấy, chúng ta từng thuộc về nhau.'' % Gemini made this up, not me.

Một tràng vỗ tay nhẹ nhàng vang lên rồi nhanh chóng chìm vào tĩnh lặng.

Azuki khẽ gật đầu, ngón tay gảy những nốt nhạc đầu tiên. Tiếng đàn guitar acoustic mộc mạc vang lên, hòa quyện với âm thanh trong trẻo của keyboard tạo nên một giai điệu vừa êm đềm vừa mang mác buồn. Rồi Honoka bắt đầu cất giọng hát. Không còn là cô nàng thích buôn chuyện ồn ào thường ngày, giọng hát của cô nàng lúc này trong vắt và đầy cảm xúc, như một cơn gió mát lành xoa dịu đi cái nóng bức của đêm hè.

Tôi đứng tựa vai vào một cây cột đèn bằng gỗ, lặng lẽ quan sát những người bạn xung quanh mình.

Ở ngay phía trước, Suzune - cô bé được mệnh danh là ``linh vật của lớp'', người luôn lao đầu vào mọi thứ thú vị và thừa năng lượng - lúc này đang đứng yên lặng đến lạ thường. Đôi mắt to tròn của cô phản chiếu ánh trăng lấp lánh, hai tay chắp trước ngực như đang cầu nguyện. Không có sự tăng động, không có những trò đùa tinh nghịch, Suzune đang đắm chìm trọn vẹn vào từng lời ca, thân hình nhỏ nhắn khẽ đung đưa theo nhịp điệu.

Đứng cạnh cô là Saki. Cô con gái nhà bán hoa nở một nụ cười hiền hậu, ánh mắt mông lung hướng về phía đài yagura. Saki luôn là một người chu đáo và tốt bụng. Nhìn cách cô ấy khẽ vuốt ve dải ruy băng buộc trên cổ tay, tôi đoán cô ấy đang liên tưởng vòng đời ngắn ngủi của những đóa hoa mùa hè với những kỷ niệm đẹp đẽ mà chúng tôi đang trải qua. Rực rỡ nhất lúc nở rộ, để rồi để lại sự nuối tiếc khi tàn phai.

Dịch sang bên trái một chút, Shiina và Yae đang đứng tựa lưng vào nhau. Shiina - cô gái được mệnh danh là mạnh nhất Aogami với tính cách nóng nảy - giờ đây khoanh tay trước ngực, đôi mắt nhắm hờ. Sự gai góc thường ngày hoàn toàn biến mất, nhường chỗ cho một nét điềm tĩnh và tĩnh lặng đến đáng ngạc nhiên. Còn Yae, vốn dĩ là một cô gái năng động của câu lạc bộ bóng rổ nhưng lại cực kỳ nhạy cảm với cảm xúc của người khác, đang khẽ sụt sịt mũi. Cô bạn cố tình lấy tay áo yukata quệt vội qua khóe mắt, cố giấu đi những giọt nước mắt đang chực trào vì giai điệu quá đỗi da diết của bài hát.

Và tất nhiên, không thể không nhắc đến Koishin. Cô bạn tsundere của câu lạc bộ Văn học đứng cách tôi một đoạn ngắn, vẻ mặt cố tỏ ra không quan tâm: ``Hừm, giai điệu cũng bình thường thôi, làm gì mà mọi người xúc động thế\dots'' Mặc dù cái miệng sắc bén lẩm bẩm như vậy, nhưng ánh mắt của Koishin lại dán chặt lên sân khấu, đặc biệt là hướng về phía Yuki.

Như cảm nhận được ánh nhìn đó, Yuki đang chơi đàn bỗng ngước lên, hướng thẳng mắt về phía Koishin và nháy mắt một cái đầy tinh nghịch. Gương mặt Koishin lập tức đỏ lựng lên như quả cà chua, cô nàng bối rối quay ngoắt mặt đi chỗ khác, hai tay nắm chặt lấy vạt áo yukata, miệng tiếp tục lầm bầm những câu mắng mỏ không rõ nghĩa nhưng khóe môi lại vô thức cong lên. Trái tim nhân hậu của cô nàng thực sự không thể giấu giếm.

Tôi thu ánh mắt về, khẽ nhắm lại để cảm nhận trọn vẹn ca từ bài hát.

\begin{center}
  ``Chúng ta gặp nhau nơi trạm dừng chân của định mệnh\\
  Như những cánh hoa trôi dạt tìm bến đỗ\\
  Ngôi nhà tạm này sẽ mãi giữ hơi ấm\\
  Dù ngày mai con đường có chia hai\dots''
\end{center}        % Again, Gemini made those.

Từng câu chữ như xoáy sâu vào tâm trí tôi. Lời bài hát do Yuki sáng tác không trực tiếp nhắc đến sự chia ly của nhà Hikawa chúng tôi hay bí mật về thế giới thực, nhưng cái ý niệm về một ``ngôi nhà tạm'' và ``sự trôi dạt'' lại miêu tả quá chính xác tình cảnh của chúng tôi lúc này. Sự đồng cảm vô tình này khiến lời ca có sức nặng lạ thường.

Lớp 1-C đối với tôi, với Kaigou, với Naomi, hay với bất kỳ ai trong chúng tôi, chính là một ngôi nhà tạm tuyệt vời nhất. Ở đây, tôi được trải qua những ngày tháng học đường bình thường mà tôi đã từng trải qua. Tôi có những người bạn sẵn sàng cãi nhau chí chóe nhưng cũng sẵn sàng bảo vệ nhau. Tôi có những ký ức không bị bôi xóa bởi những vòng lặp vô tận chết chóc.

Nhưng ngôi nhà tạm thì cuối cùng vẫn chỉ là trạm dừng chân. Nghi thức Thanh tẩy đang đến rất gần. Đêm nay, dưới ánh trăng rằm này, là giới hạn cuối cùng để chúng tôi được làm những học sinh bình thường.

Một nỗi buồn man mác len lỏi vào tận sâu trong tim, nhưng kỳ lạ thay, nó không hề mang theo sự tuyệt vọng. Khác với những nỗi đau giằng xé ở các vòng lặp trước, nỗi buồn lúc này giống như một ly trà ấm, hơi chát ở đầu lưỡi nhưng lại đọng lại vị ngọt nơi cuống họng. Bởi vì tôi biết, những gì tôi trải qua ở Aogami không phải là ảo ảnh. Tình cảm của Nanase, sự tin tưởng của Hội học sinh, nụ cười của Suzuki, hay tiếng hát lúc này của nhóm Azuki\dots\ tất cả đều là chân thật.

\begin{center}
  ``Cảm ơn vì đã là một phần tuổi trẻ của tôi\\
  Dưới bầu trời đêm nay, xin hãy nhớ về nhau\dots''
\end{center}

Đoạn điệp khúc cuối cùng vang lên, giọng hát của Honoka kết hợp cùng tiếng bè của Azuki và Yuki vút cao lên tận tầng không, xé toạc màn đêm tĩnh mịch trước khi hòa tan vào tiếng sóng biển vỗ rì rào từ xa vọng lại. Nốt nhạc cuối cùng từ phím đàn keyboard của Yuki vang lên thánh thót rồi từ từ nhỏ dần, để lại một không gian tĩnh lặng đến nghẹt thở.

Mọi người dường như vẫn còn chìm đắm trong dư âm của bài hát, không ai dám lên tiếng hay cử động, sợ rằng sẽ phá vỡ mất khoảnh khắc thiêng liêng này.

Và rồi, một giọt nước mắt rơi xuống. Yae không thể kìm nén được nữa, cô ấy ôm chầm lấy Shiina bên cạnh mà nức nở. Như một mồi lửa được châm ngòi, tiếng vỗ tay rào rào bắt đầu nổ ra từ khắp các góc sân đền. Tiếng huýt sáo, tiếng hò reo tán thưởng vang lên không ngớt.

``Tuyệt vời lắm!''

``Bài hát hay nhất tớ từng nghe!''

``Yuki, Azuki, Honoka! Các cậu là số một!''

Suzune nhảy cẫng lên vẫy tay cuồng nhiệt, Saki mỉm cười lau khóe mắt, và ngay cả Shiina cũng khẽ gật đầu mỉm cười đầy tán thưởng trong khi vẫn vỗ về lưng Yae. Koishin tuy quay mặt đi nhưng đôi vai đang run lên nhè nhẹ, có lẽ cô nàng cũng đang vất vả giấu đi sự xúc động của mình.

Tôi mở mắt ra, nhìn lên bầu trời rực rỡ ánh trăng. Khúc nhạc dạo đầu của cuộc chia ly đã kết thúc một cách hoàn mỹ nhất có thể. Giao ước của một lễ hội trọn vẹn đã được thực hiện.

Bây giờ, đã đến lúc đối mặt với thực tại.

Tôi khẽ đưa tay chạm vào chiếc đồng hồ đeo tay - nơi Game đang trú ngụ.

``\eng{Chuẩn bị tinh thần cho nghi thức đi, Kuon-user.}''

``\eng{Ừm. Tôi biết rồi.}''

\ct

Tiếng vỗ tay và những lời tán dương dành cho câu lạc bộ Âm nhạc dần lắng xuống, nhưng sự náo nhiệt của đêm hội vẫn chưa hề có dấu hiệu hạ nhiệt. Mọi người bắt đầu di chuyển khỏi khu vực đài yagura, tụ tập thành từng nhóm nhỏ dọc theo những bậc thang đá hoặc những khoảng sân trống hướng tầm nhìn lên cao.

Tôi vừa lấy lại nhịp thở bình thường thì âm thanh thông báo quen thuộc vang lên.

``Mọi người ơi! Chú ý nào!'' giọng nói nghịch ngợm của Hội trưởng Ayame vang vọng qua chiếc loa cầm tay, cắt ngang những tiếng xì xào bàn tán. ``Bây giờ đã là tám giờ ba mươi phút tối! Thời khắc mà tất cả chúng ta mong chờ nhất đêm lễ hội mùa hè đã đến! Mọi người đã chọn được vị trí đẹp chưa?''

Tiếng hò reo đồng thanh vang lên đáp lời thay cho sự đồng ý. Những ngọn đèn lồng chōchin xung quanh đền cũng được ban tổ chức tắt bớt, nhường lại một không gian tối sẫm, biến bầu trời Aogami thành một tấm màn nhung khổng lồ hoàn hảo.

``Tốt lắm! Vậy thì\dots\ tất cả cùng đếm ngược nào!'' Ayame giơ cao chiếc loa, tay kia chỉ thẳng lên bầu trời sao.

Cả sân đền, từ học sinh lớp 1-C cho đến những người dân thị trấn, đều đồng loạt hòa chung một nhịp đập.

\textbf{``Mười!''}
\textbf{``Chín!''}
\textbf{``Tám!''}

Sự phấn khích lan tỏa trong không khí đặc quánh, khiến nhịp tim tôi cũng bất giác đập nhanh hơn. Tôi nhìn sang Nanase bên cạnh, cô ấy đang nắm chặt hai tay, đôi mắt xanh thẫm hướng lên trời đầy kỳ vọng.

\textbf{``Ba!''}
\textbf{``Hai!''}
\textbf{``Một!''}
\textbf{``Không!!!''}

*\textbf{*ĐOÀNG!}

Một tiếng nổ đinh tai nhức óc vang lên từ ngọn núi, kéo theo một luồng ánh sáng chói lòa xé toạc màn đêm tĩnh mịch. Một quả cầu lửa khổng lồ vút lên cao, để lại một vệt khói xám nhạt, trước khi nổ tung thành hàng ngàn những tia sáng màu đỏ thẫm và vàng kim lấp lánh rực rỡ.

Bầu trời Aogami đêm nay không còn là bóng tối vô tận. Nó bừng sáng bởi đủ mọi gam màu: từ sắc đỏ rực lửa, xanh dương huyền bí, cho đến màu ngọc lục bảo lung linh và ánh bạc lấp lánh. Những tràng pháo hoa liên tiếp được bắn lên, nở bung như những bông hoa cúc khổng lồ vươn những cánh hoa ánh sáng rủ xuống. Tiếng nổ trầm hùng vang vọng vào các vách núi, dội lại tạo thành một bản giao hưởng tráng lệ của mùa hè.

``\vie{Naomi-chan, đẹp không em?}'' Suzuki khẽ thốt lên, vòng tay ôm trọn lấy bả vai của Naomi từ phía sau. Đôi mắt của em lúc này không còn chứa đựng hình ảnh của những đĩa thức ăn nữa, mà hoàn toàn bị thu hút bởi ánh sáng rực rỡ trên cao.

Cô bé mắt xanh lục bảo ngước nhìn bầu trời, đôi mắt mở to hết cỡ phản chiếu lại những tia sáng lấp lánh như sao sa. Em vòng tay ôm lấy cánh tay của Suzuki, chất giọng trẻ thơ cất lên, trong trẻo nhưng lại mang một sức nặng khiến trái tim tôi như thắt lại.

``\vie{Vâng\dots\ Đẹp lắm ạ. Em muốn được ở đây mãi mãi.}''

Khát khao thuần khiết của một sinh mệnh mới được tái sinh vang lên giữa tiếng pháo hoa nổ đì đùng. Tôi khẽ cắn môi, cố nén lại tiếng thở dài.

Ngay bên cạnh, Kaigou cũng đang lặng lẽ rút điện thoại ra. Cô nàng không nói một lời nào, chỉ khẽ nhếch mép tạo thành một nụ cười dịu dàng hiếm thấy, liên tục bấm máy để ghi lại khung cảnh rực rỡ này. Sự gai góc thường ngày của cô ấy dường như đã bị pháo hoa nung chảy mất rồi.

Đưa mắt sang nhóm học sinh lớp 1-C, mỗi người lại có một cách thưởng thức màn trình diễn này theo cá tính riêng của mình.

``Màu đỏ thẫm pha lẫn với chút vàng kim\dots\ Tỷ lệ lan tỏa ánh sáng thật hoàn hảo!'' Mari đưa hai bàn tay lên, tạo thành một khung hình chữ nhật hướng về phía những chùm pháo hoa nổ chậm. ``Ánh sáng rực rỡ đến mức mình muốn dùng cọ vẽ lại ngay lập tức. Cảm hứng nghệ thuật đang dâng trào!''

``\textbf{Uwaaaa! Tuyệt quá! Tuyệt đỉnh luôn!}'' Yae nhảy cẫng lên, hai tay vung vẩy loạn xạ trong không trung. ``Tớ cảm giác như lồng ngực mình cũng sắp nổ tung theo tiếng pháo luôn ấy!''

Azuki thì nhắm hờ mắt lại, hai tay chắp trước ngực như đang cảm nhận từng rung động của âm thanh: ``Độ vang của tiếng nổ dội lại từ vách núi thật tuyệt! Phải chi mình có thể ngân lên một nốt cao có độ rền như thế này trong buổi biểu diễn ban nãy!''

Trong khi mọi người đang chìm đắm trong sự phấn khích, Nanase khẽ nghiêng đầu sang phía tôi, khóe môi điểm một nụ cười đầy tự hào. Cô nàng hất nhẹ mái tóc dài, thì thầm: ``Thế nào Kuon-san? Màn pháo hoa do tập đoàn Furukawa tài trợ năm nay không làm cậu thất vọng chứ? Tôi đã tự tay kiểm duyệt thiết kế của từng quả pháo đấy.''

``Rất ấn tượng,'' tôi thành thật đáp lại, ánh mắt không giấu nổi sự ngưỡng mộ. ``Đẹp đến mức khiến người ta cảm thấy choáng ngợp.''

Ở phía sau chúng tôi, Touka đang giương ánh mắt lên bầu trời, ánh sáng pháo hoa phản chiếu lấp lánh trên mắt của cô nàng thần đồng. ``Góc bắn nghiêng 75 độ, thành phần chủ yếu là Stronti\footnote{Một kim loại kiềm thổ (cùng nhóm với Canxi, Magiê\dots) và có số nguyên tử là 38.} tạo màu đỏ và Bari\footnote{Một kim loại kiềm thổ và có số nguyên tử là 56.} tạo màu xanh lá\dots\ Dù có tính toán quỹ đạo khí động học thế nào thì vẻ đẹp ngẫu nhiên khi các hóa chất phát nổ trong không khí vẫn là một biến số cực kỳ khó lường.''

``Và đó chính là sự kỳ diệu của lễ hội đấy, Touka-chan!'' Kaoru đột nhiên nhảy ra từ phía sau, cuộn tròn một tờ ruy băng làm micro giả, chĩa thẳng về phía nhóm chúng tôi. ``Xin chào quý vị khán giả! Trực tiếp từ đền Aogami, đài phát thanh nghiệp dư Kaoru xin tường thuật màn pháo hoa hoành tráng nhất lịch sử thị trấn! Hãy nhìn những nụ cười này đi!''

*TÁCH!* *TÁCH!*

Ngay khi cô nàng tóc ngắn dứt lời, Uzuki đã nhanh tay bấm máy liên tục. ``Tuyệt lắm! Đừng chỉ nhìn lên trời, mọi người hãy cười tươi lên nào! Phải bắt được khoảnh khắc nụ cười của mọi người dưới ánh pháo hoa\dots\ Đây sẽ là những tác phẩm để đời của tớ trong cuốn kỷ yếu!''

Mitsuki đứng cạnh cô bạn bờm Thỏ, đôi mắt mơ màng chiêm ngưỡng sự pha trộn màu sắc trên cao. ``Sự chuyển màu mượt mà từ tím sang bạc\dots\ Thật là một nguồn cảm hứng tuyệt vời. Mình nghĩ mình đã phác thảo xong ý tưởng cho bộ sưu tập thời trang mùa đông của câu lạc bộ rồi na\~{}''

Trong khi những người khác không tiếc lời khen ngợi, Koishin lại khoanh tay trước ngực, cố gắng hất mặt đi nhưng khóe mắt vẫn không ngừng liếc trộm lên bầu trời. ``H-Hừm\dots\ Cũng tàm tạm thôi. So với pháo hoa ở thủ đô thì\dots\ thì\dots'' Cô nàng lắp bắp, đôi má đỏ bừng lên dưới ánh sáng đỏ rực, cuối cùng cũng chịu thua chính mình. ``Được rồi! Tôi thừa nhận là nó rất rực rỡ! Các cậu hài lòng chưa?!''

Sự ồn ào của lớp 1-C càng được đẩy lên cao trào khi Shion bắt đầu tạo dáng đầy vẻ nguy hiểm.

``Hahaha! Bầu trời đang nứt toác ra!'' Shion dang rộng hai tay, vạt áo yukata bay phấp phới. ``Ma thuật bùng nổ của Hỏa Thần đang giáng xuống để thanh tẩy màn đêm vô tận! Hãy cảm nhận đi, những thuộc hạ trung thành của bóng tối!''

Aku đi ngay phía sau, luống cuống lấy hai tay che đầu nhưng vẫn hùa theo màn kịch của Shion một cách đầy nhập tâm: ``V-Vâng, thưa Nữ hoàng! Năng lượng hủy diệt này thật đáng sợ\dots\ T-Thần sẽ giương khiên ma thuật để bảo vệ mọi người khỏi tàn lửa!''

Tôi nhìn một lượt tất cả những người bạn của mình, những gương mặt đang được thắp sáng bởi hàng ngàn tia lửa lấp lánh trên không trung. Tiếng cười đùa rôm rả, những lời trêu chọc hòa cùng tiếng pháo nổ tạo thành một bản giao hưởng tuyệt đẹp.

Nhưng trong lòng tôi lại trĩu nặng một cảm giác bồn chồn khó tả.

Pháo hoa, dù rực rỡ và lộng lẫy đến đâu, thì bản chất của nó vẫn là sự phù du. Nó nở rộ trong một chốc lát, thắp sáng cả bầu trời đêm, để rồi cuối cùng hóa thành tàn tro rơi rụng vào hư vô, trả lại không gian mịt mù tăm tối. Khoảnh khắc đẹp nhất cũng chính là khoảnh khắc chuẩn bị lụi tàn.

Nó giống hệt như khoảng thời gian của chúng tôi - những kẻ ngoại lai - ở tại thị trấn Aogami này. Chúng tôi đã có một mùa hè rực rỡ nhất, được sống trọn vẹn với những cảm xúc chân thật nhất, nhưng thời gian thì không bao giờ dừng lại.

``Kuon-san?'' Nanase khẽ gọi tên tôi, kéo tôi ra khỏi dòng suy nghĩ mông lung. Đôi mắt xanh của cô ấy nhìn tôi với vẻ dò xét, dường như đã nhận ra sự thay đổi trong nét mặt của tôi. ``Cậu sao vậy?''

Tôi hít một hơi thật sâu, để cho không khí mang đậm mùi thuốc súng và vị sương của núi cao lấp đầy buồng phổi.

Tiếng pháo hoa nổ bùm bụp trên đỉnh đầu đã tạo ra một bức tường âm thanh hoàn hảo, che lấp mọi tiếng ồn từ đám đông xung quanh. Đây chính là sân khấu lý tưởng nhất. Khung cảnh rực rỡ nhất, ồn ào nhất, nhưng cũng là lúc để đối mặt với sự thật tĩnh lặng nhất.

Tôi quay sang nhìn Nanase, rồi đánh mắt sang Inari, Ayame, Wata - những người thuộc Hội học sinh đã biết trước bí mật từ những đêm trước. Bọn họ khẽ gật đầu, ánh mắt ánh lên sự quyết tâm. Yuko-sensei đứng ở một góc xa cũng lặng lẽ quan sát chúng tôi, vẻ mặt trầm tư nhưng đầy thấu hiểu.

``\eng{Đã đến lúc rồi, Kuon-user,}'' giọng nói trầm mặc có chút hơi khàn của Game trên chiếc đồng hồ vang lên như thể đang hối thúc. ``\eng{Đã đến lúc\dots\ để họ đối diện với sự thật.}''

Tôi bước lên một bước, đứng giữa vòng tròn nhỏ của lớp 1-C. Ánh sáng từ một chùm pháo hoa khổng lồ vừa bắn lên thắp sáng toàn bộ khu vực, soi rõ từng khuôn mặt ngơ ngác của những người còn lại.

``Mọi người\dots'' Giọng tôi vang lên, trầm tĩnh nhưng đủ sức len lỏi qua tiếng nổ rền rĩ trên không trung. ``Có một chuyện\dots\ mà chúng tôi bắt buộc phải nói cho mọi người biết vào đêm nay.''

Tất cả ánh mắt đều đổ dồn về phía tôi. Tiếng pháo hoa vẫn liên tục nổ trên bầu trời, nhuộm sáng từng khuôn mặt đang ngơ ngác, nhưng vòng tròn nhỏ của chúng tôi lúc này lại tĩnh lặng đến nghẹt thở.

Tôi hít một hơi thật sâu, dồn hết can đảm để thốt ra những lời mà tôi đã chuẩn bị từ rất lâu. Những lời mà nếu nghe qua\dots\ sẽ rất là chối tai. Nhưng đó là thực tế mà họ cần được biết.

``Chúng tôi\dots\ \emph{không thuộc về thế giới này}.''

Tôi ngừng lại một nhịp, quan sát nét mặt của từng người. Vẫn là sự hoang mang, vài cái chớp mắt bối rối không hiểu chuyện gì đang xảy ra. Tôi nuốt khan, cất cao giọng hơn: ``Tôi, Kaigou-chan, và cả Suzuki nữa\dots\ Chúng tôi đến từ một dòng thời gian khác, một thế giới khác. Chúng tôi trôi dạt đến Aogami như những kẻ tị nạn không gian, với nhiệm vụ duy nhất là tìm cách thoát khỏi nơi này để quay trở về thế giới cũ. Và\dots\ việc khôi phục lại thị trấn này\dots\ là cách duy nhất để làm được điều đó.''

``Và chúng tôi\dots'' Game lúc này lên tiếng, ``\dots đã đến bước cuối cùng ngay trong đêm nay.''

``Đúng vậy,'' Kaigou bước lên đứng cạnh tôi. Gương mặt sắc lạnh và bất cần thường ngày của cô ấy biến mất, thay vào đó là một vẻ điềm tĩnh pha chút buồn bã sâu thẳm. ``Sự tồn tại của chúng tôi ở Aogami ngay từ đầu đã là một biến số lỗi. Thời gian của chúng tôi ở đây\dots\ đã hết rồi.''

Suzuki cũng chậm rãi bước tới, nắm lấy tay Kaigou. Cô nàng vốn luôn vui vẻ và phàm ăn giờ đây lại nở một nụ cười đầy gượng gạo, đôi mắt rưng rưng đỏ hoe: ``Mình không phải là học sinh chuyển trường bình thường. Bọn tớ mượn thân phận ở thế giới này chỉ để có một nơi trú ẩn tạm thời. Nhưng\dots\ tình cảm mà tớ dành cho lớp 1-C, tình cảm dành cho mọi người\dots\ hoàn toàn không phải là giả dối!''

Sự im lặng chết chóc kéo dài thêm vài giây.

``Ê này\dots\ đùa kiểu gì lạ vậy Kuon-kun?'' Honoka cười trừ, vẫy vẫy tay đầy gượng gạo. ``Cậu định đóng tiếp vở kịch khoa học viễn tưởng cho lễ hội à? Kịch bản này không có trong kế hoạch đâu nhé.''

``Đúng đấy, trò đùa này không vui chút nào đâu,'' Kaoru cũng phụ họa, từ từ hạ tờ ruy băng cuộn tròn xuống, giọng nói đã bắt đầu run rẩy khi nhìn thấy vẻ mặt cực kỳ nghiêm túc của ba chúng tôi.

Ngay lúc đó, Ayame bước lên phía trước, đứng đối diện với đám đông đang hoang mang. Wata lặng lẽ theo sát phía sau.

``Những gì Kuon-san nói\dots\ đều là sự thật,'' Ayame lên tiếng, chất giọng nghịch ngợm thường ngày được thay thế bằng sự uy quyền của một Hội trưởng Hội học sinh. ``Hội học sinh đã biết chuyện này từ trước, ngay trước ngày lễ trường học rồi.''

Inari khẽ gật đầu, rụt rè bước lên tiếp lời, đôi bàn tay đan chặt vào nhau: ``Mình\dots\ mình cũng đã biết. Nhưng dù họ có đến từ đâu, những kỷ niệm mà chúng ta có với nhau\dots\ sự giúp đỡ của Kuon-san, sự bảo bọc của Kaigou-san, nụ cười của Suzuki-san\dots\ tất cả đều là thật cơ mà!''

Để củng cố thêm cho thực tại tàn khốc ấy, Shino - cô gái giữ gìn ký ức của thị trấn - và Sakura trong bộ đồ miko vạt đỏ bước ra từ bóng tối của đài yagura.

``Với tư cách là thần tử bảo hộ của Aogami, tôi xác nhận điều này,'' Sakura gõ nhẹ cây quyền trượng xuống nền đá, âm thanh lanh lảnh vang lên xé toạc những hy vọng chối bỏ cuối cùng. ``Họ là những người ngoại lai. Một khi họ đã hoàn thành nhiệm vụ trong đêm nay\dots\ họ sẽ phải rời đi.''

``Vậy nên\dots\ Lễ hội năm nay sẽ là lần đầu tiên, và cũng là lần cuối lớp chúng ta đông đủ. Bọn tôi đã giấu chuyện này để cho các cậu không bị ảnh hưởng tâm lý,'' Nanase tiếp lời. ``Rất xin lỗi các cậu vì phải đến lúc này các cậu mới được biết.''

Shino khẽ mỉm cười buồn bã, đôi mắt ánh lên sự đồng cảm sâu sắc: ``Mình từng là nạn nhân của những vòng lặp vô tận, tôi hiểu rõ khát khao được sống một cuộc sống bình thường của họ lớn đến nhường nào. Xin mọi người\dots\ hãy hiểu cho họ.''

Đến lúc này, thực tại mới thực sự giáng một đòn chí mạng xuống tâm trí của tất cả mọi người. Lớp vỏ bọc của một đêm lễ hội hoàn hảo đã bị xé toạc, để lộ ra vết thương của sự chia ly.

Nhóm những cô gái trầm tính và giàu tình cảm như Saki, Mitsuki, và Hotaru gần như ngay lập tức đưa tay lên che miệng, nước mắt bắt đầu tuôn rơi lã chã. Saki nấc lên từng hồi, dải ruy băng trên tay cô ấy khẽ run rẩy. Miku và Kana thì ôm chầm lấy nhau, không ngừng lắc đầu như muốn từ chối tiếp nhận thông tin này.

Trong khi đó, Yae và Shiina lại đứng chết trân. Shiina siết chặt hai bàn tay thành nắm đấm đến mức các khớp xương trắng bệch. Cô nàng cắn chặt môi, đấm mạnh một cú vào cột đình bằng gỗ bên cạnh khiến lớp sơn đỏ bong tróc rớt xuống. ``Chết tiệt! Chẳng lẽ\dots\ chẳng lẽ chúng ta không thể làm gì để giúp các cậu sao?! Sức mạnh thể chất này thì có ích lợi gì khi đối mặt với những thứ không gian chết tiệt đó chứ?!''

``Kuon-san\dots\ Suzuki-san\dots\ Kaigou-san\dots'' Yae nức nở, đôi vai vận động viên rắn rỏi rung lên bần bật. ``Tại sao chứ\dots?''

Touka lùi lại một bước, vội vàng quệt đi dòng nước mắt đang làm nhòe tầm nhìn, sự điềm tĩnh và logic của một thần đồng hoàn toàn sụp đổ. Koishin cũng không còn giữ được vẻ ngoài tsundere sắc sảo nữa, cô nàng bặm môi đến bật máu, khóc nấc lên từng tiếng nặng nề.

Thậm chí cả cặp đôi chuunibyou cũng bị kéo tuột về thực tại. Shion gỡ bỏ hoàn toàn dáng vẻ uy nghi của một ``Nữ hoàng'', đôi mắt mở to thảng thốt, giọng vỡ vụn: ``Đây không phải\dots\ một trò đùa đúng không? Các cậu sẽ đi thật sao?'' Aku bên cạnh thì đã khóc òa lên thành tiếng.

Nhưng người khiến trái tim chúng tôi thắt lại đau đớn nhất, chính là Naomi.

Cô bé tám tuổi với mái tóc tết gọn gàng từ từ bước ra khỏi vòng tay của Suzuki. Đôi mắt lục bảo từng chứa đựng biết bao sự hồn nhiên và tò mò về thế giới này giờ đây ngập tràn sự hoang mang tột độ. Naomi là một sinh mệnh được chúng tôi và Game kéo về từ cõi chết, được ban cho một cơ thể con người. Đối với con bé, chúng tôi chính là gia đình duy nhất, là những người đã tạo ra hơi thở đầu tiên cho em ở thế giới này.

Naomi lảo đảo bước đi đến trước mặt tôi, bàn tay nhỏ xíu run rẩy níu lấy ống tay áo yukata của tôi.

``Onii-chan\dots\ Kaigou-oneechan\dots\ Suzuki-oneechan\dots'' giọng con bé nghẹn lại, nước mắt bắt đầu tuôn trào ồ ạt trên gò má bầu bĩnh. ``Mọi người\dots\ sẽ đi thật sao? Không đưa em đi cùng sao? Em ngoan mà\dots\ Em sẽ không đòi ăn kẹo bông gòn nữa\dots\ Mọi người đừng bỏ em lại mà\dots''

Trái tim tôi như bị ai đó bóp nghẹt thành từng mảnh. Tôi từ từ quỳ một chân xuống để tầm mắt ngang bằng với con bé, đưa tay vuốt ve mái tóc mềm mại và lau đi những giọt nước mắt nóng hổi.

``Không phải bọn anh bỏ rơi em, Naomi,'' tôi cố gắng giữ cho giọng mình không run rẩy, dù khóe mắt đã cay xè. ``Cơ thể của em được tạo ra từ linh khí của Aogami. Em đã thuộc về thế giới này. Nếu em đi cùng bọn anh qua cánh cổng không gian, không chắc điều gì sẽ xảy ra với em đâu. Lát nữa, chị Shion và chị Sakura và anh chị sẽ làm Nghi thức Thanh tẩy để em được sống như một con người thực sự.''

``Nhưng\dots\ nhưng em không muốn ở lại một mình!'' Naomi ôm chầm lấy cổ tôi, òa khóc nức nở, vùi mặt vào hõm vai tôi. ``Em muốn ở cùng mọi người cơ! Em không cần làm người nữa!''

``Em sẽ không ở một mình đâu, Nao,'' Kaigou cũng quỳ xuống, vòng tay ôm lấy cả tôi và Naomi. Những giọt nước mắt mặn chát cuối cùng cũng lăn dài trên má của cô nàng luôn khoác lên mình vẻ ngoài gai góc. ``Em còn có chị Nanase, chị Ayame, và tất cả mọi người ở đây. Bọn họ\dots\ sẽ là gia đình mới của em. Em phải sống thay cho cả phần của bọn chị nữa.''

Suzuki cũng nhào tới ôm chầm lấy chúng tôi, tạo thành một vòng tay của những kẻ ngoại đạo đang khóc than cho số phận của chính mình.

Tiếng pháo hoa trên trời vẫn tiếp tục nổ rực rỡ, nhưng dưới mặt đất, chỉ còn lại những tiếng nấc nghẹn ngào của sự chia ly. Những giọt nước mắt rơi xuống nền đá sân đền, thấm đẫm tình cảm chân thành mà chúng tôi đã dành cho nhau trong suốt mùa hè ngắn ngủi này. Dù cho có là người ngoại lai hay người bản địa, thì những giọt nước mắt này hoàn toàn không có sự phân biệt.

``\eng{Con người thật kỳ lạ,}'' giọng nói của Game từ chiếc đồng hồ vang lên, trầm mặc và không cảm xúc như một cỗ máy đang phân tích dữ liệu, nhưng lại chứa đựng một sự thấu hiểu đáng kinh ngạc về bản chất của sinh mệnh. ``\eng{Các người bỏ ra hàng nghìn giờ để tạo dựng những sợi dây liên kết, rồi lại dùng hàng vạn giọt nước mắt để tiễn đưa chúng vào hư vô. Một sự lãng phí tài nguyên cảm xúc vô nghĩa, nếu nhìn dưới góc độ logic.}''

``\eng{Thế nhưng, chính sự tan vỡ này lại là minh chứng duy nhất cho thấy sự tồn tại của các người không phải là những con số vô hồn trong dòng thời gian. Và\dots\ chia ly chính là cái giá phải trả cho một mùa hè không có vòng lặp. Dường như các người đã sẵn sàng thanh toán cái giá đó ngay từ lúc bắt đầu rồi.}''

Khi Naomi đã gục đầu vào vai Kaigou, tôi từ từ đứng dậy. Hai chị em cũng gạt nước mắt, kiên cường đứng lên theo.

Ba người chúng tôi, những kẻ ngoại lai mang trên mình án tử của những dòng thời gian, cùng bước lùi lại một bước để đối diện với toàn bộ những người bạn quan trọng nhất. Dưới ánh sáng lấp lánh của pháo hoa và ánh trăng rằm, chúng tôi đồng loạt cúi gập người một góc chín mươi độ.

``Cảm ơn mọi người\dots\ vì đã cho chúng tôi một mùa hè rực rỡ nhất trong đời,'' giọng tôi vang lên, rõ ràng và kiên định, dù cuống họng đã nghẹn đắng. ``Và xin lỗi\dots\ vì đã giấu kín chuyện này, biến mọi người thành những người bạn của những kẻ dối trá.''

``Xin lỗi mọi người,'' Suzuki thút thít, vai rung lên.

``Cảm ơn\dots\ và xin lỗi,'' Kaigou thì thầm, giọng nói vỡ òa.

Chúng tôi giữ nguyên tư thế cúi đầu đó, mặc cho những tiếng nấc từ đám đông xung quanh ngày một lớn hơn. Giao ước đã được lật mở. Sự thật đã phơi bày. Đêm lễ hội mùa hè ngày đó là đêm huy hoàng nhất, và cũng là đêm chia tay buồn bã nhất.

\ct

Mười giờ đêm.

Những tiếng sụt sùi vẫn vang lên rải rác trong không gian, nhưng sự tĩnh lặng u buồn nhanh chóng bị phá vỡ bởi một tràng vỗ tay dõng dạc.

``\textbf{Được rồi! Tất cả lau nước mắt đi!}'' Ayame lớn giọng, chiếc loa cầm tay đã được hạ xuống nhưng âm lượng của cô nàng Hội trưởng vẫn đủ sức vực dậy tinh thần của toàn bộ lớp 1-C. ``Thời gian của họ ở thế giới này không còn nhiều, chúng ta không thể lãng phí nó vào việc khóc lóc ủy mị được!''

Wata đi vòng quanh, nhẹ nhàng đưa khăn giấy cho những cô gái vẫn còn đang nức nở như Saki hay Miku. Inari cũng gật đầu mạnh mẽ, đôi mắt dù còn hoe đỏ nhưng đã tràn đầy sự quyết tâm.

``Hội trưởng nói đúng đấy,'' Nanase bước lên ngang hàng với Kuon, tà áo yukata màu xanh thẫm khẽ bay trong gió đêm. ``Và chúng ta còn một việc cực kỳ quan trọng phải làm trong đêm nay. Đó là lý do tôi đã nhờ gia đình Furukawa giữ lại một số người dân thị trấn nán lại sau màn pháo hoa. Kuon-san, Kaigou-san, Suzuki-san, hãy đưa Naomi-chan đi theo chúng tôi.''

Tôi nhìn cô nàng Tiểu thư, rồi quay sang nhìn Kaigou và Suzuki. Cả hai đều gật đầu. Tôi nắm lấy bàn tay bé nhỏ đang run rẩy của Naomi. ``Đi thôi em. Mọi người đang cố gắng vì em đấy.''

Dưới sự dẫn đường của Hội học sinh, chúng tôi rời khỏi khoảng sân chính của đền Aogami, vòng qua khu rừng bách thảo phía sau đền để tiến ra mỏm đá, tại đúng vị trí mà chiếc hố chôn cất từ lễ hiến tế. Đây là nơi mà ánh trăng rằm chiếu thẳng xuống sáng nhất, đồng nghĩa với tỷ lệ thành công cao nhất - hoặc ít nhất, là những gì mà Nanase nói với tôi như vậy và được Game đồng tình.

Dọc theo con đường mòn, những chiếc đèn lồng giấy nhỏ đã được thắp sáng từ trước, tỏa ra ánh sáng vàng vọt ấm áp. Tiếng xào xạc của lá cây xao động bởi gió đêm len lỏi vào tai tôi, xen lẫn với tiếng thì thầm đọc bản đồ của Nanase ở phía dẫn đầu.

Khi chúng tôi đến nơi, một khung cảnh kỳ bí và tĩnh lặng hiện ra trước mắt.

Câu lạc bộ Huyền bí đã có mặt ở đây từ bao giờ. Trên nền đá bằng phẳng, một vòng tròn ma thuật khổng lồ với những họa tiết hình học phức tạp đã được vẽ hoàn chỉnh bằng phấn trắng và muối biển. Những ngọn nến đỏ được đặt cẩn thận ở các điểm nút quan trọng, ngọn lửa cháy leo lét nhưng kỳ lạ thay lại không hề bị gió biển thổi tắt.

Yuka đang ngồi quỳ ở một góc, hai tay bưng một chiếc lư hương bằng đồng. Làn khói trắng mỏng manh tỏa ra mang theo mùi hương trầm cổ kính, khiến tâm trí bất cứ ai ngửi thấy đều trở nên bình lặng và thanh tĩnh lạ thường. Aku lúc này đang cầm một xô muối, cẩn thận rắc nốt những đường nét cuối cùng của vòng tròn. Sự tập trung và độ chính xác của một game thủ đỉnh cao cuối cùng cũng được ứng dụng vào thực tế.

Sakura đứng cạnh vòng tròn, bộ đồ miko truyền thống tung bay trong gió. Và ở vị trí trung tâm, người đang chỉ đạo toàn bộ quá trình thiết lập này, chính là Shion.

Tiếng bước chân của đám đông khiến Shion quay lại. Cô nàng từ từ tiến về phía chúng tôi.

Tôi khẽ chớp mắt, gần như không nhận ra cô bạn cùng lớp của mình. Shion không mặc chiếc áo choàng đen kỳ dị, cũng không diện bộ trang phục kỳ quái nào. Không có những tư thế vặn vẹo, không có những nụ cười bí hiểm ảo tưởng sức mạnh, dáng đứng của cô lúc này thẳng tắp và tràn đầy sự tự tin của một người trưởng thành.

``Mọi người đến rồi,'' Shion cất giọng. Không còn những từ ngữ cổ ngữ khó hiểu, không còn cách xưng hô ``ta - ngươi'' trịch thượng. Giọng nói của cô ấy trong trẻo, điềm tĩnh và nghiêm túc đến mức khiến cả lớp 1-C phải sững sờ.

``Shion\dots\ cậu\dots'' Tôi ngập ngừng, không biết phải phản ứng thế nào trước sự thay đổi đột ngột này.

Shion khẽ mỉm cười, một nụ cười dịu dàng và chân thật hiếm thấy. ``Đêm nay không có 'Nữ hoàng Bóng tối' nào cả, Kuon. Ở đây chỉ có Shitou Shion, một thành viên của lớp 1-C đang muốn dùng toàn bộ sức lực của mình để bảo vệ người mà các cậu thương yêu nhất.''

Cô ấy đảo mắt nhìn quanh những khuôn mặt ngỡ ngàng của mọi người, khi chỉ sau một biến cố đó, cô nàng đã thay đổi 180 độ. ``Tôi xin lỗi vì đã luôn diễn những trò kỳ quái khiến mọi người thấy phiền phức. Nhưng tôi hứa, từ giây phút này trở đi, mọi thứ tôi làm đều hoàn toàn nghiêm túc.''

Sự tĩnh lặng bao trùm lấy mỏm đá. Không ai lên tiếng trêu chọc, không ai cười cợt. Shiina khẽ mỉm cười gật đầu, còn Aku thì lén lau một giọt nước mắt tự hào khi thấy cô bạn thân của mình rũ bỏ lớp vỏ bọc chuunibyou.

``\eng{\dots Giá mà cô ấy lịch sự hơn một chút thì tốt biết mấy,}'' Game khẽ lên tiếng, chất giọng quay trở về mỉa mai.

``\eng{Không phải lúc này, Game,}'' tôi suỵt một tiếng nhẹ. ``\eng{Tập trung vào nghe đi.}''

Shion quay lại vấn đề chính, bước đến mép vòng tròn. ``Như mọi người đã biết, Nghi thức Thanh tẩy Đêm trăng rằm là giải pháp cuối cùng để giữ Naomi lại thế giới này. Thời gian qua, tôi đã cùng Sakura và Shino nghiên cứu toàn bộ hệ thống linh mạch của Aogami. Những ghi chép cổ của đền thờ kết hợp với ký ức nguyên thủy của Shino đã giúp chúng tôi hoàn thiện ma trận này.''

Shion chỉ tay vào trung tâm vòng tròn. ``Cơ thể của Naomi-chan được cấu thành từ linh khí của thị trấn Aogami. Nhưng linh hồn em ấy lại rất lỏng lẻo vì được đánh thức từ sự hư vô. Nếu em ấy bước qua cánh cổng không gian cùng Kuon, linh hồn sẽ lập tức vỡ nát do không chịu nổi áp lực thứ nguyên. Để cứu em ấy, chúng ta phải 'neo' linh hồn đó vào cốt lõi của Aogami.''

``Nghi thức sẽ diễn ra theo ba bước,'' Sakura tiếp lời, gõ nhẹ quyền trượng xuống đất. ``Đầu tiên, Naomi sẽ đứng ở tâm vòng tròn. Shino và tôi - những người mang đậm linh khí nguyên thủy của thị trấn nhất - sẽ đứng ở hai điểm Thái Cực bên trong, đóng vai trò là cột thu lôi để dẫn truyền năng lượng.''

``Bước thứ hai là nhiệm vụ của tất cả mọi người,'' Shion nhìn quanh toàn bộ lớp 1-C và nhóm người dân thị trấn đang đứng ở phía sau. ``Linh mạch của Aogami vốn dĩ là những dòng chảy ý thức tập thể. Mọi người hãy tạo thành một vòng tròn lớn bao bọc lấy ma trận này. Hãy nắm tay nhau, nhắm mắt lại và hướng mọi suy nghĩ, mọi sự yêu thương về phía Naomi. Hãy để thị trấn này biết rằng, các cậu chấp nhận em ấy là một phần máu thịt của Aogami. Sự kết nối chân thành của con người chính là loại phép thuật mạnh mẽ nhất.''

``Và cuối cùng,'' Yuka cất giọng đều đều đằng sau làn khói trầm hương. ``Khi mặt trăng lên đến đỉnh đầu\dots\ ánh trăng rằm sẽ hoạt động như một chiếc khóa vạn năng, đóng chặt phong ấn lại. Linh hồn của Naomi sẽ dung hợp hoàn toàn với quy luật của thế giới này.''

Tôi nghe xong, cảm thấy một tia hy vọng bùng lên mạnh mẽ. ``Vậy tôi, Kaigou-chan và Suzuki phải đứng ở vị trí nào?''

Shion khẽ lắc đầu, ánh mắt lộ vẻ tiếc nuối. ``Đó là vấn đề lớn nhất. Ba người các cậu là những thực thể ngoại lai. Năng lượng của các cậu hoàn toàn dị biệt với thế giới này. Nếu các cậu bước vào vòng tròn hay thậm chí là tham gia truyền ý thức, nó sẽ gây nhiễu loạn linh mạch và làm nghi thức thất bại. Các cậu\dots\ chỉ có thể đứng ngoài quan sát.''

Tôi khựng lại, bàn tay đang nắm lấy tay Naomi vô thức siết chặt hơn. Sự bất lực lại một lần nữa bủa vây lấy tôi. Ngay cả trong khoảnh khắc quyết định sinh mạng của con bé, tôi cũng không thể tự tay bảo vệ em ấy.

``Không sao đâu,'' một bàn tay ấm áp đặt lên vai tôi. Kaigou mỉm cười, nụ cười thanh thản nhất mà tôi từng thấy ở chị ấy. ``Nhiệm vụ của chúng ta là đưa con bé đến thế giới này. Phần còn lại\dots\ hãy để thế giới này ôm lấy con bé.''

Suzuki cũng gật đầu, lau đi những giọt nước mắt còn đọng lại trên mi. ``Đúng vậy. Lớp 1-C chưa bao giờ làm chúng ta thất vọng, đúng không Kuon?''

Tôi hít một hơi thật sâu, dằn lại cảm giác nghẹn ngào nơi cổ họng, rồi từ từ quỳ một chân xuống trước mặt Naomi. Cô bé tám tuổi chớp chớp đôi mắt lục bảo, sự hoảng loạn ban nãy đã vơi đi phần nào khi nhìn thấy sự quyết tâm của mọi người xung quanh.

``Naomi này,'' tôi đưa tay vuốt lại mái tóc tết của con bé. ``Em nghe rõ rồi chứ? Bọn anh không thể đứng trong vòng tròn để bảo vệ em, nhưng bọn anh sẽ đứng ngay đây, không rời nửa bước. Em hãy nhìn xem,'' Tôi chỉ tay về phía những người bạn đang xắn tay áo, đan tay vào nhau để chuẩn bị. ``Tất cả những người yêu thương em đều đang ở đây. Đừng sợ hãi. Cứ dũng cảm bước vào đó, và ngày mai, em sẽ thức dậy đón ánh nắng của mùa hè như một cô bé bình thường.''

Naomi nhìn tôi, rồi nhìn Kaigou và Suzuki. Con bé cắn môi, những giọt nước mắt lại ứa ra nhưng lần này, em đưa tay tự mình lau đi. Naomi gật đầu thật mạnh.

``Vâng! Em sẽ dũng cảm! Em sẽ sống thật tốt để mọi người không phải lo lắng!''

Dứt lời, Naomi buông tay tôi ra, lững thững bước về phía trung tâm vòng tròn ma thuật. Chiếc yukata màu hoa anh đào của con bé nổi bật giữa những ánh nến đỏ rực. Shino và Sakura cũng bước vào, đứng đối diện nhau ở hai đầu vòng tròn.

Bên ngoài, lớp 1-C đã nhanh chóng vào vị trí. Shiina nắm lấy tay Saki, Yae đan tay với Azuki. Koishin tuy miệng vẫn lầm bầm phàn nàn nhưng hai tay lại siết chặt lấy tay Touka và Honoka. Vòng tròn người cứ thế mở rộng ra, bao quanh toàn bộ khu vực mỏm đá. Những người dân thị trấn cũng hòa vào vòng tròn ấy, nhắm nghiền mắt lại.

Shion lấy từ trong túi áo ra một nén hương, châm lửa từ ngọn nến trung tâm rồi cắm xuống bệ đá trước mặt Naomi. Gió biển đột ngột ngừng thổi. Khói từ nén hương không tạt sang hai bên mà bay thẳng đứng lên bầu trời đêm, hướng về phía mặt trăng đang từ từ tiến đến điểm cao nhất.

Cô nàng tóc xám-trắng hít một hơi, bước lùi lại để hòa vào vòng tròn cùng với mọi người. Cô nhắm mắt lại, âm giọng phát ra trầm ấm vang vọng khắp mỏm đá.

``Thời khắc đã đến. Mọi người\dots\ Bắt đầu thôi.''

Tiếng xào xạc của lá cây như chìm hẳn xuống, nhường chỗ cho một sự tĩnh lặng đến nghẹt thở. Hàng chục, hàng trăm con người - từ những học sinh lớp 1-C cho đến những người dân thị trấn - đồng loạt nhắm mắt lại. Không một tiếng động nào phát ra từ họ, nhưng trong không khí lại bắt đầu rung lên một thứ âm thanh trầm thấp, giống như tiếng ngân nga của một chiếc chuông đồng khổng lồ bị chôn vùi dưới lòng đất.

Đó là sự cộng hưởng của ý thức.

Từ hai điểm Thái Cực bên trong vòng tròn, Shino và Sakura bắt đầu cất giọng đọc chú. Đó không phải là những câu thần chú chuunibyou mà Shion hay dùng để trêu đùa, mà là những lời cầu nguyện cổ xưa của Thần đạo, được truyền lại từ những thế hệ đầu tiên khai sinh ra Aogami. Giọng của Shino dịu dàng và mang tính vỗ về, trong khi giọng của Sakura lại lanh lảnh và uy nghiêm. Hai âm sắc hòa quyện vào nhau, tạo thành một sợi dây vô hình kết nối toàn bộ những người đang đứng ở vòng ngoài.

``\eng{Báo cáo trạng thái: Mật độ linh khí xung quanh đang tăng lên mức 40\%. Các luồng năng lượng ý thức từ các cá thể con người đang tập trung về phía trung tâm ma trận.}''

Giọng nói của Game vang lên đều đều từ chiếc đồng hồ, chỉ đủ lớn để tôi, Kaigou và Suzuki nghe thấy. Sự phân tích lạnh lùng của hệ thống đối lập hoàn toàn với bầu không khí tâm linh đang diễn ra, nhưng lại là bằng chứng vững chắc nhất cho thấy nghi thức thực sự có tác dụng.

Dưới nền đá, những đường nét được vẽ bằng phấn và muối biển bắt đầu phát ra một thứ ánh sáng mờ ảo. Ánh sáng ấy không chói lóa, mà mang màu bàng bạc giống hệt như ánh trăng đang chiếu rọi từ trên cao. Khói trầm hương từ chiếc lư đồng của Yuka không còn bay lơ lửng nữa, mà cuộn lại thành từng dải ruy băng mỏng, xoay vòng quanh vị trí Naomi đang đứng.

Cô bé tám tuổi nhắm nghiền mắt, hai bàn tay nhỏ xíu nắm chặt lấy vạt áo yukata màu hoa anh đào. Dù không có gió, tà áo của con bé vẫn khẽ bay lật phật. Đôi mắt lục bảo thỉnh thoảng hé mở, phản chiếu ánh sáng bạc của ma trận bên dưới.

``\eng{Cảnh báo: Tần số không gian tại tâm vòng tròn đang dao động dữ dội. Tỷ lệ đồng bộ linh hồn của thực thể "Naomi" hiện đạt 65\%. Hệ thống linh mạch của Aogami đang bắt đầu tiến trình tiếp nhận.}''

``Cố lên, Naomi\dots'' Suzuki thì thầm, hai tay đan vào nhau cầu nguyện. Kaigou đứng bên cạnh, hai khoanh tay trước ngực nhưng những ngón tay lại đang bấu chặt vào bắp tay đến mức gân xanh nổi lên.

Mặt trăng dần di chuyển đến vị trí cao nhất trên bầu trời. Ánh sáng của nó như quy tụ lại thành một luồng sáng duy nhất, chiếu thẳng đứng xuống mỏm đá, bao trùm lấy toàn bộ vòng tròn ma thuật.

Nhiệt độ xung quanh dường như giảm xuống, nhưng những giọt mồ hôi lại bắt đầu lấm tấm trên trán của những người đang đứng ở vòng ngoài. Shiina cắn chặt môi, cơ bắp căng cứng như đang dùng sức mạnh thể chất để ép ý chí của mình truyền đi. Saki và Mitsuki nắm tay nhau chặt đến mức các đốt ngón tay trắng bệch. Dù không ai mở mắt, nhưng tôi có thể cảm nhận được khao khát mãnh liệt của họ: ``\emph{Hãy để con bé ở lại đây.}''

Shino và Sakura đồng loạt giơ hai tay lên cao, hướng về phía Naomi. Mặt đất dưới chân chúng tôi khẽ rung chuyển. Không phải một trận động đất, mà là cảm giác một thực thể khổng lồ đang cựa mình thức giấc. Linh mạch của Aogami đã hoàn toàn mở ra.

``\eng{Năng lượng thứ nguyên vượt ngưỡng 90\%. Đang tiến hành quá trình khóa vĩnh viễn (Permanent Anchor). Yêu cầu duy trì sự ổn định ý thức trong 15 giây nữa.}''

Nhưng ngay lúc đó, một sự cố xảy ra.

Naomi đột ngột thở dốc. Cơ thể nhỏ bé của em ấy từ từ nhấc bổng lên cách mặt đất chừng vài chục xăng-ti-mét. Những luồng ánh sáng bạc xung quanh đột nhiên trở nên hỗn loạn, quất mạnh vào nhau tạo ra những tia lửa điện nhỏ. Linh hồn của Naomi vốn được kéo về từ cõi chết, từ sự hư vô, nên lực đẩy tự nhiên của thế giới này đang cố gắng bài xích sự tồn tại dị biệt đó ở khoảnh khắc dung hợp cuối cùng.

``Á\dots\ đau quá\dots!'' Naomi thốt lên, khuôn mặt nhăn nhó vì đau đớn.

``Naomi!'' Tôi hét lên, vô thức bước một chân về phía trước, định lao vào vòng tròn để kéo con bé ra.

``\eng{\textbf{Đừng lại gần đó, Kuon-user!} Nếu cậu bước vào vùng ma trận bây giờ, năng lượng ngoại lai từ cơ thể cậu sẽ tạo ra phản ứng chuỗi, xé toạc linh hồn cô bé thành từng mảnh ngay lập tức!}''

Tiếng quát điện tử vang lên sắc lạnh như một gáo nước lạnh tát thẳng vào mặt tôi, đóng băng đôi chân đang định bước đi. Tôi chết sững tại chỗ, chỉ biết giương mắt nhìn Naomi đang vật vã giữa không trung. Bất lực. Sự bất lực tột cùng lại một lần nữa gặm nhấm tâm can tôi.

``\textbf{Đừng bỏ cuộc!}'' Nanase đột nhiên hét lớn từ vòng ngoài, đôi mắt vẫn nhắm nghiền nhưng giọng nói lại vang vọng đầy uy quyền. ``Mọi người, tập trung hơn nữa! Nghĩ về nụ cười của em ấy! Nghĩ về việc ngày mai chúng ta sẽ cùng em ấy đi ăn kẹo bông gòn!''

``\textbf{Đúng vậy! Bọn chị chưa dạy em cách làm bánh pancake đâu đấy, Naomi-chan!}'' Touka cũng gào lên, phá vỡ hoàn toàn sự điềm tĩnh thường ngày.

``Chị còn hứa sẽ may cho em một bộ váy công chúa mà!'' Miku nức nở tiếp lời.

Từng tiếng nói, từng lời hứa hẹn của lớp 1-C vang lên giữa đêm khuya. Những dòng năng lượng vô hình nhưng mãnh liệt bùng phát từ vòng ngoài, xông thẳng vào trung tâm ma trận. Những tia lửa điện hỗn loạn dần bị dập tắt, thay vào đó là một thứ ánh sáng rực rỡ và thuần khiết bao bọc lấy cơ thể đang lơ lửng của Naomi.

Shino và Sakura dồn toàn bộ sức lực, dập mạnh hai tay xuống nền đá.

``\eng{Tỷ lệ đồng bộ đạt\dots\ 95\%\dots\ 98\%\dots\ 100\%.}''

Một tiếng nổ trầm đục vang lên, không phải từ không khí mà từ sâu trong tâm trí của mỗi người. Luồng ánh sáng bạc bùng lên chói lóa đến mức tôi phải đưa tay lên che mắt, rồi ngay lập tức vỡ vụn thành hàng triệu đốm sáng nhỏ li ti, rơi lất phất xuống mỏm đá như một cơn mưa sao băng tuyệt đẹp.

Gió rừng thổi trở lại, cuốn tan những tàn khói trầm hương cuối cùng.

``\eng{Phong ấn linh hồn hoàn tất. Lực bài xích đã bị triệt tiêu hoàn toàn. Thực thể ``Naomi-chan'' hiện đã được hệ thống thứ nguyên định danh là một cư dân hợp lệ của không-thời gian này.}''

Tiếng báo cáo đều đều của Game chưa bao giờ nghe êm tai đến thế.

Naomi từ từ rơi xuống, đôi chân trần chạm nhẹ lên nền đá. Ánh sáng từ vòng tròn ma thuật dần vụt tắt. Cô bé lảo đảo một chút, rồi chớp chớp mắt nhìn xung quanh.

``\textbf{Naomi!}''

Tôi, Kaigou và Suzuki đồng loạt lao tới. Suzuki là người ôm chầm lấy con bé đầu tiên, theo sau là tôi và Kaigou. Vòng tay của chúng tôi ôm siết lấy cơ thể bé nhỏ ấy. Không còn cảm giác lạnh lẽo của một thực thể mong manh nữa. Cơ thể Naomi bây giờ mang hơi ấm của sự sống, của một con người thực sự, trái tim con bé đập những nhịp đập vững vàng áp vào ngực tôi.

``Em làm được rồi\dots'' Naomi thều thào, khóe môi nhợt nhạt vẽ lên một nụ cười rạng rỡ. ``Em thấy\dots\ ấm lắm\dots''

Xung quanh chúng tôi, những tiếng thở phào nhẹ nhõm đồng loạt vang lên. Một số người ngã khuỵu xuống nền đá vì kiệt sức sau khi vắt kiệt tinh thần, nhưng trên môi ai cũng nở một nụ cười mãn nguyện. Aku nằm bẹp xuống đất thở hồng hộc, trong khi Shion chỉ khẽ đẩy nhẹ mái tóc ướt đẫm mồ hôi của mình, mỉm cười đầy tự hào.

Chúng tôi đã làm được. Lớp 1-C đã làm được. Một sinh mạng đã được giữ lại với thế giới này bằng chính phép màu của tình người.

Thế nhưng, ngay trong khoảnh khắc hạnh phúc ngắn ngủi ấy, khi tôi vừa định cõng Naomi lên lưng để cùng mọi người quay về nhà thì chiếc đồng hồ trên tay tôi đột ngột rung lên bần bật.

Nó không phải là kiểu rung nhẹ nhàng báo hiệu tin nhắn hay phân tích thông thường. Lõi từ của chiếc đồng hồ nóng rực lên, và màn hình hiển thị nhấp nháy một màu đỏ chót đầy đe dọa: ``\eng{CẢNH BÁO ĐỎ! PHÁT HIỆN BIẾN ĐỘNG NĂNG LƯỢNG THỨ NGUYÊN KHỔNG LỒ!}''

Giọng nói của Game mất đi sự bình thản thường ngày, tốc độ nói nhanh hơn và mang đầy tính cấp bách.

Tôi sững người, nhịp tim vừa dịu xuống lại bắt đầu đập loạn nhịp. ``\eng{Cái gì? Có chuyện gì vậy Game? Chúng ta vừa mới phong ấn thành công rồi cơ mà!?}''

``\eng{Việc neo giữ một linh hồn được tái sinh bằng linh khí của toàn bộ thị trấn đã tạo ra một vụ nổ năng lượng vi mô ở tầng sâu của linh mạch. Năng lượng này đóng vai trò như một chất xúc tác, cộng hưởng với năng lượng ngoại lai tích tụ từ cơ thể các cậu.}''

Kaigou nhíu mày, bước tới gần tôi: ``\eng{Nói ngôn ngữ con người đi! Cụ thể là chuyện chết tiệt gì đang xảy ra?}''

``\eng{Nói cách khác, một lỗ xoáy không gian - cánh cổng để đưa các cậu rời khỏi thế giới này - đang được gia tốc hình thành.}''

Nghe đến đây, cả ba chúng tôi gần như vô thức trút ra một hơi thở phào nhẹ nhõm. Một sự gia tốc đồng nghĩa với việc chúng tôi sẽ sớm thoát khỏi chuỗi ngày trốn chạy vô tận này. Khác với ba lần trước đó, nơi lỗ xoáy cứ lù lù như vậy xuất hiện khiến chúng tôi không kịp trở tay, lần này cánh cổng sẽ hình thành từ từ, cho chúng tôi thêm thời gian để chuẩn bị.

Kaigou khẽ thả lỏng đôi vai đang gồng cứng từ nãy đến giờ, trong khi Suzuki mỉm cười quay sang nhìn tôi, ánh mắt lấp lánh hy vọng rằng chúng tôi vẫn còn ít nhất vài ngày ngắn ngủi, đủ để cùng lớp 1-C tận hưởng nốt kỳ nghỉ hè và tổ chức một buổi tiệc chia tay tử tế.

``\eng{Nhưng\dots}'' giọng nói vô cảm của Game lại tiếp tục vang lên, như vừa chặn đứng khả năng trên, ``\eng{thời gian dự kiến mở cổng không phải là vài tháng hay vài ngày nữa\dots\ mà là bình minh ngày mai, lúc sáu giờ sáng.}'' Câu nói của ông ta như một lưỡi dao sắc lẹm cắt đứt bầu không khí ăn mừng vừa mới nhen nhóm.

``\eng{Bình minh\dots\ ngày mai sao?}'' Suzuki thốt lên, đôi mắt mở to bàng hoàng. ``\eng{Nhưng\dots\ bây giờ đã là nửa đêm rồi! Tức là chúng ta chỉ còn\dots}''

``\eng{Chính xác. Sáu tiếng đồng hồ nữa. Khi tia sáng đầu tiên của ngày mới chiếu xuống Aogami, cánh cổng sẽ mở ra. Nếu các cậu không bước qua nó, chúng ta sẽ không thể nào biết được tới khi nào nó sẽ xuất hiện trở lại đâu. Sức mạnh của nơi đây kết hợp với của các cậu cũng chỉ đủ để duy trì cổng trong vòng mười lăm phút thôi.}''

Tôi cúi gầm mặt, nghiến chặt hai hàm răng. Cứ tưởng rằng chúng tôi có thể thư thả tận hưởng những ngày hè còn lại, nhưng không, số phận không bao giờ cho phép những kẻ tị nạn không gian như chúng tôi được nghỉ ngơi.

Tất cả chỉ còn lại một đêm nay. Một đêm cuối cùng.

\end{document}